% ThinkAlgo — Number Sense — Grade 7
% Edition: teacher (with solutions), EN.  Generated by build_grade7.py — edit body/ and rebuild.
\documentclass[11pt]{article}
\newif\ifsol \soltrue
\usepackage{iftex}
\ifPDFTeX
  \usepackage[utf8]{inputenc}\usepackage[T5]{fontenc}
\else
  \usepackage{fontspec}
\fi
\usepackage[a4paper,top=2.95cm,headheight=1.5cm,headsep=0.45cm,bottom=3.1cm,left=1.9cm,right=1.9cm]{geometry}
\usepackage{amsmath,amssymb}
\usepackage{tikz}\usetikzlibrary{calc,shadows,shapes.geometric}
\usepackage[most]{tcolorbox}
\usepackage{enumitem}
\usepackage{array}
\usepackage{tabularx}
\usepackage{colortbl}
\usepackage{comment}
\usepackage{needspace}
\usepackage{fancyhdr}
\usepackage{fontawesome5}
\usepackage{graphicx}
\usepackage{microtype}
% PDF outline (bookmarks) so each chapter shows in the reader's outline — does not change the page layout
\usepackage[hidelinks,unicode,bookmarksnumbered=false,bookmarksopen=true,bookmarksopenlevel=0]{hyperref}
\usepackage{bookmark}
\usetikzlibrary{decorations.pathreplacing,arrows.meta,angles,quotes,patterns,positioning,through,intersections}

% ---------------- palette: navy family + subtle light tint; orange ONLY in the logo ----------------
% ----- EXACT palette from the user's original main (36).tex (vivid, not pale) -----
\definecolor{tmBlue}{RGB}{66,133,244}   % problem / header / section
\definecolor{tmGreen}{RGB}{52,168,83}   % solution
\definecolor{tmOrange}{RGB}{255,152,0}  % example
\colorlet{Navy}{tmBlue}                 % header rule / section heading / badge / divider
\definecolor{NavyD}{HTML}{222222}       % solution body text (dark, readable)
\colorlet{Gold}{Navy}
\colorlet{ProbBlue}{tmBlue}             % problem boxes
\colorlet{SolGreen}{tmGreen}            % solution boxes
\definecolor{AnsRed}{HTML}{E02424}      % vivid red for answers / totals
\definecolor{CardGreen}{HTML}{34A853}   % digit-card green
\definecolor{StarGold}{HTML}{F2B134}    % difficulty star (gold)
\definecolor{KidOrange}{HTML}{DE8A2C}   % logo accent only
\definecolor{Ink}{HTML}{1F2733}       % body text
\definecolor{Slate}{HTML}{5B6A82}     % muted labels
\definecolor{Line}{HTML}{C6D0DF}      % light hairlines / student boxes
\definecolor{Tint}{HTML}{EDF1F7}      % pale panel fill (worked solution / answer)
% aliases so older names keep working
\colorlet{TextDark}{Ink}\colorlet{TextGray}{Slate}\colorlet{Soft}{Slate}
\colorlet{CardLine}{Line}\colorlet{Teal}{Navy}\colorlet{ProbB}{Navy}\colorlet{SolG}{Navy}
% bright, kid-friendly colours for bar-model (sơ đồ đoạn thẳng) blocks
\definecolor{tmPurple}{RGB}{124,77,226}
\definecolor{tmTeal}{RGB}{0,150,136}
\definecolor{tmPink}{RGB}{233,30,99}
\AtBeginDocument{\color{Ink}}
\raggedbottom  % natural flow, even spacing (no stretched gaps)
\setlength{\parindent}{0pt}\setlength{\parskip}{0.18em}
\renewcommand{\arraystretch}{1.35}
\providecommand{\AnsWord}{Đáp số}

% ---------------- brand logo (original mark) ----------------
\newcommand{\logomark}[1]{\begin{tikzpicture}[baseline=-0.55ex,scale=#1]
  \draw[KidOrange,line width=1.5pt] (-0.27,0) .. controls (-0.27,0.23) and (0,0.23) .. (0,0)
        .. controls (0,-0.23) and (-0.27,-0.23) .. (-0.27,0);
  \draw[Navy,line width=1.5pt] (0.27,0) .. controls (0.27,0.23) and (0,0.23) .. (0,0)
        .. controls (0,-0.23) and (0.27,-0.23) .. (0.27,0);
\end{tikzpicture}}
% real brand logo image (gold infinity + "think algo" wordmark)
\newcommand{\logohead}{\IfFileExists{logo.png}{\includegraphics[height=1.35cm]{logo.png}}%
  {\raisebox{0.35cm}{\logomark{1.5}\hspace{0.5em}{\Large\bfseries\color{KidOrange}think\,\color{Navy}algo}}}}

% ---------------- running header (logo + Name/Date + topic) on EVERY page, single navy rule ----------------
\pagestyle{fancy}\fancyhf{}
\renewcommand{\headrulewidth}{1.4pt}\renewcommand{\footrulewidth}{0pt}
\renewcommand{\headrule}{{\color{Navy}\hrule height 1.3pt}}
\fancyhead[L]{\raisebox{-0.95cm}{\logohead}}
\fancyhead[R]{\setlength{\tabcolsep}{0pt}%
  \begin{tabular}{@{}r@{}}
    {\small\color{Ink}{\color{tmBlue}\faUserEdit}~\NameWord:~\rule{3.0cm}{0.5pt}\quad{\color{tmBlue}\faCalendarDay}~\DateWord:~\rule{1.9cm}{0.5pt}}\\[5pt]
    {\large\bfseries\color{Navy}\CurTopic}
  \end{tabular}}

% ---------------- footer (rounded page border + Topic/Grade + thinkalgo.com + page + tagline) ----------------
% NOTE: a remember-picture overlay must be placed DIRECTLY in \fancyfoot, NOT inside \def/\newcommand.
\fancyfoot[C]{\begin{tikzpicture}[remember picture, overlay]
  \draw[Navy!40, line width=1.4pt, rounded corners=12pt]
     ($(current page.north west)+(0.8cm,-0.8cm)$) rectangle ($(current page.south east)+(-0.8cm,0.45cm)$);
  \coordinate (FL) at ($(current page.south west)+(1.2cm, 2.18cm)$);
  \coordinate (FR) at ($(current page.south east)+(-1.2cm, 2.18cm)$);
  \draw[Navy!55, line width=1pt] (FL) -- (FR);
  \node[fill=white, draw=Navy, line width=1.2pt, rounded corners=8pt, font=\bfseries\small,
        text=Navy, inner xsep=11pt, inner ysep=4pt] at ($(current page.south)+(0,2.18cm)$) {{\color{tmBlue}\faGlobe}\ \ thinkalgo.com};
  \node[text=Slate, font=\footnotesize, anchor=north west, align=left, inner sep=0pt, text width=7.6cm]
        at ($(FL)+(0.1cm,-0.30cm)$) {{\color{tmBlue}\faTag}\ \TopicWord: \textbf{\color{Ink}\CurTopic}\\[1.5pt]{\color{tmBlue}\faGraduationCap}\ \DocGrade};
  \node[text=Slate, font=\footnotesize] at ($(current page.south)+(0,1.20cm)$) {\thepage};
  \node[text=Slate, font=\scriptsize, anchor=north east, align=right, inner sep=0pt, text width=7.4cm]
        at ($(FR)+(-0.1cm,-0.34cm)$) {{\normalsize\bfseries\color{Navy}ThinkAlgo}\ — From Math Thinking to Code};
\end{tikzpicture}}

% ---------------- page-1 lead-in (topic is already in the header; keep a slim spacer) ----------------
\newcommand{\titlehead}{\vspace*{-0.2cm}}
\providecommand{\CurTopic}{}

% chapter banner + sets the running-header topic. #1 = number, #2 = title
\newcommand{\chapterhead}[2]{\renewcommand{\CurTopic}{#2}%
  \setcounter{seccnt}{0}% each chapter is independent: section badges restart at 1
  \phantomsection\pdfbookmark[0]{#1. #2}{chap#1}\label{chap:#1}% add this chapter to the PDF outline + TOC ref
  \par\vspace{0.1cm}\noindent\begin{tikzpicture}[baseline=-0.5ex]
    \node[circle,fill=tmBlue,text=white,font=\Large\bfseries,minimum size=0.95cm,inner sep=0pt](cn){#1};
    \node[anchor=west,font=\LARGE\bfseries,text=tmBlue] at ([xshift=12pt]cn.east){#2};
  \end{tikzpicture}\par\vspace{4pt}\noindent\textcolor{tmBlue}{\rule{\linewidth}{2pt}}\par\vspace{0.5cm}}

% ---------------- section heading (outline navy number badge, no fill) ----------------
\newcounter{seccnt}
\newcommand{\shead}[1]{\par\vspace{0.7em}\refstepcounter{seccnt}%
  \noindent\begin{tikzpicture}[baseline=-0.6ex]
    \node[circle,draw=Gold,line width=1.2pt,text=Gold,font=\bfseries,minimum size=1.05em,inner sep=1pt]{\theseccnt};
  \end{tikzpicture}\hspace{0.4em}{\large\bfseries\color{Navy} #1}\par\vspace{0.25em}}

% ============ boxes EXACTLY per the original main (36).tex: solid-blue PROBLEM, solid-orange EXAMPLE,
%              outline-green SOLUTION, all with drop shadow + rounded corners ============
% ---------------- knowledge box: white, tmBlue border + left accent + drop shadow ----------------
\newtcolorbox{knowbox}{enhanced,breakable,colback=white,colframe=tmBlue!70,boxrule=0.9pt,arc=5pt,
  left=12pt,right=11pt,top=6pt,bottom=6pt,drop shadow,
  borderline west={3.5pt}{0pt}{tmBlue}}

% ---------------- EXAMPLE (Ví dụ): SOLID ORANGE tab, white text ----------------
\newtcolorbox{exbox}[1]{enhanced,breakable,colback=white,colframe=tmOrange,boxrule=1pt,arc=5pt,
  left=12pt,right=11pt,top=14pt,bottom=7pt,drop shadow,
  attach boxed title to top left={xshift=10pt,yshift=-\tcboxedtitleheight/2},
  boxed title style={colback=tmOrange,colframe=tmOrange,arc=4pt,boxrule=0pt,drop shadow},
  title={#1},fonttitle=\bfseries\color{white}}
\newcommand{\guide}{\par\vspace{2pt}{\bfseries\color{tmBlue}\GuideWord}\par\nopagebreak}

% ---------------- PROBLEM (Bài): SOLID BLUE tab, white text ----------------
\newtcolorbox{probbox}[2]{enhanced,breakable,colback=white,colframe=tmBlue,boxrule=1pt,arc=5pt,
  left=12pt,right=10pt,top=14pt,bottom=7pt,drop shadow,
  attach boxed title to top left={xshift=10pt,yshift=-\tcboxedtitleheight/2},
  boxed title style={colback=tmBlue,colframe=tmBlue,arc=4pt,boxrule=0pt,drop shadow},
  fonttitle=\bfseries\color{white},title={\ProbWord\ #1\ifx\relax#2\relax\else#2\fi}}

% ---------------- SOLUTION (Lời giải): OUTLINE GREEN tab (white fill, green border + text) ----------------
% Teacher: a green tcolorbox. Student: the body is dropped entirely (comment pkg — robust with tables/\\).
\ifsol
  \newtcolorbox{solbox}{enhanced,breakable,colback=white,colframe=tmGreen,boxrule=1pt,arc=5pt,
    left=12pt,right=10pt,top=14pt,bottom=7pt,drop shadow,
    attach boxed title to top left={xshift=10pt,yshift=-\tcboxedtitleheight/2},
    boxed title style={colback=white,colframe=tmGreen,boxrule=1pt,arc=4pt,drop shadow},
    fonttitle=\bfseries\color{tmGreen!88!black},title={\SolTitle},before upper={\color{NavyD}}}
\else
  \excludecomment{solbox}
\fi

% ---------------- answer highlight: label green, value RED ----------------
\newcommand{\ans}[1]{\par\smallskip{\bfseries\color{tmGreen!88!black}\AnsWord:}\ {\bfseries\color{AnsRed}#1}\par}

% ---------------- digit cards (exact original style: fill=c!10, border=c!60!black, drop shadow) ----------------
% \digitcards{0/gray,1/cyan,2/green,3/cyan} : value/colour pairs
\newcommand{\digitcards}[1]{\begin{center}\begin{tikzpicture}[scale=0.7,transform shape,font=\sffamily]
  \def\dx{0}\foreach \v/\c in {#1}{
    \node[draw=\c!60!black,line width=1.5pt,fill=\c!10,rounded corners=4pt,
          minimum width=1cm,minimum height=1.4cm,drop shadow] at (\dx,0){\Huge\bfseries \v};
    \pgfmathsetmacro\dx{\dx+1.5}\xdef\dx{\dx}}
\end{tikzpicture}\end{center}}

% ---------------- student working space (white, light border) ----------------
\newcommand{\worksp}[1]{\ifsol\else\begin{tcolorbox}[colback=white,colframe=Line,boxrule=0.8pt,
  arc=5pt,height=\dimexpr#1+0.05cm\relax,left=8pt,right=8pt,top=4pt,bottom=4pt]{\footnotesize\color{Slate}\itshape\WorkWord}\end{tcolorbox}\fi}

% ---------------- \probsol{solution} : Teacher shows the green solution box; Student shows a BLANK
%   "Bài làm" box whose height matches that solution's height (writing room = work needed), so both
%   versions have identical per-problem height and paginate the same. ----------------
\newsavebox{\solmeasure}\newlength{\solht}
\newcommand{\probsol}[1]{%
  \begin{lrbox}{\solmeasure}\begin{minipage}{\dimexpr\linewidth-24pt\relax}#1\end{minipage}\end{lrbox}%
  \ifsol
    \begin{solbox}#1\end{solbox}%
  \else
    \setlength{\solht}{\dimexpr\ht\solmeasure+\dp\solmeasure+0.85cm\relax}%
    \ifdim\solht<1.35cm\setlength{\solht}{1.35cm}\fi
    \ifdim\solht>\StudentMaxWork\setlength{\solht}{\StudentMaxWork}\fi
    \begin{tcolorbox}[colback=white,colframe=Line,boxrule=0.8pt,arc=5pt,
      left=8pt,right=8pt,top=4pt,bottom=4pt,height=\solht]{\footnotesize\color{Slate}\itshape\WorkWord}\end{tcolorbox}%
  \fi}

% ---------------- practice divider: blue left accent bar, no fill ----------------
\newcommand{\dividerbar}[1]{\par\vspace{0.6em}\noindent
  \begin{tcolorbox}[enhanced,frame hidden,interior hidden,borderline west={5pt}{0pt}{ProbBlue},
    colback=white,left=12pt,right=12pt,top=5pt,bottom=5pt,arc=4pt]
  {\large\bfseries\color{ProbBlue} #1}\end{tcolorbox}\vspace{0.4em}}

% place-value choice-count diagram (Cambridge style): #1 = "label/count/note/color" list, #2 = result
% notes must NOT contain commas or slashes (they separate \foreach fields)
\newcommand{\countboxes}[2]{\begin{center}\begin{tikzpicture}[scale=0.62,transform shape,font=\sffamily]
  \def\xx{0}%
  \foreach \lab/\cnt/\note/\col [count=\i] in {#1}{%
    \ifnum\i>1\relax \node[text=gray,font=\Huge] at (\xx-1.8,0){$\times$};\fi
    \node[draw=\col!65,fill=\col!8,rounded corners=3pt,line width=1pt,minimum width=2.7cm,minimum height=1.7cm] at (\xx,0){};
    \node[text=\col!85,font=\bfseries] at (\xx,0.58){\lab};
    \node[text=\col!88,font=\Huge\bfseries] at (\xx,-0.1){\cnt};
    \node[text=gray,font=\scriptsize] at (\xx,-1.05){\note};
    \pgfmathsetmacro\xx{\xx+3.6}\xdef\xx{\xx}}%
  \node[text=gray,font=\Huge] at (\xx-1.8,0){$=$};
  \node[text=red,font=\Huge\bfseries] at (\xx,0){#2};
\end{tikzpicture}\end{center}}

% ======== bar model (sơ đồ đoạn thẳng): lightly-rounded blocks, darker border, tiny gaps, bright colours ========
% one block: \blk{x1}{x2}{yc}{height}{colour}{label}  (a small gap is left on each side)
\newcommand{\blk}[6]{%
  \draw[rounded corners=1.3pt,fill=#5!16,draw=#5!80!black,line width=1pt]
    (#1+0.045,#3-#4/2) rectangle (#2-0.045,#3+#4/2);
  \node[font=\small\bfseries,text=#5!52!black] at ({(#1+#2)/2},#3){#6};}
% measurement with a double-headed ARROW (not a brace) + label above the span: \brcup{x1}{x2}{y}{label}
\newcommand{\brcup}[4]{%
  \draw[<->,>=stealth,tmBlue!72!black,line width=0.7pt] (#1+0.05,#3)--(#2-0.05,#3);
  \node[above=-1.5pt,font=\footnotesize,text=tmBlue!58!black] at ({(#1+#2)/2},#3){#4};}
% measurement with a double-headed ARROW + label below the span: \brcdn{x1}{x2}{y}{label}
\newcommand{\brcdn}[4]{%
  \draw[<->,>=stealth,tmBlue!72!black,line width=0.7pt] (#1+0.05,#3)--(#2-0.05,#3);
  \node[below=-1.5pt,font=\footnotesize,text=tmBlue!58!black] at ({(#1+#2)/2},#3){#4};}
% row label at the left of a bar: \rowlab{x}{y}{text}
\newcommand{\rowlab}[3]{\node[anchor=east,font=\small\bfseries,text=Ink] at (#1,#2){#3};}
% vertical right-pointing CURLY BRACE spanning two stacked bars: \vbrace{x}{ytop}{ybot}{label}
\newcommand{\vbrace}[4]{%
  \pgfmathsetmacro\ym{(#2+#3)/2}%
  \draw[tmBlue!72!black,line width=0.9pt,line cap=round]
    (#1,#2) .. controls (#1+0.18,#2) and (#1+0.18,#2-0.18) .. (#1+0.18,#2-0.22)
    -- (#1+0.18,{\ym+0.22})
    .. controls (#1+0.18,\ym) and (#1+0.30,{\ym+0.08}) .. (#1+0.42,\ym)
    .. controls (#1+0.30,{\ym-0.08}) and (#1+0.18,\ym) .. (#1+0.18,{\ym-0.22})
    -- (#1+0.18,#3+0.22)
    .. controls (#1+0.18,#3+0.18) and (#1+0.18,#3) .. (#1,#3);
  \node[right=4pt,font=\footnotesize,text=tmBlue!58!black] at (#1+0.42,\ym){#4};}
% blank 3x3 grid for a problem statement (student fills it in)
\newcommand{\blankgrid}{\begin{tikzpicture}[font=\sffamily]
  \foreach \r in {0,1,2}{\foreach \c in {0,1,2}{
    \draw[rounded corners=1.3pt,fill=Tint,draw=Slate!55,line width=0.9pt]
      (\c+0.05,\r+0.05) rectangle (\c+0.95,\r+0.95);}}
\end{tikzpicture}}
% wrapper: LEFT-aligned (not centred), a little air above/below
\newenvironment{barfig}{\par\vspace{3pt}\noindent\kern0.35cm\begin{tikzpicture}[font=\sffamily,>=stealth]}%
  {\end{tikzpicture}\par\vspace{3pt}}
% wrapper: CENTRED (for illustrative diagrams in theory / examples)
\newenvironment{barfigC}{\par\vspace{3pt}\begin{center}\begin{tikzpicture}[font=\sffamily,>=stealth]}%
  {\end{tikzpicture}\end{center}\vspace{2pt}}

% ---------------- difficulty stars (nicely drawn) : \stars{1}=hard, \stars{2}=very hard ----------------
\newcommand{\diffstar}{\tikz[baseline=-0.55ex]{\node[star,star points=5,star point ratio=0.46,
  draw=StarGold!72!black,fill=StarGold,line width=0.3pt,inner sep=1.75pt]{};}}
\newcommand{\stars}[1]{\ifcase#1\or\kern2pt\diffstar\or\kern2pt\diffstar\kern1.5pt\diffstar\fi}

\newcommand{\nd}[1]{\overline{#1}}

% number line: #1 last index, #2 ticks "i/label", #3 open (excluded) indices, #4 filled (solution) indices
\newcommand{\numlinex}[4]{\begin{center}\begin{tikzpicture}[x=1.55cm,>=stealth]
  \draw[->,thick,Navy](-0.45,0)--(#1+0.45,0);
  \foreach \i/\lab in {#2}{\draw[Navy,line width=0.8pt](\i,0.08)--(\i,-0.08);
     \node[below,font=\small,text=Ink]at(\i,-0.1){\lab};}
  \foreach \i in {#3}{\draw[Navy,fill=white,line width=1pt](\i,0)circle(2.3pt);}
  \foreach \i in {#4}{\fill[KidOrange](\i,0)circle(2.7pt);}
\end{tikzpicture}\end{center}}


% ======================================================================
% Grade 7 additions (shared by all four editions)
% ======================================================================
% student edition: blank "work" box height = min(solution height, this cap)
\newlength{\StudentMaxWork}\setlength{\StudentMaxWork}{7.5cm}

% ThinkAlgo tip box (label set per language by \TipWord)
\newtcolorbox{tipbox}{enhanced,breakable,colback=StarGold!8,colframe=StarGold!75!black,boxrule=1pt,arc=4pt,
  left=11pt,right=10pt,top=6pt,bottom=6pt,
  before upper={\textbf{\color{StarGold!65!black}\faLightbulb\ \TipWord}\ }}

% "Recall" box: re-states a piece of theory inside a solution / example
\newtcolorbox{recallbox}{enhanced,breakable,colback=tmBlue!5,colframe=tmBlue!45,boxrule=0.6pt,arc=3pt,
  left=8pt,right=8pt,top=4pt,bottom=4pt,
  before upper={\textbf{\color{tmBlue!80!black}\faBookOpen\ \RecallWord}\ }}

% "Why is it true? / Proof" box
\newtcolorbox{proofbox}{enhanced,breakable,colback=tmPurple!4,colframe=tmPurple!45,boxrule=0.6pt,arc=3pt,
  left=8pt,right=8pt,top=4pt,bottom=4pt,
  before upper={\textbf{\color{tmPurple!85!black}\faSearch\ \ProofWord}\ },
  after upper={\hfill$\square$}}

% "Common mistake" box
\newtcolorbox{warnbox}{enhanced,breakable,colback=AnsRed!4,colframe=AnsRed!45,boxrule=0.6pt,arc=3pt,
  left=8pt,right=8pt,top=4pt,bottom=4pt,
  before upper={\textbf{\color{AnsRed!85!black}\faExclamationTriangle\ \WarnWord}\ }}

% "Note for parents" box (short hint on how to explain the idea at home)
\newtcolorbox{parentbox}{enhanced,breakable,colback=tmTeal!5,colframe=tmTeal!50,boxrule=0.6pt,arc=3pt,
  left=8pt,right=8pt,top=4pt,bottom=4pt,
  before upper={\textbf{\color{tmTeal!80!black}\faUsers\ \ParentWord}\ }}

% numbered step inside a solution:  \step{1}{Short title}  -> "Bước 1. Short title"
\newcommand{\step}[2]{\par\smallskip\noindent{\bfseries\color{tmBlue!85!black}\StepWord~#1.}\ {\bfseries #2}\par\nopagebreak}
% verification line:  \verify{...}
\newcommand{\verify}[1]{\par\smallskip\noindent{\bfseries\color{tmGreen!80!black}\faCheckCircle\ \CheckWord}\ #1\par}

% multiple-choice options
% \choices{..}{..}{..}{..}{..}      : five short options on one line
% \choicesII{..}{..}{..}{..}{..}    : two/three per line (medium length)
% \choicesL{..}{..}{..}{..}{..}     : one option per line (long text)
\newcommand{\choices}[5]{\par\smallskip\noindent
  \begin{tabularx}{\linewidth}{@{}XXXXX@{}}(A)\ #1&(B)\ #2&(C)\ #3&(D)\ #4&(E)\ #5\end{tabularx}\par}
\newcommand{\choicesII}[5]{\par\smallskip\noindent
  \begin{tabularx}{\linewidth}{@{}XXX@{}}(A)\ #1&(B)\ #2&(C)\ #3\\(D)\ #4&(E)\ #5&\end{tabularx}\par}
\newcommand{\choicesL}[5]{\par\smallskip\noindent
  \begin{tabular}{@{}l@{}}(A)\ #1\\(B)\ #2\\(C)\ #3\\(D)\ #4\\(E)\ #5\end{tabular}\par}

% centred figure inside a problem / solution
\newenvironment{figC}{\par\vspace{2pt}\begin{center}\begin{tikzpicture}[font=\sffamily\small,>=stealth,line join=round]}%
  {\end{tikzpicture}\end{center}\vspace{1pt}}

\def\DocGrade{Grade 7}
\def\ProbWord{Problem}
\def\SolWord{Solution.}
\def\GuideWord{Guided solution.}
\def\PageWord{Page}
\def\NameWord{Name}
\def\DateWord{Date}
\def\SolTitle{Detailed solution}
\def\WorkWord{Your work:}
\def\TopicWord{Topic}
\def\AnsWord{Answer}
\def\TipWord{ThinkAlgo Tip.}
\def\RecallWord{Recall.}
\def\ProofWord{Why is it true?}
\def\WarnWord{Common mistake.}
\def\ParentWord{Note for parents.}
\def\StepWord{Step}
\def\CheckWord{Check:}

\def\DocTopic{Number Sense}
\begin{document}
\titlehead
\chapterhead{4}{Number Sense}

% ---------------- part a ----------------
% =====================================================================
% Topic 4 — Number Sense — part a: theory, worked examples, problems 1–15
% =====================================================================
\begin{knowbox}
\textbf{Big idea.} ``Number sense'' is the ability to \emph{understand} numbers, not just to \emph{compute} with them:
knowing which number is larger, whether a result is reasonable, how to round for a quick estimate, and seeing the
hidden structure inside a number (its digits, its divisors, its multiples). In this topic you will review and connect:
\begin{itemize}[leftmargin=1.4em,itemsep=1pt,topsep=2pt]
  \item \textbf{decimals and negative numbers}: comparing and ordering them on the number line;
  \item the \textbf{order of operations} and \textbf{estimating} results;
  \item \textbf{fractions, ratios, percentages}: ``a fraction of a number'', equivalent fractions, discounts;
  \item \textbf{speed} -- distance -- time and unit conversion (km/h, m/s);
  \item \textbf{divisibility, common divisors, last digits, place value} (solving digit puzzles);
  \item the \textbf{mean, the median} and \textbf{number patterns} (Fibonacci, arithmetic sequences).
\end{itemize}
The general secret: \emph{estimate before you calculate; check after you calculate}. A student with good number
sense can often eliminate several unreasonable options in a multiple-choice question without any exact calculation.
\end{knowbox}
\vspace{0.3cm}

% ---------------------------------------------------------------------
\shead{Decimals: reading, comparing, ordering, adding and subtracting}
\begin{knowbox}
\textbf{1. Place value.} In a decimal, each digit sits in a \emph{place}. After the decimal point come the tenths,
the hundredths, the thousandths, \ldots
\begin{center}
\renewcommand{\arraystretch}{1.25}
\begin{tabular}{|c|c|c|c|c|c|}
\hline
\rowcolor{tmBlue!10} Tens & Ones & . & Tenths & Hundredths & Thousandths\\ \hline
 & $2$ & . & $2$ & $3$ & \\ \hline
 & $2$ & . & $3$ & $0$ & \\ \hline
\end{tabular}
\end{center}
Example: $2.23 = 2 + \dfrac{2}{10} + \dfrac{3}{100}$, while $2.3 = 2 + \dfrac{3}{10} = 2.30$.

\textbf{2. Adding or removing zeros at the end of the decimal part does not change the value:}
$2.3 = 2.30 = 2.300$. Reason: $\dfrac{3}{10} = \dfrac{30}{100} = \dfrac{300}{1000}$ (multiply numerator and denominator by $10$).

\textbf{3. Rule for comparing two positive decimals.}
\begin{itemize}[leftmargin=1.4em,itemsep=1pt,topsep=2pt]
  \item Compare the \emph{whole-number parts} first: the one with the larger whole part is larger.
  \item If the whole parts are equal, \emph{add zeros} so that both decimal parts have the same number of digits,
  then compare place by place from left to right: tenths, then hundredths, \ldots
\end{itemize}
Example: compare $0.304$ and $0.34$. Write $0.34 = 0.340$. Tenths: $3 = 3$; hundredths: $0 < 4$.
So $0.304 < 0.340 = 0.34$.

\textbf{4. Adding and subtracting decimals:} write the numbers in a column \emph{with the decimal points lined up}
(places lined up), then add/subtract as with whole numbers.

\textbf{5. Rounding decimals.} To round to the nearest hundredth (two digits after the point), look at the digit just
to the right of that place (the thousandths digit): if it is less than $5$, keep the hundredths digit; if it is $5$ or more,
add $1$ to the hundredths digit; then drop all digits after it. Example: $4.6875 \approx 4.69$ (the thousandths digit is
$7 \ge 5$), while $3.1416 \approx 3.14$ (because $1 < 5$). The sign ``$\approx$'' reads ``is approximately''.
\end{knowbox}
\begin{proofbox}
Why can't we compare ``by length'' (the number with more digits is larger)? Write both in thousandths:
$0.304 = \dfrac{304}{1000}$ and $0.34 = \dfrac{340}{1000}$. The fractions have the same denominator and $304 < 340$, so
$0.304 < 0.34$, even though $0.304$ ``looks longer''. Adding zeros is exactly bringing both to the same denominator.
\end{proofbox}
\begin{warnbox}
Many students think $2.23 > 2.3$ because ``$23 > 3$''. Wrong! Compare place by place: $2.23$ has $2$ tenths,
while $2.3$ has $3$ tenths, so $2.23 < 2.3$.
\end{warnbox}
\begin{tipbox}
To find the smallest (largest) number in a set of decimals, write all of them in \emph{one column} with the decimal
points lined up, adding zeros to make them equally long. Then you just compare them like whole numbers.
\end{tipbox}

% ---------------------------------------------------------------------
\shead{Negative numbers and the number line}
\begin{knowbox}
\textbf{1. The number line.} The integers are placed on a straight line: $0$ in the middle, positive numbers to the
right, negative numbers to the left.
\numlinex{8}{0/$-4$,1/$-3$,2/$-2$,3/$-1$,4/$0$,5/$1$,6/$2$,7/$3$,8/$4$}{}{4}
\textbf{2. Comparison rule.} On a (horizontal) number line, \emph{a number to the left is smaller than a number to the
right}. Therefore:
\begin{itemize}[leftmargin=1.4em,itemsep=1pt,topsep=2pt]
  \item every negative number is less than $0$, and $0$ is less than every positive number;
  \item every negative number is less than every positive number;
  \item of two negative numbers, the one \emph{farther from $0$} is \emph{smaller}: $-8 < -2$ because $8 > 2$.
\end{itemize}
\textbf{3. Everyday pictures.} A temperature of $-8^\circ$C is colder than $-2^\circ$C; owing $8$ thousand dong
leaves you ``poorer'' than owing $2$ thousand dong.

\textbf{4. The coordinate plane.} A point is given by a pair of numbers $(x;y)$: $x$ is the horizontal coordinate
(moving right/left), $y$ is the vertical coordinate (moving up/down). A horizontal move changes only $x$; a vertical
move changes only $y$.
\end{knowbox}
\begin{proofbox}
Why is $-8 < -2$? We have $-2 = -8 + 6$: to go from $-8$ to $-2$ on the number line we must move $6$ units \emph{to the
right}. Moving right means increasing, so $-2$ is larger than $-8$.
\end{proofbox}
\begin{warnbox}
Do not compare negative numbers ``by the feel of positive numbers''. $-8$ is \emph{less} than $-2$, even though $8$ is
greater than $2$. When in doubt, sketch a quick number line.
\end{warnbox}
\begin{tipbox}
\textbf{Parity (even -- odd)} is a powerful ``number sense'' tool: even $+$ even $=$ even, odd $+$ odd $=$ even,
even $+$ odd $=$ odd; a product of odd numbers is always odd. For instance, the sum of $5$ odd numbers is odd, and the sum
of $4$ odd numbers is even. When a point moves on a coordinate grid, looking at the parity of its coordinates often tells
you at once whether a position can be reached.
\end{tipbox}

% ---------------------------------------------------------------------
\shead{Order of operations and estimation}
\begin{knowbox}
\textbf{1. Order of operations.}
\begin{enumerate}[leftmargin=1.6em,itemsep=1pt,topsep=2pt]
  \item \emph{Brackets} first: $(\;)$, then $[\;]$, then $\{\;\}$.
  \item Next, \emph{powers}.
  \item Then \emph{multiplication and division}, from left to right.
  \item Finally \emph{addition and subtraction}, from left to right.
\end{enumerate}
Example: $6\times 2 - 3 = 12 - 3 = 9$ (multiply first, then subtract), whereas $6\times(2-3) = 6\times(-1) = -6$.

\textbf{2. Rounding to estimate.} To know quickly ``roughly how much'' a result is, round the factors to the nearest
ten or hundred and compute mentally. Example: $99 \times 9 \approx 100 \times 9 = 900$.

\textbf{3. A quick-calculation trick:} $99 \times 9 = (100-1)\times 9 = 900 - 9 = 891$ (multiplication distributes
over subtraction).
\end{knowbox}
\begin{proofbox}
Why is $(100-1)\times 9 = 100\times 9 - 1\times 9$? Imagine $9$ rows of $100$ sweets each, but $1$ sweet is taken from
every row. The sweets left are $9$ rows $\times$ $99$ sweets, which also equals ($9\times 100$ sweets) minus (the $9$
sweets taken). This is the \emph{distributive property}: $a\times(b-c) = a\times b - a\times c$.
\end{proofbox}
\begin{warnbox}
Working ``from left to right regardless of the operation'' is wrong: $2 + 3\times 4 \ne 5\times 4$. The correct
value is $2 + 3\times 4 = 2 + 12 = 14$.
\end{warnbox}
\begin{tipbox}
In a multiple-choice question asking ``which number is closest to\ldots'', the options often differ by factors of $10$
($10$; $100$; $1000$; \ldots). Then a rough estimate is enough to choose correctly.
\end{tipbox}

% ---------------------------------------------------------------------
\shead{Fractions: a fraction of a number, equivalent fractions}
\begin{knowbox}
\textbf{1. A fraction is a part of a whole.} $\dfrac{a}{b}$ means: split the whole into $b$ equal parts and take $a$ of them.

\textbf{2. Finding a fraction of a number:} to find $\dfrac{m}{n}$ of $a$, compute $a \times \dfrac{m}{n}$
(split $a$ into $n$ parts and take $m$ of them). Example: $\dfrac{1}{3}$ of $396$ is $396 \div 3 = 132$.

\textbf{3. The part that is left.} If $\dfrac{m}{n}$ of a job is done, the part left is $1 - \dfrac{m}{n} = \dfrac{n-m}{n}$.

\textbf{4. Equivalent fractions (the basic property).} Multiplying (or dividing) the numerator and the denominator by the
same non-zero number gives an equal fraction:
\[ \frac{3}{7} = \frac{3\times 9}{7\times 9} = \frac{27}{63}, \qquad \frac{12}{36} = \frac{12\div 12}{36\div 12} = \frac{1}{3}. \]
\textbf{5. Simplifying a fraction:} divide the numerator and the denominator by their greatest common divisor to get a
fraction in lowest terms.

\textbf{6. ``Sum and ratio'' problems.} If two numbers are in the ratio $\dfrac{5}{8}$, we can think of the first number as
$5$ equal parts and the second as $8$ such parts. Knowing the sum (or the difference) gives the value of one part.

\textbf{7. A fraction of a fraction.} ``Half of a quarter'' is $\dfrac12\times\dfrac14 = \dfrac18$: multiply numerator by
numerator, denominator by denominator. This happens often when a shape is halved several times in a row: each halving
multiplies the area fraction by $\dfrac12$ again.

\textbf{8. Comparing fractions.} Bring them to a common denominator and compare numerators, or convert them to
decimals. Examples: $\dfrac23 = \dfrac{8}{12} > \dfrac{7}{12}$; \ $\dfrac35 = 0.6 < \dfrac34 = 0.75$.
\end{knowbox}
\begin{barfigC}
  \rowlab{-0.2}{0}{First number}
  \foreach \i in {0,...,4}{\blk{\i*0.9}{\i*0.9+0.9}{0}{0.55}{tmBlue}{}}
  \rowlab{-0.2}{-0.8}{Second number}
  \foreach \i in {0,...,7}{\blk{\i*0.9}{\i*0.9+0.9}{-0.8}{0.55}{tmOrange}{}}
  \vbrace{7.35}{0.3}{-1.1}{sum $= 13$ parts}
\end{barfigC}
\begin{proofbox}
Why does multiplying the numerator and denominator by the same number leave the fraction unchanged? $\dfrac{3}{7}$ is
``cut the cake into $7$ parts, take $3$''. If we cut \emph{each} part into $9$ smaller pieces, the cake has
$7\times 9 = 63$ pieces and our share has $3\times 9 = 27$ pieces. The amount of cake has not changed, so
$\dfrac{3}{7} = \dfrac{27}{63}$.
\end{proofbox}
\begin{warnbox}
Read the question carefully: does it ask for the part \emph{done} or the part \emph{left}? Many students correctly find
$\dfrac13$ of $396$ to be $132$ and circle it straight away, while the question asks for the pages \emph{left},
$396 - 132 = 264$.
\end{warnbox}
\begin{parentbox}
With fractions, a picture beats any explanation. Draw a strip (a chocolate bar) split into equal parts and shade the
part used; your child will see the part left by themselves. This is the familiar ``bar model'' from primary school.
\end{parentbox}

% ---------------------------------------------------------------------
\shead{Percentages, unit price and discounts}
\begin{knowbox}
\textbf{1. Percent means ``out of one hundred''.} $80\% = \dfrac{80}{100} = 0.8$.

\textbf{2. Finding $p\%$ of a number $a$:} compute $a \times \dfrac{p}{100}$. Example: $80\%$ of $50$ is
$50 \times \dfrac{80}{100} = 40$.

\textbf{3. Expressing $a$ as a percentage of $b$:} compute $\dfrac{a}{b}\times 100\%$. Example: $32$ correct answers out of
$40$ questions is $\dfrac{32}{40}\times 100\% = 80\%$.

\textbf{4. Unit price.} The unit price is the price of \emph{one} unit (one kilogram, one book, \ldots):
\[ \text{Total cost} = \text{unit price} \times \text{quantity}, \qquad \text{unit price} = \text{total cost} \div \text{quantity}. \]
\textbf{5. A discount of $p\%$.} New price $=$ old price $\times \left(1 - \dfrac{p}{100}\right)$. A $10\%$ discount
leaves $90\%$: multiply by $0.9$.

\textbf{6. Successive discounts.} Each time the price is reduced by $10\%$ \emph{of the current price}, we multiply by
$0.9$ once more. Example: a price of $100$ thousand dong after two reductions: $100 \to 90 \to 81$ (thousand dong),
\emph{not} $80$.
\end{knowbox}
\begin{proofbox}
Why are two $10\%$ reductions not the same as one $20\%$ reduction? The second reduction is $10\%$ of $90$ (that is $9$),
not $10\%$ of $100$ (that is $10$). Since the price has already dropped, the next reduction is smaller too: in total the
price drops by only $10 + 9 = 19$ thousand dong, i.e.\ $19\%$.
\end{proofbox}
\begin{tipbox}
Finding $10\%$ very fast: just divide by $10$ (move the decimal point one place to the left). From $10\%$ you get
$5\%$ (halve it), $20\%$ (double it), $80\% = 100\% - 20\%$, \ldots
\end{tipbox}

% ---------------------------------------------------------------------
\shead{Speed, distance, time and unit conversion}
\begin{knowbox}
\textbf{1. Three linked formulas.} For motion at a constant speed:
\[ s = v \times t, \qquad v = s \div t, \qquad t = s \div v, \]
where $s$ is the distance, $v$ the speed and $t$ the time.

\textbf{2. Who is faster?} The faster person covers a \emph{longer} distance in the \emph{same amount of time}. The
safest way: compute each person's speed in the same unit (for example m/s) and compare.

\textbf{3. Running laps.} Running $n$ laps of a track $L$ metres long means a distance of $n \times L$ metres.

\textbf{4. Converting km/h to m/s.} $1$ km $= 1000$ m and $1$ hour $= 3600$ seconds, so
\[ 1 \text{ km/h} = \frac{1000 \text{ m}}{3600 \text{ s}} = \frac{5}{18} \text{ m/s}, \qquad
   18 \text{ km/h} = 18\times\frac{5}{18} = 5 \text{ m/s}, \qquad 1 \text{ m/s} = 3.6 \text{ km/h}. \]
\end{knowbox}
\begin{figC}
  \draw[tmBlue,line width=1.2pt,rounded corners=8pt] (0,0) rectangle (3.6,1.6);
  \node[draw=tmOrange,fill=tmOrange!15,rounded corners=3pt,minimum width=1.6cm] at (1.8,0.8) {$s = v\times t$};
  \node[text=Ink,align=left,anchor=west] at (4.2,0.8) {Cover the quantity you want:\\ $v = s \div t$ \quad and \quad $t = s \div v$};
\end{figC}
\begin{warnbox}
The quantities must use \emph{matching units}: if the distance is in metres and the time in seconds, the speed comes
out in m/s; do not divide metres by hours. Before comparing two speeds, convert them to the same unit.
\end{warnbox}

% ---------------------------------------------------------------------
\shead{Divisibility, common divisors, last digits and place value}
\begin{knowbox}
\textbf{1. Multiples and divisors.} $a$ is a \emph{multiple} of $b$ (or $b$ is a \emph{divisor} of $a$) if
$a = b\times k$ for some whole number $k$. Example: $45 = 9\times 5$, so $45$ is a multiple of $9$.

\textbf{2. Divisibility rules.}
\begin{itemize}[leftmargin=1.4em,itemsep=1pt,topsep=2pt]
  \item by $2$: the last digit is $0, 2, 4, 6, 8$; \quad by $5$: the last digit is $0$ or $5$; \quad by $10$: the last digit is $0$;
  \item by $3$: the digit sum is divisible by $3$; \quad by $9$: \emph{the digit sum is divisible by $9$}.
\end{itemize}
\textbf{3. Common divisors.} A number that divides both $a$ and $b$ is a common divisor of $a$ and $b$. Every common divisor
is a divisor of the \emph{greatest common divisor} $\gcd(a,b)$. Example: $\gcd(40,72)=8$, so the common divisors of $40$
and $72$ are $1, 2, 4, 8$.

\textbf{4. The last digit of a product} depends only on the last digits of the factors:
$\ldots 3 \times \ldots 7$ ends in the same digit as $3\times 7 = 21$, namely $1$.

\textbf{5. Place value.} A number written with digits can be expanded by place:
\[ \nd{abc} = 100a + 10b + c, \qquad \nd{PQQ} = 100P + 10Q + Q = 100P + 11Q. \]
(The bar on top shows that this is the number with digits $a, b, c$, not the product $a\times b\times c$.)
\end{knowbox}
\begin{proofbox}
Why does the divisibility rule for $9$ work? Take the number $\nd{abc} = 100a + 10b + c$ and split it:
\[ 100a + 10b + c = (99a + 9b) + (a + b + c) = 9\times(11a + b) + (a+b+c). \]
The part $9\times(11a+b)$ is always divisible by $9$. So $\nd{abc}$ is divisible by $9$ exactly when the remaining part
$a+b+c$ (the digit sum) is divisible by $9$. For numbers with more digits the argument is the same, because
$10 = 9+1$, $100 = 99+1$, $1000 = 999+1$, \ldots
\end{proofbox}
\begin{tipbox}
For a ``digit puzzle'' (each letter stands for a digit) such as $\nd{PQQ}+\nd{PPQ}+\nd{QQQ}=876$, work \emph{from the ones
column to the left}, watching out for carries. The ones column usually has the fewest possibilities, so the first digit
is found very quickly. You can also use place value to turn the puzzle into one equation. When you have to count or find
all possibilities (sets of digits, outcomes of tossing coins, \ldots), \emph{list them systematically} (in increasing order,
or with a table/tree diagram) so that nothing is missed or repeated.
\end{tipbox}
\begin{knowbox}
\textbf{Systematic listing and chance.}
\begin{itemize}[leftmargin=1.4em,itemsep=1pt,topsep=2pt]
  \item To count the three-digit numbers whose digits have a given \emph{product}, find the sets of three digits (written in
  increasing order so as not to repeat), then count the arrangements of each set: three different digits can be arranged in
  $3\times 2\times 1 = 6$ ways; two equal digits and one other in $3$ ways. Example: product $12$ with the set $\{1;3;4\}$
  gives $6$ numbers $134, 143, 314, 341, 413, 431$; the set $\{2;2;3\}$ gives $3$ numbers $223, 232, 322$.
  \item When all outcomes are \emph{equally likely}, the probability of an event is
  \[ \frac{\text{number of favourable outcomes}}{\text{total number of possible outcomes}}. \]
  Example: tossing $2$ coins has $4$ equally likely outcomes: HH, HT, TH, TT (H: heads, T: tails). The probability that both
  coins show the same face is $\dfrac{2}{4} = \dfrac12$. Use a tree diagram to list them:
\end{itemize}
\begin{figC}
  \node[draw=Navy,circle,inner sep=1.5pt] (r) at (0,0) {\scriptsize start};
  \node[draw=tmBlue,fill=tmBlue!10,circle,inner sep=2pt] (a) at (2,0.8) {H};
  \node[draw=tmOrange,fill=tmOrange!10,circle,inner sep=2pt] (b) at (2,-0.8) {T};
  \foreach \n/\y/\l/\c in {aa/1.2/H/tmBlue,ab/0.4/T/tmOrange}{\node[draw=\c,fill=\c!10,circle,inner sep=2pt] (\n) at (4,\y) {\l};}
  \foreach \n/\y/\l/\c in {ba/-0.4/H/tmBlue,bb/-1.2/T/tmOrange}{\node[draw=\c,fill=\c!10,circle,inner sep=2pt] (\n) at (4,\y) {\l};}
  \draw[Slate] (r)--(a) (r)--(b) (a)--(aa) (a)--(ab) (b)--(ba) (b)--(bb);
  \node[right,font=\scriptsize] at (4.3,1.2) {HH \ (same face)};
  \node[right,font=\scriptsize] at (4.3,0.4) {HT};
  \node[right,font=\scriptsize] at (4.3,-0.4) {TH};
  \node[right,font=\scriptsize] at (4.3,-1.2) {TT \ (same face)};
\end{figC}
\end{knowbox}
\begin{warnbox}
``The digit sum is divisible by $9$'' is the rule for $9$; it does \emph{not} work for $4$, $6$ or $8$. And a number
divisible by $3$ is not necessarily divisible by $9$: $12$ is divisible by $3$ but not by $9$.
\end{warnbox}

% ---------------------------------------------------------------------
\shead{Number patterns, the mean and the median}
\begin{knowbox}
\textbf{1. The Fibonacci sequence.} $1, 1, 2, 3, 5, 8, 13, 21, 34, 55, \ldots$ Each number (from the third on) is
\emph{the sum of the two numbers just before it}: $2 = 1+1$, $3 = 1+2$, $5 = 2+3$, \ldots

\textbf{2. Arithmetic sequences (equally spaced).} Each number (from the second on) equals the previous number plus a fixed
number $d$ (the common difference). Example: $1, 4, 7, 10$ (with $d = 3$). The $n$-th number is $a_1 + (n-1)\times d$.

\textbf{3. The mean (average).} The mean of $n$ numbers $=$ (their sum) $\div n$. Hence
\[ \text{Sum} = \text{mean} \times \text{how many numbers}. \]
\textbf{4. The ``deviation'' method.} The deviations from the mean always add up to $0$: the ``excess'' of the numbers
above the mean exactly equals the ``shortfall'' of the numbers below the mean.

\textbf{5. The median.} Arrange the numbers in increasing order; the number in the middle is the median (when there is
an odd number of numbers). Example: the median of $42, 43, 46, 49, 50$ is $46$. With an even number of numbers, the median
is the mean of the two middle numbers: the median of $6, 7, 9, 10$ is $(7+9)\div 2 = 8$. Note: when a number is added or
changed, the ``middle'' position can move, so re-order the list.

\textbf{6. Consecutive integers.} For an odd number of consecutive integers, the mean is the middle number:
the mean of $19, 20, 21, 22, 23$ is $21$.
\end{knowbox}
\begin{proofbox}
Why do the deviations add up to $0$? Suppose $4$ numbers have mean $80$, so their sum is $4\times 80 = 320$. Write
each number as $80 + $ (its deviation). Adding up: $320 + $ (sum of deviations) $= 320$, so the sum of the deviations
must be $0$.
\end{proofbox}
\begin{figC}
  \draw[Slate!60] (-0.3,0)--(5.3,0);
  \node[left,font=\scriptsize,text=Slate] at (-0.3,0) {mean $=40$};
  \foreach \x/\v/\c in {0/38/tmOrange,1.5/43/tmGreen,3/41/tmGreen,4.5/38/tmOrange}{
    \pgfmathsetmacro\h{(\v-40)*0.3}
    \fill[\c!70] (\x,0) rectangle (\x+0.7,\h);
    \node[font=\scriptsize] at (\x+0.35,{\h+0.25*sign(\h+0.001)}) {$\v$};}
\end{figC}
\begin{center}\small The four numbers $38, 43, 41, 38$ have mean $40$. Deviations from $40$: $-2,\ +3,\ +1,\ -2$; sum $= 0$.\end{center}
\begin{parentbox}
You can illustrate the mean with stacks of coins: ``level out'' the tall stacks into the short ones until all stacks are
the same height. That common height is the mean. The coins moved out of the tall stacks exactly equal the coins added to
the short ones -- this is the idea ``the deviations add up to $0$''.
\end{parentbox}

% =====================================================================
% WORKED EXAMPLES
% =====================================================================
\par\vspace{0.7em}\noindent{\large\bfseries\color{tmOrange!90!black}\faPencil*\ \ Worked examples}\par\vspace{0.3em}

\begin{exbox}{Example 1}
Consider the set $\{0.5;\ 0.45;\ 0.405;\ 0.054;\ 0.54\}$.
\begin{enumerate}[label=\alph*),leftmargin=1.6em,itemsep=0pt,topsep=2pt]
  \item Arrange the numbers in increasing order.
  \item Find the sum of the largest and the smallest number.
\end{enumerate}
\guide
\step{1}{Write the numbers in a column, adding zeros to get three digits after the point}
Adding zeros at the end of the decimal part does not change the value, so:
\[ 0.5 = 0.500;\quad 0.45 = 0.450;\quad 0.405;\quad 0.054;\quad 0.54 = 0.540. \]
\step{2}{Compare them like whole numbers}
Dropping ``$0.$'', we compare $500;\ 450;\ 405;\ 54;\ 540$:
\[ 54 < 405 < 450 < 500 < 540. \]
So the increasing order is: $0.054 < 0.405 < 0.45 < 0.5 < 0.54$.
\step{3}{Add the largest and the smallest (decimal points lined up)}
\[ \begin{array}{r} 0.540\\ +\ 0.054\\ \hline 0.594 \end{array} \]
\verify{$0.594 - 0.54 = 0.054$. Correct.}
\ans{a) $0.054 < 0.405 < 0.45 < 0.5 < 0.54$;\quad b) $0.594$}
\end{exbox}

\begin{exbox}{Example 2}
Hoa cycles along a $45$ km route. In the morning she covers $\dfrac{2}{5}$ of the route. How many kilometres does she
still have to ride in the afternoon, and what fraction of the whole route is that?
\guide
\step{1}{Draw a bar model}
Split the route into $5$ equal parts; the morning covers $2$ parts, the afternoon $3$ parts.
\begin{barfigC}
  \foreach \i in {0,1}{\blk{\i*1.6}{\i*1.6+1.6}{0}{0.6}{tmBlue}{$9$ km}}
  \foreach \i in {2,3,4}{\blk{\i*1.6}{\i*1.6+1.6}{0}{0.6}{tmOrange}{$9$ km}}
  \brcup{0}{3.2}{0.5}{morning}
  \brcup{3.2}{8}{0.5}{afternoon}
  \brcdn{0}{8}{-0.5}{$45$ km}
\end{barfigC}
\step{2}{Find the value of one part}
One part is $45 \div 5 = 9$ (km).
\step{3}{Morning and afternoon distances}
Morning: $2\times 9 = 18$ (km). Afternoon: $45 - 18 = 27$ (km), which is also $3\times 9 = 27$ (km).
\step{4}{The fraction left}
$1 - \dfrac{2}{5} = \dfrac{5}{5} - \dfrac{2}{5} = \dfrac{3}{5}$. Check: $\dfrac{27}{45} = \dfrac{27\div 9}{45\div 9} = \dfrac{3}{5}$.
\ans{$27$ km, which is $\dfrac{3}{5}$ of the route}
\end{exbox}

\begin{exbox}{Example 3}
a) Evaluate $20 - 4\times 3 + 8 \div 2$ \quad and \quad $(20-4)\times(3+8) \div 2$.\\
b) Without calculating exactly, decide which number $101 \times 49$ is closest to: $50$; $500$; $5000$; $50\,000$? Then
compute it exactly with a quick method.
\guide
\step{1}{No brackets: multiply and divide first; add and subtract after}
\[ 20 - \underbrace{4\times 3}_{12} + \underbrace{8 \div 2}_{4} = 20 - 12 + 4 = 8 + 4 = 12. \]
(Subtraction and addition are done from left to right: $20 - 12 = 8$, then $8 + 4 = 12$.)
\step{2}{With brackets: brackets first}
\[ (20 - 4)\times(3+8) \div 2 = 16 \times 11 \div 2 = 176 \div 2 = 88. \]
\step{3}{Estimate}
$101 \approx 100$ and $49 \approx 50$, so $101 \times 49 \approx 100\times 50 = 5000$. Choose $5000$.
\step{4}{Exact value using the distributive property}
$101 \times 49 = (100 + 1)\times 49 = 4900 + 49 = 4949$, indeed very close to $5000$.
\ans{a) $12$ and $88$;\quad b) closest to $5000$; $101\times 49 = 4949$}
\end{exbox}

\begin{exbox}{Example 4}
a) Rice costs $18$ thousand dong per kilogram. Mum buys $15$ kg and gets a $10\%$ discount on the whole bill. How much does she pay?\\
b) A shirt costs $400$ thousand dong. After every week the shop reduces the \emph{current price} by $10\%$. After at least how
many weeks is the price of the shirt below $300$ thousand dong?
\guide
\step{1}{Total before the discount}
Total cost $=$ unit price $\times$ quantity $= 18 \times 15 = 270$ (thousand dong).
\step{2}{Subtract $10\%$}
$10\%$ of $270$ is $270 \div 10 = 27$ (thousand dong). Mum pays $270 - 27 = 243$ (thousand dong) $= 243\,000$ dong.
(Or: $90\%$ remains, i.e.\ $270 \times 0.9 = 243$.)
\step{3}{Successive discounts: multiply by $0.9$ every week}
\begin{center}
\begin{tabular}{|c|c|c|c|c|}
\hline
\rowcolor{tmBlue!10} Time & Start & After 1 week & After 2 weeks & After 3 weeks\\ \hline
Price (thousand dong) & $400$ & $360$ & $324$ & $291.6$\\ \hline
Reduction & -- & $40$ & $36$ & $32.4$\\ \hline
\end{tabular}
\end{center}
After $2$ weeks the price is $324 > 300$; after $3$ weeks it is $291.6 < 300$.
\begin{warnbox}
Do not think ``the price drops by $40$ thousand every week'' (because $10\%$ of $400$ is $40$). Each week the reduction is
$10\%$ of \emph{the price at that time}, so the reductions get smaller and smaller: $40$, then $36$, then $32.4$, \ldots
\end{warnbox}
\ans{a) $243\,000$ dong;\quad b) $3$ weeks}
\end{exbox}

\begin{exbox}{Example 5}
a) Bình cycles $3$ laps around a stadium with a perimeter of $400$ m in $4$ minutes. An cycles at $15$ km/h. Who is faster?\\
b) A car travels at $72$ km/h. How many metres does it travel in $3$ seconds?
\guide
\step{1}{Bình's speed (m/s)}
Distance: $3 \times 400 = 1200$ (m). Time: $4$ minutes $= 4\times 60 = 240$ (seconds).
\[ v_{\text{Bình}} = 1200 \div 240 = 5 \text{ (m/s)}. \]
\step{2}{Convert An's speed to m/s}
\begin{recallbox}$1$ km/h $= \dfrac{1000 \text{ m}}{3600 \text{ s}} = \dfrac{5}{18}$ m/s.\end{recallbox}
$v_{\text{An}} = 15 \times \dfrac{5}{18} = \dfrac{75}{18} \approx 4.17$ (m/s). Since $5 > 4.17$, \textbf{Bình is faster}.
(Another way: $5$ m/s $= 5\times 3.6 = 18$ km/h $> 15$ km/h.)
\step{3}{The car}
$72$ km/h $= 72 \times \dfrac{5}{18} = 20$ m/s. In $3$ seconds it travels $20 \times 3 = 60$ (m).
\ans{a) Bình is faster ($5$ m/s against about $4.17$ m/s);\quad b) $60$ m}
\end{exbox}

\begin{exbox}{Example 6}
a) Which of the numbers $72;\ 84;\ 117;\ 135$ are divisible by $9$?\\
b) Find the digit $x$ such that the number $\nd{3x5}$ is divisible by $9$.\\
c) How many whole numbers are divisors of both $24$ and $36$?
\guide
\step{1}{Use the rule for $9$: compute the digit sum}
$72$: $7+2 = 9$ (divisible by $9$). \quad $84$: $8 + 4 = 12$ (not divisible by $9$).\\
$117$: $1+1+7 = 9$ (divisible). \quad $135$: $1+3+5 = 9$ (divisible).\\
So $72$, $117$, $135$ are divisible by $9$; $84$ is not ($84 = 9\times 9 + 3$).
\step{2}{Find the digit $x$}
The digit sum of $\nd{3x5}$ is $3 + x + 5 = 8 + x$. Since $x$ is a digit, $8 \le 8 + x \le 17$. In this range only $9$ is
divisible by $9$, so $8 + x = 9$, i.e.\ $x = 1$. The number is $315 = 9 \times 35$.
\step{3}{Common divisors of $24$ and $36$}
Divisors of $24$: $1, 2, 3, 4, 6, 8, 12, 24$. Divisors of $36$: $1, 2, 3, 4, 6, 9, 12, 18, 36$.\\
The numbers in both lists: $1, 2, 3, 4, 6, 12$ -- these are exactly the divisors of $\gcd(24,36) = 12$.
\ans{a) $72;\ 117;\ 135$;\quad b) $x = 1$;\quad c) $6$ numbers}
\end{exbox}

\begin{exbox}{Example 7}
In the calculation $\nd{AB} + \nd{AB} + \nd{AB} = \nd{1AB}$, each letter stands for a digit. Find $A$ and $B$.
\guide
\textbf{Method 1 -- Place value.} We have $\nd{1AB} = 100 + \nd{AB}$. So
\[ 3\times\nd{AB} = 100 + \nd{AB} \;\Rightarrow\; 2\times\nd{AB} = 100 \;\Rightarrow\; \nd{AB} = 50. \]
So $A = 5$, $B = 0$.

\textbf{Method 2 -- Column by column.} Ones column: $B + B + B = 3B$ ends in the digit $B$, so $3B - B = 2B$ ends in $0$,
hence $B = 0$ or $B = 5$.
\begin{itemize}[leftmargin=1.4em,itemsep=1pt,topsep=1pt]
  \item If $B = 5$: $3\times 5 = 15$, write $5$ carry $1$. Tens column: $3A + 1 = \nd{1A} = 10 + A$, so $2A = 9$ -- no digit
  works. Rejected.
  \item If $B = 0$: no carry. Tens column: $3A = 10 + A$, so $2A = 10$, $A = 5$.
\end{itemize}
\verify{$50 + 50 + 50 = 150 = \nd{1AB}$ with $A = 5$, $B = 0$. Correct.}
\ans{$A = 5$, $B = 0$}
\end{exbox}

\begin{exbox}{Example 8}
a) Nam's four test scores are $6$; $9$; $10$ and $x$. The mean of the four scores is $8$. Find $x$.\\
b) The mean of five consecutive whole numbers is $30$. Find the smallest of them.
\guide
\step{1}{Using the sum}
\begin{recallbox}Sum $=$ mean $\times$ how many numbers.\end{recallbox}
The four scores add up to $8 \times 4 = 32$. So $x = 32 - 6 - 9 - 10 = 7$.
\step{2}{Using deviations (faster)}
Compared with $8$: the $6$ is $2$ short, the $9$ is $1$ over, the $10$ is $2$ over. Sum of deviations: $-2 + 1 + 2 = +1$.
For all deviations to add up to $0$, the fourth score must be $1$ short: $x = 8 - 1 = 7$.
\step{3}{Five consecutive numbers}
For five consecutive numbers, the mean is the middle number. The middle number is $30$, so the five numbers are
$28, 29, 30, 31, 32$. The smallest is $28$.
\verify{$28 + 29 + 30 + 31 + 32 = 150 = 5\times 30$. Correct.}
\ans{a) $x = 7$;\quad b) $28$}
\end{exbox}

% =====================================================================
\clearpage
\dividerbar{Core practice}

% ---------------------------------------------------------------- Problem 1
\begin{probbox}{1}{\relax}
The smallest number in the set $\{3.2;\ 2.3;\ 3;\ 2.23;\ 3.22\}$ is
\choices{$3.2$}{$2.3$}{$3$}{$2.23$}{$3.22$}
\end{probbox}
\probsol{
\textbf{Understanding the problem.} There are $5$ decimals. We need the \emph{smallest} one.
\begin{recallbox}Comparing decimals: compare the whole parts first; if they are equal, compare the tenths, then the
hundredths, \ldots\ You may add zeros at the end of the decimal part without changing the value.\end{recallbox}
\step{1}{Compare the whole parts}
Whole parts: $3.2 \to 3$; \ $2.3 \to 2$; \ $3 \to 3$; \ $2.23 \to 2$; \ $3.22 \to 3$.\\
The smallest number must have the smallest whole part, namely $2$. Only two candidates remain: $2.3$ and $2.23$.
\step{2}{Use the same number of decimal places and compare}
$2.3 = 2.30$ \ and \ $2.23$. Tenths: $3$ against $2$; since $2 < 3$, we get $2.23 < 2.30$.
\step{3}{Conclusion}
Ordering everything to be sure: $2.23 < 2.3 < 3 < 3.2 < 3.22$. The smallest number is $2.23$.
\begin{warnbox}Do not choose $2.3$ just because ``$3$ is less than $23$''. The digit $3$ in $2.3$ means $3$ tenths, while
$2.23$ has only $2$ tenths.\end{warnbox}
\verify{$2.23 = \dfrac{223}{100}$ and $2.3 = \dfrac{230}{100}$; $223 < 230$.}
The number $2.23$ is option \textbf{(D)}.
\ans{(D) $2.23$}
}

% ---------------------------------------------------------------- Problem 2
\begin{probbox}{2}{\relax}
Minh goes on a $36$ km kayak trip. Before lunch he paddles the first $12$ km. After lunch, what fraction of the whole
trip does Minh still have to paddle?
\choices{$\dfrac{1}{2}$}{$\dfrac{5}{6}$}{$\dfrac{3}{4}$}{$\dfrac{2}{3}$}{$\dfrac{3}{5}$}
\end{probbox}
\probsol{
\textbf{Understanding the problem.} The whole trip is $36$ km and $12$ km are done. We are asked for the part \emph{left}
(not yet paddled), as a fraction.
\step{1}{Distance left}
$36 - 12 = 24$ (km).
\step{2}{Form the fraction ``part left over the whole''}
\begin{recallbox}The fraction ``what part of the whole'' $=\dfrac{\text{the part}}{\text{the whole}}$. Simplify by dividing
numerator and denominator by their greatest common divisor.\end{recallbox}
\[ \frac{24}{36} = \frac{24 \div 12}{36 \div 12} = \frac{2}{3}. \]
\step{3}{Picture it}
Split the $36$ km into $3$ sections of $12$ km: Minh has done $1$ section, $2$ sections are left.
\begin{barfigC}
  \blk{0}{2.4}{0}{0.6}{tmBlue}{$12$ km}
  \blk{2.4}{4.8}{0}{0.6}{tmOrange}{$12$ km}
  \blk{4.8}{7.2}{0}{0.6}{tmOrange}{$12$ km}
  \brcup{0}{2.4}{0.5}{before lunch}
  \brcup{2.4}{7.2}{0.5}{after lunch: $\frac{2}{3}$}
\end{barfigC}
\begin{warnbox}The fraction $\dfrac{12}{36} = \dfrac{1}{3}$ is the part \emph{done}. The question asks for the part
\emph{left}, so the answer is $1 - \dfrac13 = \dfrac23$.\end{warnbox}
\verify{$\dfrac{2}{3}\times 36 = 24$ km; $24 + 12 = 36$ km. Correct.}
\begin{parentbox}Ask your child two questions: ``What part is done?'' and ``What part is left?''. Getting into the habit of
answering both helps avoid the trap of choosing $\dfrac13$.\end{parentbox}
The fraction $\dfrac{2}{3}$ is option \textbf{(D)}.
\ans{(D) $\dfrac{2}{3}$}
}

% ---------------------------------------------------------------- Problem 3
\begin{probbox}{3}{\relax}
Which of the following statements is \textbf{true}?
\choicesII{$0$ is less than $-5$}{$7$ is less than $-1$}{$10$ is less than $\dfrac{1}{4}$}{$-1$ is less than $-3$}{$-8$ is less than $-2$}
\end{probbox}
\probsol{
\textbf{Understanding the problem.} Each option is a comparison. We check them one by one to find the true one.
\begin{recallbox}On a horizontal number line, a number to the left is smaller than a number to the right. Every negative
number is less than $0$ and less than every positive number. Of two negative numbers, the one farther from $0$ is smaller.\end{recallbox}
\step{1}{Place the numbers on a number line}
\begin{figC}
  \draw[->,thick,Navy] (-9.4*0.6,0)--(10.6*0.6,0);
  \foreach \x in {-9,...,10}{\draw[Navy] (\x*0.6,0.08)--(\x*0.6,-0.08);}
  \foreach \x in {-8,-5,-3,-2,-1,0,7,10}{\node[below,font=\scriptsize] at (\x*0.6,-0.1) {$\x$};}
  \foreach \x in {-8,-2}{\fill[tmOrange] (\x*0.6,0) circle (3pt);}
  \draw[->,tmGreen!70!black,thick] (-8*0.6,0.35) to[bend left=30] node[above,font=\scriptsize]{$6$ units to the right} (-2*0.6,0.35);
\end{figC}
\step{2}{Check each option}
(A) $0$ is less than $-5$? False: $-5$ is to the left of $0$, so $-5 < 0$.\\
(B) $7$ is less than $-1$? False: a positive number is always greater than a negative number.\\
(C) $10$ is less than $\dfrac14$? False: $\dfrac14 = 0.25 < 10$.\\
(D) $-1$ is less than $-3$? False: $-3$ is farther from $0$, so $-3 < -1$.\\
(E) $-8$ is less than $-2$? \textbf{True}: $-8$ is to the left of $-2$ (farther from $0$).
\begin{warnbox}Many students choose (D), thinking ``$1 < 3$ so $-1 < -3$''. For negative numbers it is the other way round:
$-1 > -3$ (a temperature of $-1^\circ$C is warmer than $-3^\circ$C).\end{warnbox}
\verify{$-2 - (-8) = 6 > 0$, which means $-2$ is greater than $-8$.}
The only true statement is \textbf{(E)}.
\ans{(E) $-8$ is less than $-2$}
}

% ---------------------------------------------------------------- Problem 4
\begin{probbox}{4}{\relax}
The value of $6 \times 2 - 3$ is
\choices{$9$}{$-6$}{$12$}{$15$}{$10$}
\end{probbox}
\probsol{
\textbf{Understanding the problem.} The expression has two operations, multiplication and subtraction, and no brackets.
\begin{recallbox}In an expression without brackets: multiply and divide first; add and subtract after.\end{recallbox}
\step{1}{Do the multiplication first}
$6\times 2 = 12$.
\step{2}{Then the subtraction}
$12 - 3 = 9$.
\begin{warnbox}Subtracting first would give $6\times(2-3) = 6\times(-1) = -6$ -- option (B), a ``trap'' for anyone who
forgets the order of operations. The expression has \emph{no brackets}, so we must multiply first.\end{warnbox}
\verify{Read the expression as a story: ``$6$ boxes with $2$ cakes each, $3$ cakes eaten'' leaves $12 - 3 = 9$ cakes.}
The value $9$ is option \textbf{(A)}.
\ans{(A) $9$}
}

% ---------------------------------------------------------------- Problem 5
\begin{probbox}{5}{\relax}
A test is marked out of $50$. Mai scores $80\%$. How many marks does Mai get?
\choices{$40$}{$62.5$}{$10$}{$45$}{$35$}
\end{probbox}
\probsol{
\textbf{Understanding the problem.} The maximum is $50$ marks; Mai gets $80\%$ of $50$. We need Mai's marks.
\begin{recallbox}To find $p\%$ of a number $a$, compute $a \times \dfrac{p}{100}$.\end{recallbox}
\step{1}{Write the calculation}
$80\%$ of $50$ is $50 \times \dfrac{80}{100}$.
\step{2}{Calculate}
$50 \times \dfrac{80}{100} = \dfrac{50\times 80}{100} = \dfrac{4000}{100} = 40$.
\step{3}{A quick mental method}
$10\%$ of $50$ is $5$; so $80\% = 8\times 10\%$ is $8\times 5 = 40$.
\begin{warnbox}Option (B) $62.5$ comes from dividing $50 \div 0.8$ -- wrong, since $62.5$ is even more than the maximum $50$!
Option (C) $10$ is the number of marks Mai \emph{lost} ($20\%$ of $50$).\end{warnbox}
\verify{$\dfrac{40}{50} = \dfrac{80}{100} = 80\%$. Correct.}
The value $40$ is option \textbf{(A)}.
\ans{(A) $40$}
}

% ---------------------------------------------------------------- Problem 6
\begin{probbox}{6}{\relax}
Minced beef sells for $120\,000$ dong per kilogram. How much do $12$ kg of minced beef cost?
\choicesII{$120\,000$ dong}{$12\,000$ dong}{$1\,440\,000$ dong}{$132\,000$ dong}{$10\,000$ dong}
\end{probbox}
\probsol{
\textbf{Understanding the problem.} The unit price (price of $1$ kg) is $120\,000$ dong; the quantity is $12$ kg. We need the total cost.
\begin{recallbox}Total cost $=$ unit price $\times$ quantity.\end{recallbox}
\step{1}{Write the calculation}
Total cost $= 120\,000 \times 12$ (dong).
\step{2}{Calculate (split $12 = 10 + 2$)}
$120\,000\times 10 = 1\,200\,000$; \quad $120\,000 \times 2 = 240\,000$.\\
$120\,000 \times 12 = 1\,200\,000 + 240\,000 = 1\,440\,000$ (dong).
\step{3}{Estimate to check}
$12$ kg at more than $100$ thousand dong per kg must cost more than $1\,200\,000$ dong. Only option (C) fits.
\begin{warnbox}The wrong options come from using the wrong operation: (A) is the price of $1$ kg; (D) is $120\,000 + 12\,000$
(adding instead of multiplying); (E) is $120\,000 \div 12$ (dividing instead of multiplying).\end{warnbox}
\verify{$1\,440\,000 \div 12 = 120\,000$ dong per kg. Correct.}
The amount $1\,440\,000$ dong is option \textbf{(C)}.
\ans{(C) $1\,440\,000$ dong}
}

% ---------------------------------------------------------------- Problem 7
\begin{probbox}{7}{\relax}
The value of $(8 \times 4) + 3$ is
\choices{$96$}{$15$}{$56$}{$35$}{$28$}
\end{probbox}
\probsol{
\textbf{Understanding the problem.} The brackets contain a multiplication; then $3$ is added.
\begin{recallbox}Expressions with brackets: work out the brackets first.\end{recallbox}
\step{1}{Work out the brackets}
$8\times 4 = 32$.
\step{2}{Add $3$}
$32 + 3 = 35$.
\begin{warnbox}Misreading it as $8 \times (4 + 3) = 8\times 7 = 56$ leads to option (C). The brackets surround $8\times 4$,
not $4 + 3$. Option (A) $96 = 8\times 4\times 3$ comes from mistaking the plus sign for a times sign.\end{warnbox}
\verify{Without brackets the result is still $8\times 4 + 3 = 35$, because multiplication always comes before addition; here the
brackets only add emphasis.}
The value $35$ is option \textbf{(D)}.
\ans{(D) $35$}
}

% ---------------------------------------------------------------- Problem 8
\begin{probbox}{8}{\relax}
Which of the following integers is closest to $99 \times 9$?
\choices{$10\,000$}{$100$}{$100\,000$}{$1000$}{$10$}
\end{probbox}
\probsol{
\textbf{Understanding the problem.} The options are $10$; $100$; $1000$; $10\,000$; $100\,000$ -- each $10$ times the
previous. An estimate of $99\times 9$ is enough, but we also compute it exactly to be sure.
\step{1}{Estimate}
$99 \approx 100$, so $99 \times 9 \approx 100\times 9 = 900$.
\step{2}{Exact value using the distributive property}
\begin{recallbox}Distributive property: $(a - b)\times c = a\times c - b\times c$.\end{recallbox}
$99\times 9 = (100 - 1)\times 9 = 900 - 9 = 891$.
\step{3}{Compare the distances to the nearest options}
The distance from $891$ to $1000$ is $1000 - 891 = 109$.\\
The distance from $891$ to $100$ is $891 - 100 = 791$.\\
The other options ($10$; $10\,000$; $100\,000$) are much farther away. So $1000$ is the closest.
\begin{figC}
  \draw[->,thick,Navy] (-0.3,0)--(10.6,0);
  \foreach \x/\l in {0/0,1/100,10/1000}{\draw[Navy] (\x,0.1)--(\x,-0.1); \node[below,font=\scriptsize] at (\x,-0.1) {$\l$};}
  \fill[tmOrange] (8.91,0) circle (3pt); \node[above,font=\scriptsize] at (8.91,0.1) {$891$};
  \draw[<->,tmGreen!70!black] (8.91,-0.55)--(10,-0.55) node[midway,below,font=\scriptsize]{$109$};
  \draw[<->,AnsRed!70] (1,-0.55)--(8.91,-0.55) node[midway,below,font=\scriptsize]{$791$};
\end{figC}
\verify{$891 \div 9 = 99$. Correct.}
The number $1000$ is option \textbf{(D)}.
\ans{(D) $1000$}
}

% ---------------------------------------------------------------- Problem 9
\begin{probbox}{9}{\relax}
Huy is reading a $396$-page book. He has read only the first $\dfrac{1}{3}$ of the book. How many pages does he still
have to read to finish the book?
\choices{$264$}{$124$}{$250$}{$199$}{$244$}
\end{probbox}
\probsol{
\textbf{Understanding the problem.} The whole book: $396$ pages. Read so far: $\dfrac13$ of the book. Asked: the number of
pages \emph{left}.
\begin{recallbox}To find $\dfrac{m}{n}$ of a number, divide it by $n$ and multiply by $m$.\end{recallbox}
\step{1}{Pages read}
$\dfrac13$ of $396$ is $396 \div 3 = 132$ (pages).
\step{2}{Pages left}
$396 - 132 = 264$ (pages).
\step{3}{Another way: compute the part left directly}
The part left is $1 - \dfrac13 = \dfrac23$ of the book, i.e.\ $396 \div 3 \times 2 = 132\times 2 = 264$ (pages).
\begin{barfigC}
  \blk{0}{2.6}{0}{0.6}{tmBlue}{$132$ pages}
  \blk{2.6}{5.2}{0}{0.6}{tmOrange}{$132$ pages}
  \blk{5.2}{7.8}{0}{0.6}{tmOrange}{$132$ pages}
  \brcup{0}{2.6}{0.5}{read}
  \brcup{2.6}{7.8}{0.5}{still to read}
  \brcdn{0}{7.8}{-0.5}{$396$ pages}
\end{barfigC}
\begin{warnbox}Stopping at $132$ answers the wrong question (that is the number of pages \emph{read}; and $132$ is not even
among the options). Always re-read the question before choosing an answer.\end{warnbox}
\verify{$264 + 132 = 396$, and $132\times 3 = 396$. Correct.}
The number $264$ is option \textbf{(A)}.
\ans{(A) $264$}
}

% ---------------------------------------------------------------- Problem 10
\begin{probbox}{10}{\relax}
Which of the following numbers is a multiple of $9$?
\choices{$50$}{$40$}{$35$}{$45$}{$55$}
\end{probbox}
\probsol{
\textbf{Understanding the problem.} We need the number divisible by $9$ (a multiple of $9$) among the five given numbers.
\begin{recallbox}A number is divisible by $9$ exactly when its digit sum is divisible by $9$.\end{recallbox}
\step{1}{Digit sum of each number}
\begin{center}
\begin{tabular}{|c|c|c|c|c|c|}
\hline
\rowcolor{tmBlue!10} Number & $50$ & $40$ & $35$ & $45$ & $55$\\ \hline
Digit sum & $5$ & $4$ & $8$ & $9$ & $10$\\ \hline
Divisible by $9$? & no & no & no & \textbf{yes} & no\\ \hline
\end{tabular}
\end{center}
\step{2}{Check with the $9$ times table}
Multiples of $9$ nearby: $36,\ 45,\ 54$. Only $45 = 9\times 5$ appears among the options.
\begin{proofbox}Why does the rule work for a two-digit number $\nd{ab}$? We have $\nd{ab} = 10a + b = 9a + (a + b)$. The part $9a$
is divisible by $9$, so $\nd{ab}$ is divisible by $9$ exactly when $a + b$ is. For $45$: $45 = 9\times 4 + 9$.\end{proofbox}
\verify{$45 \div 9 = 5$, remainder $0$.}
The number $45$ is option \textbf{(D)}.
\ans{(D) $45$}
}

% ---------------------------------------------------------------- Problem 11
\begin{probbox}{11}{\relax}
Consider the set $\{0.34;\ 0.304;\ 0.034;\ 0.43\}$. The sum of the smallest and the largest numbers in the set is
\choices{$0.77$}{$0.734$}{$0.077$}{$0.464$}{$0.338$}
\end{probbox}
\probsol{
\textbf{Understanding the problem.} Part 1: find the smallest and the largest number. Part 2: add them.
\begin{recallbox}To compare decimals, add zeros so the decimal parts have equal length, then compare place by place from
left to right.\end{recallbox}
\step{1}{Write the numbers with three decimal places}
\begin{center}
\begin{tabular}{|c|c|c|c|}
\hline
\rowcolor{tmBlue!10} Original & Rewritten & Tenths & Hundredths \\ \hline
$0.34$ & $0.340$ & $3$ & $4$ \\ \hline
$0.304$ & $0.304$ & $3$ & $0$ \\ \hline
$0.034$ & $0.034$ & $0$ & $3$ \\ \hline
$0.43$ & $0.430$ & $4$ & $3$ \\ \hline
\end{tabular}
\end{center}
\step{2}{Find the smallest and the largest}
Comparing like whole numbers: $34 < 304 < 340 < 430$, so $0.034 < 0.304 < 0.34 < 0.43$.\\
Smallest: $0.034$ (its tenths digit is $0$). Largest: $0.43$ (its tenths digit is $4$).
\step{3}{Add, with the decimal points lined up}
\[ \begin{array}{r} 0.430\\ +\ 0.034\\ \hline 0.464 \end{array} \]
\begin{warnbox}Misaligning the columns (adding $43$ and $34$ ``by the last digits'') gives $0.77$ or $0.077$ -- options (A)
and (C). Keep the decimal points in line: $4$ tenths plus $0$ tenths, $3$ hundredths plus $3$ hundredths, \ldots\end{warnbox}
\verify{$0.464 - 0.43 = 0.034$. Correct.}
The sum $0.464$ is option \textbf{(D)}.
\ans{(D) $0.464$}
}

% ---------------------------------------------------------------- Problem 12
\begin{probbox}{12}{\relax}
Lan and Hà run on a circular track $150$ m long. Lan runs $3$ laps in $120$ seconds, and Hà runs $5$ laps in $160$
seconds. Who is the faster runner, and approximately what is her speed?
\choicesII{Lan, $3.75$ m/s}{Lan, $2.4$ m/s}{Hà, $3.3$ m/s}{Hà, $4.69$ m/s}{Hà, $3.75$ m/s}
\end{probbox}
\probsol{
\textbf{Understanding the problem.} One lap is $150$ m. Lan: $3$ laps in $120$ s. Hà: $5$ laps in $160$ s. We compute
each girl's speed (m/s), compare, and choose the option.
\begin{recallbox}Speed $=$ distance $\div$ time. The distance of $n$ laps $= n\times$ the length of one lap.\end{recallbox}
\step{1}{Distance run by each girl}
Lan: $3\times 150 = 450$ (m). \quad Hà: $5\times 150 = 750$ (m).
\step{2}{Lan's speed}
$v_{\text{Lan}} = 450 \div 120 = 3.75$ (m/s).
\step{3}{Hà's speed}
$v_{\text{Hà}} = 750 \div 160 = 4.6875 \approx 4.69$ (m/s).
\step{4}{Compare}
$4.69 > 3.75$, so \textbf{Hà is faster}, running at about $4.69$ m/s.
\begin{center}
\begin{tabular}{|c|c|c|c|}
\hline
\rowcolor{tmBlue!10} & Distance & Time & Speed\\ \hline
Lan & $450$ m & $120$ s & $3.75$ m/s\\ \hline
Hà & $750$ m & $160$ s & $\approx 4.69$ m/s\\ \hline
\end{tabular}
\end{center}
\begin{warnbox}Option (E) ``Hà, $3.75$ m/s'' pairs the right name with the slower girl's speed. Option (B) $2.4$ m/s comes from
the reversed calculation $120 \div 50$ (time divided by distance). Always write the units: m $\div$ s $=$ m/s.\end{warnbox}
\verify{$4.6875 \times 160 = 750$ m $= 5$ laps; $3.75\times 120 = 450$ m $= 3$ laps. Correct.}
The conclusion ``Hà, $4.69$ m/s'' is option \textbf{(D)}.
\ans{(D) Hà, $4.69$ m/s}
}

% ---------------------------------------------------------------- Problem 13
\begin{probbox}{13}{\relax}
In the addition shown, $P$ and $Q$ each represent a digit. The value of $P + Q$ is
\[ \begin{array}{r@{\;}c@{\;}c@{\;}c} & P & Q & Q\\ & P & P & Q\\ + & Q & Q & Q\\ \hline & 8 & 7 & 6 \end{array} \]
\choices{$3$}{$5$}{$7$}{$6$}{$4$}
\end{probbox}
\probsol{
\textbf{Understanding the problem.} The three numbers $\nd{PQQ}$, $\nd{PPQ}$, $\nd{QQQ}$ add up to $876$. We need the digits $P$, $Q$,
then $P + Q$. We go column by column, starting with the ones.
\step{1}{Ones column: $Q + Q + Q = 3Q$ ends in $6$}
Since $Q$ is a digit, $3Q$ is one of $0, 3, 6, \ldots, 27$. Of these, only $6$ ends in $6$ ($16$ and $26$ are not multiples
of $3$). So $3Q = 6$, $Q = 2$, and there is \emph{no carry} into the tens.
\step{2}{Tens column: $Q + P + Q = P + 4$ ends in $7$}
$P + 4$ lies between $4$ and $13$; ending in $7$ means $P + 4 = 7$, so $P = 3$. No carry into the hundreds.
\step{3}{Hundreds column: check $P + P + Q = 8$}
$3 + 3 + 2 = 8$. This matches the hundreds digit of $876$.
\begin{recallbox}Place value: $\nd{PQQ} = 100P + 11Q$, \ $\nd{PPQ} = 110P + Q$, \ $\nd{QQQ} = 111Q$.\end{recallbox}
\step{4}{Another way using place value}
Adding: $(100P + 11Q) + (110P + Q) + 111Q = 210P + 123Q = 876$. Try $Q = 2$: $210P = 876 - 246 = 630$, so $P = 3$.
(Other values of $Q$ do not give a digit for $P$.)
\verify{$322 + 332 + 222 = 876$. Correct.}
\begin{parentbox}A digit puzzle is like a detective game: each column is a ``clue''. Let your child write each column as a
sentence (``three times $Q$ ends in $6$'') and try the digits $0$ to $9$. Do not forget to check the carries.\end{parentbox}
So $P + Q = 3 + 2 = 5$, which is option \textbf{(B)}.
\ans{(B) $5$}
}

% ---------------------------------------------------------------- Problem 14
\begin{probbox}{14}{\relax}
The mean (average) of the four integers $78$; $83$; $82$ and $x$ is $80$. Which of the following statements is true?
\choicesL{$x$ is $2$ greater than the mean}{$x$ is $1$ less than the mean}{$x$ is $2$ less than the mean}{$x$ is $3$ less than the mean}{$x$ is equal to the mean}
\end{probbox}
\probsol{
\textbf{Understanding the problem.} We know the mean of four numbers is $80$ and we know three of them. Find $x$ and compare it with $80$.
\begin{recallbox}Sum of the numbers $=$ mean $\times$ how many numbers.\end{recallbox}
\step{1}{Sum of the four numbers}
$80\times 4 = 320$.
\step{2}{Find $x$}
$78 + 83 + 82 = 243$, so $x = 320 - 243 = 77$.
\step{3}{Compare with the mean}
$80 - 77 = 3$: $x$ is $3$ less than the mean.
\step{4}{A faster way -- the deviation method}
Compared with $80$: $78$ is $2$ short; $83$ is $3$ over; $82$ is $2$ over. Sum of these three deviations: $-2 + 3 + 2 = +3$.
Since all the deviations must add up to $0$, $x$ must be $3$ ``short'': $x = 80 - 3 = 77$.
\begin{proofbox}Each number can be written as $80 + $(its deviation). The sum of the four numbers is $4\times 80 + $(sum of deviations).
But the sum of the four numbers is exactly $4\times 80$, so the sum of the deviations is $0$.\end{proofbox}
\verify{$(78 + 83 + 82 + 77) \div 4 = 320 \div 4 = 80$. Correct.}
\begin{parentbox}The ``deviation'' method is great for working out average marks mentally in everyday life. You can ask your
child: ``Three tests scored $7$, $9$, $8$; to average $8$, what does the fourth test need?'' (Answer: $8$, because
$-1 + 1 + 0 = 0$.)\end{parentbox}
The statement ``$x$ is $3$ less than the mean'' is option \textbf{(D)}.
\ans{(D)}
}

% ---------------------------------------------------------------- Problem 15
\begin{probbox}{15}{\relax}
In the Fibonacci sequence $1, 1, 2, 3, 5, \ldots$, each number from the $2$ onwards is the sum of the two numbers just
before it. For example, the next number in the sequence is $3 + 5 = 8$. Which of the following numbers is in the sequence?
\choices{$20$}{$21$}{$22$}{$23$}{$24$}
\end{probbox}
\probsol{
\textbf{Understanding the problem.} Rule: next number $=$ sum of the two before it. We continue the sequence until it
passes $24$ (the largest option), then see which option appears.
\step{1}{Continue the sequence}
\begin{center}
\begin{tabular}{|c|c|c|c|c|c|c|c|c|c|}
\hline
\rowcolor{tmBlue!10} Position & 1 & 2 & 3 & 4 & 5 & 6 & 7 & 8 & 9\\ \hline
Number & $1$ & $1$ & $2$ & $3$ & $5$ & $8$ & $13$ & $\mathbf{21}$ & $34$\\ \hline
\end{tabular}
\end{center}
In detail: $5 + 8 = 13$; \ $8 + 13 = 21$; \ $13 + 21 = 34$.
\step{2}{Compare with the options}
$34$ is already greater than $24$, so all later numbers are greater than $24$. Between $13$ and $34$ there is exactly one
number of the sequence, $21$. So $20$, $22$, $23$, $24$ are not in the sequence; only $21$ is.
\begin{tipbox}The Fibonacci sequence keeps increasing (from the second number on) because we keep adding positive numbers.
So once the sequence has passed a number, that number can never appear later.\end{tipbox}
\verify{$8 + 13 = 21$. Correct.}
The number $21$ is option \textbf{(B)}.
\ans{(B) $21$}
}

% ---------------- part b ----------------
% ============================================================
% Topic 4 -- Number Sense -- part b: problems 16--30 (EN)
% ============================================================

\begin{probbox}{16}{\relax}
The mean (average) of five consecutive integers is $21$. The smallest of the five integers is
\choices{$17$}{$21$}{$1$}{$18$}{$19$}
\end{probbox}
\probsol{
\textbf{Understanding the problem.} We have \emph{five consecutive integers} (each one is exactly $1$ more than the one before). Their mean is $21$. We must find the \emph{smallest} one.

\begin{recallbox}Mean of a group of numbers $=$ (sum of the numbers) $\div$ (how many numbers there are). So: sum $=$ mean $\times$ how many numbers.\end{recallbox}

\step{1}{Find the sum of the five numbers}
The sum of the five numbers is $21\times 5=105$.

\step{2}{Name the five numbers using the middle one}
Call the middle (third) number $n$. Then the five consecutive numbers are
\[n-2,\quad n-1,\quad n,\quad n+1,\quad n+2.\]
Adding: $(n-2)+(n-1)+n+(n+1)+(n+2)=5n$ (the $-2$ and $+2$, and the $-1$ and $+1$, cancel each other).

\step{3}{Find the middle number, then the smallest}
We have $5n=105$, so $n=105\div 5=21$. The five numbers are $19,\ 20,\ 21,\ 22,\ 23$. The smallest is $n-2=19$.

\begin{figC}
\draw[->,thick,Navy] (-0.5,0)--(6.6,0);
\foreach \x/\lab in {0/18,1/19,2/20,3/21,4/22,5/23,6/24}{\draw[Navy,line width=0.8pt] (\x,0.08)--(\x,-0.08); \node[below,text=Ink] at (\x,-0.1) {$\lab$};}
\foreach \x in {1,2,3,4,5}{\fill[KidOrange] (\x,0) circle (2.7pt);}
\draw[<->,tmBlue] (1,0.45)--(3,0.45) node[midway,above,font=\footnotesize]{back $2$};
\draw[<->,tmBlue] (3,0.45)--(5,0.45) node[midway,above,font=\footnotesize]{forward $2$};
\node[tmGreen!70!black,font=\footnotesize\bfseries] at (3,1.1) {middle number $=$ mean};
\end{figC}

\begin{proofbox}With an \emph{odd} number of consecutive integers, the numbers placed symmetrically around the middle one ``balance out'' (one is $1$ less, one is $1$ more; one is $2$ less, one is $2$ more\dots). So the sum equals (middle number) $\times$ (how many numbers), which means the mean is exactly the middle number.\end{proofbox}

\begin{warnbox}Do not choose $21$ straight away: that is the \emph{middle} number, not the smallest. And do not take $21-4=17$ either: from the middle number to the smallest is only $2$ steps back, not $4$.\end{warnbox}

\verify{$19+20+21+22+23=105$ and $105\div 5=21$. Correct.}
The smallest number is $19$, which is option \textbf{(E)}.
\ans{(E) $19$}
}

\begin{probbox}{17}{\relax}
How many positive whole numbers, including $1$, divide exactly into both $40$ and $72$?
\choices{$9$}{$12$}{$4$}{$2$}{$5$}
\end{probbox}
\probsol{
\textbf{Understanding the problem.} A number ``divides exactly into both $40$ and $72$'' when $40$ is divisible by it \emph{and} $72$ is divisible by it. These are the \emph{common divisors} (common factors) of $40$ and $72$. We must count them.

\begin{recallbox}$d$ is a divisor of $a$ if $a$ is divisible by $d$. A common divisor of two numbers is a divisor of both. The greatest common divisor (GCD) is the largest of the common divisors.\end{recallbox}

\step{1}{List the divisors of $40$}
Find them in pairs whose product is $40$: $1\times 40$, $2\times 20$, $4\times 10$, $5\times 8$. So the divisors of $40$ are
\[\{1,\ 2,\ 4,\ 5,\ 8,\ 10,\ 20,\ 40\}.\]

\step{2}{List the divisors of $72$}
Pairs with product $72$: $1\times 72$, $2\times 36$, $3\times 24$, $4\times 18$, $6\times 12$, $8\times 9$. So the divisors of $72$ are
\[\{1,\ 2,\ 3,\ 4,\ 6,\ 8,\ 9,\ 12,\ 18,\ 24,\ 36,\ 72\}.\]

\step{3}{Pick the numbers that appear in both lists}
Comparing the lists, the common numbers are $1,\ 2,\ 4,\ 8$. So there are $4$ common divisors.

\begin{tipbox}A faster way: the common divisors of two numbers are exactly the divisors of their GCD. We have $40=2^3\cdot 5$ and $72=2^3\cdot 3^2$, so $\text{GCD}(40,72)=2^3=8$. The divisors of $8$ are $1,2,4,8$: exactly $4$ numbers.\end{tipbox}

\begin{proofbox}If $d$ divides both $40$ and $72$, then $d$ also divides the difference $72-40=32$, and then $40-32=8$. Conversely, every divisor of $8$ divides $40=8\cdot 5$ and $72=8\cdot 9$. So ``common divisors of $40$ and $72$'' and ``divisors of $8$'' are the same set.\end{proofbox}

\begin{warnbox}Option $12$ is the number of divisors of $72$ alone; options $9$ or $5$ usually come from missing or miscounting some. The question asks for \emph{common} divisors.\end{warnbox}

\verify{$40\div1,\ 40\div2=20,\ 40\div4=10,\ 40\div8=5$; $72\div1,\ 72\div2=36,\ 72\div4=18,\ 72\div8=9$ -- all exact. But $3,6,9$ do not divide $40$, and $5,10$ do not divide $72$.}
There are $4$ such numbers, which is option \textbf{(C)}.
\ans{(C) $4$}
}

\begin{probbox}{18}{\relax}
The number that goes into the box to make $\dfrac{3}{7}=\dfrac{\square}{63}$ true is
\choices{$27$}{$9$}{$59$}{$63$}{$3$}
\end{probbox}
\probsol{
\textbf{Understanding the problem.} The two fractions are equal; the denominator has changed from $7$ to $63$. We need the new numerator.

\begin{recallbox}\textbf{Basic property of fractions:} multiplying both the numerator and the denominator by the same non-zero number gives an equal fraction: $\dfrac{a}{b}=\dfrac{a\cdot m}{b\cdot m}\ (m\neq 0)$.\end{recallbox}

\step{1}{What was the denominator multiplied by?}
$63\div 7=9$, so the denominator $7$ was multiplied by $9$ to become $63$.

\step{2}{Multiply the numerator by the same number}
To keep the fraction the same, the numerator must also be multiplied by $9$: $3\times 9=27$. So
\[\frac{3}{7}=\frac{3\cdot 9}{7\cdot 9}=\frac{27}{63}.\]

\begin{warnbox}Option $9$ is the \emph{multiplier}, not the new numerator. Option $59$ comes from the ``adding'' mistake: $63-7=56$ and then $3+56=59$ -- adding the same number to the numerator and denominator does \emph{not} keep the fraction the same.\end{warnbox}

\verify{Cross-multiply: $3\times 63=189$ and $7\times 27=189$. The products are equal, so the fractions are equal.}
The missing number is $27$, which is option \textbf{(A)}.
\ans{(A) $27$}
}

\begin{probbox}{19}{\relax}
A fraction is equal to $\dfrac{5}{8}$. Its numerator and denominator add up to $91$. What is the difference between the denominator and the numerator of this fraction?
\choices{$21$}{$3$}{$33$}{$13$}{$19$}
\end{probbox}
\probsol{
\textbf{Understanding the problem.} The fraction equals $\dfrac58$, so its numerator and denominator ``share'' in the ratio $5:8$. We know numerator $+$ denominator $=91$. We want denominator $-$ numerator.

\begin{recallbox}Every fraction equal to $\dfrac{5}{8}$ has the form $\dfrac{5k}{8k}$ for some positive integer $k$ (multiply the numerator and denominator by $k$).\end{recallbox}

\step{1}{Draw a bar model}
Think of the numerator as $5$ equal parts and the denominator as $8$ such parts.
\begin{barfig}
\rowlab{-0.15}{0.5}{Numerator}
\foreach \i in {0,...,4}{\blk{\i*0.8}{\i*0.8+0.8}{0.5}{0.5}{tmBlue}{}}
\rowlab{-0.15}{-0.3}{Denominator}
\foreach \i in {0,...,7}{\blk{\i*0.8}{\i*0.8+0.8}{-0.3}{0.5}{tmOrange}{}}
\brcup{0}{4}{0.9}{$5$ parts}
\brcdn{0}{6.4}{-0.7}{$8$ parts}
\vbrace{6.6}{0.75}{-0.55}{sum $=91$}
\end{barfig}

\step{2}{Find the value of one part}
There are $5+8=13$ parts in total, worth $91$. One part is $91\div 13=7$.

\step{3}{Find the numerator, the denominator and the difference}
Numerator $=5\times 7=35$; denominator $=8\times 7=56$. The fraction is $\dfrac{35}{56}$.\\
Denominator $-$ numerator $=56-35=21$ (which is also $3$ parts: $3\times 7=21$).

\begin{warnbox}Option $3$ is the difference $8-5$ of the original fraction $\dfrac58$, not yet multiplied by the value of one part. Option $13$ is the number of parts, not the difference.\end{warnbox}

\verify{$35+56=91$; simplifying, $\dfrac{35}{56}=\dfrac{35\div7}{56\div7}=\dfrac58$. Correct.}
The difference is $21$, which is option \textbf{(A)}.
\ans{(A) $21$}
}

\begin{probbox}{20}{\relax}
Lan is driving at $80$ km/h along a road with equally spaced power poles. Her car passes a pole every $4$ seconds. Which of the following is closest to the distance between two neighbouring poles?
\choices{$50$ m}{$60$ m}{$70$ m}{$80$ m}{$90$ m}
\end{probbox}
\probsol{
\textbf{Understanding the problem.} In $4$ seconds the car travels exactly the distance between two neighbouring poles. So we need the distance travelled in $4$ seconds at $80$ km/h. Careful: the speed is in \emph{km and hours}, the time is in \emph{seconds}, and the answers are in \emph{metres} -- we must convert units.

\begin{recallbox}Distance $=$ speed $\times$ time. $1$ km $=1000$ m; $1$ hour $=60$ minutes $=3600$ seconds.\end{recallbox}

\step{1}{Convert the speed to metres per second}
$80$ km/h means $80\,000$ m in $3600$ seconds. In $1$ second the car travels
\[80\,000\div 3600=\frac{80\,000}{3600}=\frac{200}{9}\approx 22.2\ \text{m}.\]

\step{2}{Distance travelled in $4$ seconds}
\[\frac{200}{9}\times 4=\frac{800}{9}\approx 88.9\ \text{m}.\]

\step{3}{Compare with the options}
$88.9$ m is about $1.1$ m from $90$ m and about $8.9$ m from $80$ m. So the closest is $90$ m.

\begin{tipbox}Another way to think: in $1$ hour $=3600$ seconds the car passes $3600\div 4=900$ gaps between poles. These $900$ gaps make up $80\,000$ m altogether, so each gap is $80\,000\div 900\approx 88.9$ m.\end{tipbox}

\begin{warnbox}If you forget to convert and compute $80\times 4=320$, no option matches. Converting $1$ hour as $60$ seconds (wrong -- it is $3600$ seconds) also gives an absurd result.\end{warnbox}

\verify{$90$ m every $4$ seconds is $22.5$ m/s, i.e.\ $22.5\times 3600=81\,000$ m per hour $=81$ km/h, very close to $80$ km/h. Reasonable.}
Choose option \textbf{(E)}.
\ans{(E) $90$ m}
}

\begin{probbox}{21}{\relax}
Each of the integers $334$ and $419$ has digits whose product is $36$ (since $3\cdot 3\cdot 4=36$ and $4\cdot 1\cdot 9=36$). How many $3$-digit positive integers have digits whose product is $36$?
\choices{$21$}{$15$}{$18$}{$24$}{$12$}
\end{probbox}
\probsol{
\textbf{Understanding the problem.} A $3$-digit number $\nd{abc}$ has digit product $a\cdot b\cdot c=36$. No digit can be $0$ (a $0$ would make the product $0$). A tidy method: \emph{first} find the sets of three digits (ignoring order) with product $36$, \emph{then} count how many numbers each set can make.

\step{1}{Find every set of three digits (from $1$ to $9$) with product $36$}
List each set in order $a\le b\le c$ and look at the smallest digit $a$.
\begin{itemize}[leftmargin=1.4em,itemsep=1pt]
\item $a=1$: we need $b\cdot c=36$ with $1\le b\le c\le 9$: $36=4\cdot 9=6\cdot 6$ (while $1\cdot 36$, $2\cdot 18$, $3\cdot 12$ use a number bigger than $9$). We get $\{1,4,9\}$ and $\{1,6,6\}$.
\item $a=2$: we need $b\cdot c=18$ with $2\le b\le c\le 9$: $18=2\cdot 9=3\cdot 6$. We get $\{2,2,9\}$ and $\{2,3,6\}$.
\item $a=3$: we need $b\cdot c=12$ with $3\le b\le c$: $12=3\cdot 4$. We get $\{3,3,4\}$.
\item $a\ge 4$: then $a\cdot b\cdot c\ge 4\cdot 4\cdot 4=64>36$, no sets.
\end{itemize}
So there are exactly $5$ sets: $\{1,4,9\},\ \{1,6,6\},\ \{2,2,9\},\ \{2,3,6\},\ \{3,3,4\}$.

\begin{recallbox}\textbf{Counting arrangements:} three \emph{different} digits can be arranged into $3\cdot 2\cdot 1=6$ numbers (hundreds digit: $3$ choices, tens: $2$ choices, units: $1$ choice). Three digits with \emph{two equal} can be arranged into only $3$ numbers (we just choose the position of the odd one out).\end{recallbox}

\step{2}{Count the numbers made from each set}
\begin{center}
\begin{tabular}{|c|c|l|c|}
\hline
\rowcolor{Tint}\textbf{Digit set} & \textbf{Type} & \textbf{Numbers formed} & \textbf{Count}\\
\hline
$\{1,4,9\}$ & all different & $149,194,419,491,914,941$ & $6$\\
$\{2,3,6\}$ & all different & $236,263,326,362,623,632$ & $6$\\
$\{1,6,6\}$ & two equal & $166,616,661$ & $3$\\
$\{2,2,9\}$ & two equal & $229,292,922$ & $3$\\
$\{3,3,4\}$ & two equal & $334,343,433$ & $3$\\
\hline
\multicolumn{3}{|r|}{\textbf{Total}} & $\mathbf{21}$\\
\hline
\end{tabular}
\end{center}

\step{3}{Add up}
$6+6+3+3+3=21$ numbers.

\begin{warnbox}Two common mistakes: (1) forgetting the sets containing the digit $1$ (like $\{1,4,9\}$, $\{1,6,6\}$) because you only think of splitting $36$ into three numbers ``bigger than $1$''; (2) assuming a set with two equal digits also gives $6$ numbers -- this double counts (for example, $334$ would be counted twice).\end{warnbox}

\verify{The table lists all $21$ numbers, including $334$ and $419$ from the question. A computer check (going through every number from $100$ to $999$) also gives exactly $21$.}
Choose option \textbf{(A)}.
\ans{(A) $21$}
}

\begin{probbox}{22}{\relax}
An \emph{arithmetic sequence} is a sequence in which each term after the first is obtained by adding a constant to the previous term. For example, $2,4,6,8$ and $1,4,7,10$ are arithmetic sequences.

In the grid shown, the numbers in each row (read from left to right) must form an arithmetic sequence, and the numbers in each column (read from top to bottom) must form an arithmetic sequence. The value of $x$ is
\begin{figC}
\foreach \rr in {0,1,2,3}{\foreach \cc in {0,1,2,3}{\draw[Navy,line width=0.9pt,fill=white] (\cc*1.1,-\rr*0.8) rectangle (\cc*1.1+1.1,-\rr*0.8-0.8);}}
\node at (0.55,-0.4) {$1$};
\node at (0.55,-1.2) {$4$}; \node at (3.85,-1.2) {$25$};
\node at (0.55,-2.0) {$7$}; \node at (3.85,-2.0) {$x$};
\node at (0.55,-2.8) {$10$}; \node at (2.75,-2.8) {$36$};
\end{figC}
\choices{$37$}{$28$}{$36$}{$43.75$}{$46$}
\end{probbox}
\probsol{
\textbf{Understanding the problem.} The grid has $4$ rows and $4$ columns. Every row and every column is an arithmetic sequence. We need $x$ in row $3$, column $4$. Idea: use the rows that already have two known numbers to find more numbers, then use column $4$ to find $x$.

\begin{recallbox}An arithmetic sequence with $4$ terms can be written $a,\ a+d,\ a+2d,\ a+3d$ ($d$ is the constant difference). From the first to the last term there are $3$ steps, so $d=(\text{last}-\text{first})\div 3$. If we know the first and the third terms, then $d=(\text{third}-\text{first})\div 2$.\end{recallbox}

\step{1}{Row $2$: $4,\ \_,\ \_,\ 25$}
From $4$ to $25$ there are $3$ steps: $d=(25-4)\div 3=21\div 3=7$. Row $2$ is $4,\ 11,\ 18,\ 25$.

\step{2}{Row $4$: $10,\ \_,\ 36,\ \_$}
From the first term $10$ to the third term $36$ there are $2$ steps: $d=(36-10)\div 2=26\div 2=13$. Row $4$ is $10,\ 23,\ 36,\ 49$. So the last cell of row $4$ is $49$.

\step{3}{Column $4$: $\_,\ 25,\ x,\ 49$}
In column $4$, $x$ lies between $25$ and $49$, the same distance from each. So $x$ is their mean:
\[x=\frac{25+49}{2}=\frac{74}{2}=37.\]
(The difference in column $4$ is $37-25=12$, so the top of the column is $25-12=13$.)

\step{4}{Fill in the whole grid to check}
\begin{figC}
\foreach \rr in {0,1,2,3}{\foreach \cc in {0,1,2,3}{\draw[Navy,line width=0.9pt,fill=white] (\cc*1.1,-\rr*0.8) rectangle (\cc*1.1+1.1,-\rr*0.8-0.8);}}
\foreach \cc/\val in {0/1,1/5,2/9,3/13}{\node at (\cc*1.1+0.55,-0.4) {$\val$};}
\foreach \cc/\val in {0/4,1/11,2/18,3/25}{\node at (\cc*1.1+0.55,-1.2) {$\val$};}
\foreach \cc/\val in {0/7,1/17,2/27}{\node at (\cc*1.1+0.55,-2.0) {$\val$};}
\node[text=AnsRed,font=\bfseries] at (3.85,-2.0) {$37$};
\foreach \cc/\val in {0/10,1/23,2/36,3/49}{\node at (\cc*1.1+0.55,-2.8) {$\val$};}
\node[anchor=west,font=\footnotesize,text=Slate] at (4.6,-0.4) {$+4$};
\node[anchor=west,font=\footnotesize,text=Slate] at (4.6,-1.2) {$+7$};
\node[anchor=west,font=\footnotesize,text=Slate] at (4.6,-2.0) {$+10$};
\node[anchor=west,font=\footnotesize,text=Slate] at (4.6,-2.8) {$+13$};
\end{figC}
Column $1$: $1,4,7,10$ (step $3$); column $2$: $5,11,17,23$ (step $6$); column $3$: $9,18,27,36$ (step $9$); column $4$: $13,25,37,49$ (step $12$). Every row and every column is arithmetic.

\begin{warnbox}You cannot find $x$ by looking at row $3$ first: row $3$ has only one known number ($7$), not enough to find its difference. You must go ``around'' through row $2$ and row $4$ to reach column $4$.\end{warnbox}

\verify{Row $3$: $7,17,27,37$ has constant difference $10$; column $4$: $25,37,49$ has constant difference $12$. Correct.}
$x=37$, which is option \textbf{(A)}.
\ans{(A) $37$}
}

\begin{probbox}{23}{\relax}
In the six-digit number $\nd{1ABCDE}$, each letter represents a digit. Given that $\nd{1ABCDE}\times 3=\nd{ABCDE1}$, the value of $A+B+C+D+E$ is
\choices{$29$}{$26$}{$22$}{$30$}{$28$}
\end{probbox}
\probsol{
\textbf{Understanding the problem.} When $\nd{1ABCDE}$ is multiplied by $3$, the leading digit $1$ ``moves'' to the end, and the block $ABCDE$ shifts one place to the left. We must find the digits $A,B,C,D,E$ and add them. There are two methods; both are well worth learning.

\begin{recallbox}\textbf{Place value:} writing a digit $c$ to the right of a number $N$ gives $10\cdot N+c$; writing the digit $1$ to the left of a five-digit number $N$ gives $100\,000+N$.\end{recallbox}

\textbf{Method 1: One unknown for the whole block.}
\step{1}{Let $N=\nd{ABCDE}$ (a five-digit number)}
Then $\nd{1ABCDE}=100\,000+N$ and $\nd{ABCDE1}=10N+1$.
\step{2}{Write the equation and find $N$}
\[3\cdot(100\,000+N)=10N+1\ \Longrightarrow\ 300\,000+3N=10N+1.\]
Subtract $3N$ and $1$ from both sides: $299\,999=7N$, so $N=299\,999\div 7=42\,857$.
\step{3}{Read off the digits}
$A=4,\ B=2,\ C=8,\ D=5,\ E=7$. Sum: $4+2+8+5+7=26$.

\textbf{Method 2: Multiply column by column, from right to left.}
\[\begin{array}{r@{\;}c@{\;}c@{\;}c@{\;}c@{\;}c@{\;}c}
 & 1 & A & B & C & D & E\\
\times & & & & & & 3\\ \hline
 & A & B & C & D & E & 1
\end{array}\]
\begin{itemize}[leftmargin=1.4em,itemsep=1pt]
\item Units: $3\cdot E$ ends in $1$ $\Rightarrow$ $E=7$ ($3\cdot 7=21$, write $1$ carry $2$).
\item Tens: $3\cdot D+2$ ends in $E=7$ $\Rightarrow$ $3D$ ends in $5$ $\Rightarrow$ $D=5$ ($15+2=17$, write $7$ carry $1$).
\item Hundreds: $3\cdot C+1$ ends in $D=5$ $\Rightarrow$ $3C$ ends in $4$ $\Rightarrow$ $C=8$ ($24+1=25$, write $5$ carry $2$).
\item Thousands: $3\cdot B+2$ ends in $C=8$ $\Rightarrow$ $3B$ ends in $6$ $\Rightarrow$ $B=2$ ($6+2=8$, carry $0$).
\item Ten thousands: $3\cdot A+0$ ends in $B=2$ $\Rightarrow$ $A=4$ ($12$, write $2$ carry $1$).
\item Hundred thousands: $3\cdot 1+1=4$, exactly $A=4$. The multiplication matches completely.
\end{itemize}
(At each step only one digit works, because $3\cdot 0, 3\cdot 1,\dots,3\cdot 9$ all have different last digits.)

\begin{tipbox}$142\,857$ is a famous number: $142\,857\times 2=285\,714$, $\times 3=428\,571$, $\times 4=571\,428$, $\times 5=714\,285$, $\times 6=857\,142$ -- always the digits $1,4,2,8,5,7$ ``rotated''. It is linked to $\dfrac17=0.\overline{142857}$.\end{tipbox}

\verify{$142\,857\times 3=428\,571$. This is indeed $\nd{ABCDE1}$ with $ABCDE=42857$.}
$A+B+C+D+E=26$, which is option \textbf{(B)}.
\ans{(B) $26$}
}

\begin{probbox}{24}{\relax}
On a coordinate grid, a dot starts at the point $(20,19)$. On its first move it goes $1$ unit horizontally or vertically, that is, to one of the points $(21,19)$, $(19,19)$, $(20,20)$ or $(20,18)$. On its second move it goes $2$ units in a direction perpendicular to the first move (either way). Each move after that is $1$ unit longer than the previous one ($3$ units, $4$ units, $5$ units,\dots) and must be perpendicular to the direction of the move just before it. The dot stops after $10$ moves. Which of the following final positions is \textbf{not} possible?
\choices{$(27,33)$}{$(30,40)$}{$(21,21)$}{$(42,44)$}{$(37,37)$}
\end{probbox}
\probsol{
\textbf{Understanding the problem.} ``Perpendicular to the previous move'' means the moves \emph{alternate}: horizontal, vertical, horizontal, vertical,\dots\ (or vertical, horizontal, vertical,\dots). The move lengths are $1,2,3,\dots,10$. We track separately the change in the $x$-coordinate (from horizontal moves) and in the $y$-coordinate (from vertical moves).

\begin{recallbox}For the point $(x,y)$: $x$ is the horizontal coordinate (moving right increases $x$, moving left decreases it), $y$ is the vertical coordinate (moving up increases $y$, moving down decreases it). A sum of an odd number of odd numbers is odd; a sum (or difference) of even numbers is always even.\end{recallbox}

\step{1}{Split the moves into two groups}
Moves $1,3,5,7,9$ are all in one direction (horizontal or vertical), with lengths $1,3,5,7,9$ (all odd). Moves $2,4,6,8,10$ are in the other direction, with lengths $2,4,6,8,10$ (all even).
\begin{itemize}[leftmargin=1.4em,itemsep=1pt]
\item If the first move is \emph{horizontal}: $x$ changes by $\pm1\pm3\pm5\pm7\pm9$ and $y$ changes by $\pm2\pm4\pm6\pm8\pm10$.
\item If the first move is \emph{vertical}: the other way round, $x$ changes by $\pm2\pm4\pm6\pm8\pm10$ and $y$ by $\pm1\pm3\pm5\pm7\pm9$.
\end{itemize}
(The sign $+$ or $-$ depends on which way the dot goes.)

\step{2}{Odd and even}
$\pm1\pm3\pm5\pm7\pm9$ combines $5$ odd numbers, so it is always \emph{odd}, and at most $25$. And $\pm2\pm4\pm6\pm8\pm10$ is always \emph{even}, at most $30$. So of the two changes (in $x$ and in $y$), exactly one is odd (from the odd group) and the other is even (from the even group).

\begin{proofbox}Start with the ``all plus'' total $1+3+5+7+9=25$. Switching the sign of a number $k$ from $+$ to $-$ lowers the total by $2k$. So every possible value is $25-2s$, where $s$ is the sum of the numbers whose sign was switched. Likewise, the even group gives the values $30-2s$.\end{proofbox}

\step{3}{Compute the changes for each option}
Subtract $(20,19)$ from the final point:
\begin{center}
\small\begin{tabularx}{\linewidth}{|c|c|c|X|}
\hline
\rowcolor{Tint}\textbf{Option} & \textbf{Change in $x$} & \textbf{Change in $y$} & \textbf{Can it be made?}\\
\hline
(A) $(27,33)$ & $7$ & $14$ & $1+3+5+7-9=7$; $2+4+6-8+10=14$ \\
(B) $(30,40)$ & $10$ & $21$ & we would need $\pm1\pm3\pm5\pm7\pm9=21$ \dots\\
(C) $(21,21)$ & $1$ & $2$ & $1-3+5+7-9=1$; $2-4+6+8-10=2$\\
(D) $(42,44)$ & $22$ & $25$ & $2-4+6+8+10=22$; $1+3+5+7+9=25$\\
(E) $(37,37)$ & $17$ & $18$ & $-1-3+5+7+9=17$; $2+4-6+8+10=18$\\
\hline
\end{tabularx}
\end{center}

\step{4}{Why (B) is impossible}
In (B) the odd change is $21$, so we would need $\pm1\pm3\pm5\pm7\pm9=21$, i.e.\ $25-2s=21$, so $s=2$: we would have to switch the signs of some odd numbers adding up to $2$. But a single odd number is never $2$ (it can only be $1,3,5,\dots$), and two different odd numbers of the group add up to at least $1+3=4$. There is no way! So the dot cannot stop at $(30,40)$.

\step{5}{An example path (option D)}
First move up: up $1$, right $2$, up $3$, left $4$, up $5$, right $6$, up $7$, right $8$, up $9$, right $10$. The dot goes from $(20,19)$ to $(42,44)$.
\begin{figC}\begin{scope}[scale=0.17]
\draw[Line,line width=0.3pt] (-3,-1) grid (23,26);
\draw[tmBlue,line width=1.4pt,->] (0,0)--(0,1)--(2,1)--(2,4)--(-2,4)--(-2,9)--(4,9)--(4,16)--(12,16)--(12,25)--(22,25);
\fill[tmGreen] (0,0) circle (10pt); \fill[AnsRed] (22,25) circle (10pt);
\node[below right,font=\scriptsize,text=tmGreen!70!black] at (0,0) {$(20,19)$};
\node[below,font=\scriptsize,text=AnsRed] at (22,24.6) {$(42,44)$};
\end{scope}\end{figC}

\begin{warnbox}Checking only ``is the point too far away?'' is not enough: $(30,40)$ is no farther than $(42,44)$. What decides it is odd/even and which totals can actually be made.\end{warnbox}

\verify{Every equation in the table can be recomputed, e.g.\ $-1-3+5+7+9=17$. A computer search over all $2\cdot 2^{10}=2048$ ways of moving also confirms that only $(30,40)$ cannot be reached.}
Choose option \textbf{(B)}.
\ans{(B) $(30,40)$}
}

\begin{probbox}{25}{\relax}
In the ``Coin Game'', you toss three coins at the same time. You win only if the $3$ coins all show heads or all show tails. If you play the game once only, what is the probability of winning?
\choices{$\dfrac16$}{$\dfrac14$}{$\dfrac{2}{27}$}{$\dfrac23$}{$\dfrac13$}
\end{probbox}
\probsol{
\textbf{Understanding the problem.} Each coin has $2$ faces: heads (H) or tails (T), equally likely. You win with HHH or TTT. We list all possible outcomes and count.

\begin{recallbox}When all outcomes are equally likely: the \emph{probability} of an event $=$ (number of favourable outcomes) $\div$ (total number of possible outcomes).\end{recallbox}

\step{1}{List all outcomes with a tree diagram}
\begin{figC}\begin{scope}[level distance=1.9cm,grow'=right,
  level 1/.style={sibling distance=2.4cm},level 2/.style={sibling distance=1.2cm},level 3/.style={sibling distance=0.6cm},
  every node/.style={font=\footnotesize\bfseries}, edge from parent/.style={draw=tmBlue!70,thick}]
\node[text=Ink] {Toss}
  child {node[text=tmOrange] {H}
    child {node[text=tmOrange] {H} child {node[text=AnsRed] {H \ $\to$ HHH (win)}} child {node[text=Slate] {T \ $\to$ HHT}}}
    child {node[text=tmOrange] {T} child {node[text=Slate] {H \ $\to$ HTH}} child {node[text=Slate] {T \ $\to$ HTT}}}}
  child {node[text=tmPurple] {T}
    child {node[text=tmPurple] {H} child {node[text=Slate] {H \ $\to$ THH}} child {node[text=Slate] {T \ $\to$ THT}}}
    child {node[text=tmPurple] {T} child {node[text=Slate] {H \ $\to$ TTH}} child {node[text=AnsRed] {T \ $\to$ TTT (win)}}}};
\node[font=\scriptsize,text=Slate] at (1.9,2.55) {coin $1$};
\node[font=\scriptsize,text=Slate] at (3.8,2.55) {coin $2$};
\node[font=\scriptsize,text=Slate] at (5.7,2.55) {coin $3$};
\end{scope}\end{figC}
The first coin has $2$ possibilities, for each of them the second coin has $2$, and then the third coin has $2$: altogether $2\cdot2\cdot2=8$ equally likely outcomes.

\step{2}{Count the favourable outcomes}
Only $2$ outcomes win: HHH and TTT.

\step{3}{Compute the probability}
\[P(\text{win})=\frac{2}{8}=\frac14.\]

\begin{warnbox}A common \emph{wrong} argument: ``The number of heads can be $0,1,2,3$ -- four cases, and we win in $2$ of them ($0$ and $3$), so the probability is $\dfrac24$.'' This is wrong because these four cases are \emph{not} equally likely: ``$1$ head'' happens in $3$ ways (HTT, THT, TTH), while ``$3$ heads'' happens in only $1$ way.\end{warnbox}

\begin{tipbox}A quick way: the first coin can show anything. To win, the second coin must match the first (probability $\dfrac12$) and the third must match too (probability $\dfrac12$): $\dfrac12\cdot\dfrac12=\dfrac14$.\end{tipbox}

\verify{The $8$ outcomes in the tree are all different and cover everything; exactly $2$ of them win. $\dfrac28=\dfrac14$.}
Choose option \textbf{(B)}.
\ans{(B) $\dfrac14$}
}

\begin{probbox}{26}{\relax}
In the right-angled triangle $PQR$, the right angle is at $Q$ and $PQ=QR$. The segments $QS$, $TU$ and $VW$ are perpendicular to $PR$, and the segments $ST$ and $UV$ are perpendicular to $QR$, as shown. What fraction of the area of triangle $PQR$ is shaded?
\begin{figC}\begin{scope}[scale=0.42]
\coordinate (Q) at (0,0); \coordinate (R) at (8,0); \coordinate (P) at (0,8);
\coordinate (S) at (4,4); \coordinate (T) at (4,0); \coordinate (U) at (6,2); \coordinate (V) at (6,0); \coordinate (W) at (7,1);
\fill[tmBlue!35] (S)--(T)--(U)--cycle; \fill[tmBlue!35] (U)--(V)--(W)--cycle;
\draw[Navy,line width=1pt] (P)--(Q)--(R)--cycle;
\draw[Navy,line width=0.8pt] (Q)--(S)--(T)--(U)--(V)--(W);
\draw[Navy,line width=0.5pt] (0,0.5)--(0.5,0.5)--(0.5,0);
\node[above] at (P) {$P$}; \node[below left] at (Q) {$Q$}; \node[below right] at (R) {$R$};
\node[above right] at (S) {$S$}; \node[below] at (T) {$T$}; \node[above right] at (U) {$U$}; \node[below] at (V) {$V$}; \node[above right] at (W) {$W$};
\end{scope}\end{figC}
\choices{$\dfrac{3}{16}$}{$\dfrac38$}{$\dfrac{5}{16}$}{$\dfrac{5}{32}$}{$\dfrac{7}{32}$}
\end{probbox}
\probsol{
\textbf{Understanding the problem.} Triangle $PQR$ is a \emph{right isosceles} triangle with the right angle at $Q$. The picture is built by repeatedly ``dropping a perpendicular''. The shaded part consists of two triangles: $STU$ and $UVW$. Key idea: \emph{dropping a perpendicular from the right-angle vertex of a right isosceles triangle cuts it into two equal halves, and each half is again a right isosceles triangle.}

\begin{recallbox}In an isosceles triangle, the altitude from the apex (to the base) is also the median, and it divides the triangle into two congruent triangles. A right isosceles triangle has two acute angles of $45^\circ$.\end{recallbox}

\begin{proofbox}Take triangle $XYZ$, right isosceles at $X$ ($XY=XZ$), and draw $XH\perp YZ$. The right triangles $XHY$ and $XHZ$ have equal hypotenuses $XY=XZ$ and the common leg $XH$, so they are congruent (hypotenuse--leg). So each has half the area of $XYZ$. Moreover $\widehat{Y}=45^\circ$ and $\widehat{XHY}=90^\circ$, so $\widehat{HXY}=45^\circ$: triangle $XHY$ is right isosceles at $H$. The same holds for $XHZ$.\end{proofbox}

\step{1}{Keep halving}
Let the area of triangle $PQR$ be $1$ (the whole triangle).
\begin{itemize}[leftmargin=1.4em,itemsep=1pt]
\item $QS\perp PR$ cuts $PQR$ into two halves: triangle $QSR$ has area $\dfrac12$ and is right isosceles at $S$.
\item $ST\perp QR$ cuts $QSR$ into two halves: triangle $STR$ has area $\dfrac14$ and is right isosceles at $T$.
\item $TU\perp SR$ cuts $STR$ into two halves: triangle $STU$ \textbf{(shaded)} has area $\dfrac18$; triangle $TUR$ has area $\dfrac18$ and is right isosceles at $U$.
\item $UV\perp TR$ cuts $TUR$ into two halves: triangle $UVR$ has area $\dfrac1{16}$ and is right isosceles at $V$.
\item $VW\perp UR$ cuts $UVR$ into two halves: triangle $UVW$ \textbf{(shaded)} has area $\dfrac1{32}$.
\end{itemize}

\step{2}{Add the shaded areas}
\[\frac18+\frac1{32}=\frac{4}{32}+\frac{1}{32}=\frac{5}{32}.\]

\step{3}{Check with a square grid}
Put $Q$ at the origin with legs of $8$ units: $Q(0,0)$, $R(8,0)$, $P(0,8)$. Then $S(4,4)$, $T(4,0)$, $U(6,2)$, $V(6,0)$, $W(7,1)$ (as in the figure). Triangle $PQR$ has area $\dfrac{8\cdot 8}{2}=32$ squares. Triangle $STU$ has side $ST=4$ (vertical) and $U$ is $2$ units away from $ST$, so its area is $\dfrac{4\cdot 2}{2}=4$ squares. Triangle $UVW$ has $UV=2$ and $W$ is $1$ unit from $UV$, area $\dfrac{2\cdot1}{2}=1$ square. Shaded: $4+1=5$ squares out of $32$, i.e.\ $\dfrac5{32}$.

\begin{recallbox}Area of a triangle $=\dfrac{\text{base}\times\text{height}}{2}$.\end{recallbox}

\begin{warnbox}Many students pick $\dfrac{3}{16}$ or $\dfrac38$ by ``estimating from the picture''. Always compute with the chain of halvings: each smaller triangle is half of the triangle containing it.\end{warnbox}

\verify{Both methods (halving and counting squares) give $\dfrac{5}{32}$.}
Choose option \textbf{(D)}.
\ans{(D) $\dfrac{5}{32}$}
}

\begin{probbox}{27}{\relax}
Mai went to a second-hand market where, every $15$ minutes, the price of every item was reduced by $10\%$ of its \emph{current} price. At $9$:$00$ a.m., the price of a rug was $100\,000$ dong. At $9$:$15$ a.m., the price was reduced to $90\,000$ dong. As soon as the price of the rug fell \emph{below} $80\,000$ dong, Mai bought it. At what time did Mai buy the rug?
\choicesII{$9$:$45$ a.m.}{$9$:$15$ a.m.}{$9$:$30$ a.m.}{$10$:$15$ a.m.}{$10$:$00$ a.m.}
\end{probbox}
\probsol{
\textbf{Understanding the problem.} Each reduction is $10\%$ of the \emph{current price} (not of the original $100\,000$ dong). We compute the price after each $15$-minute mark until it is less than $80\,000$ dong.

\begin{recallbox}Reducing a number by $10\%$ means taking away $\dfrac{10}{100}=\dfrac1{10}$ of it; what remains is $90\%$ of it. So: new price $=$ old price $\times 0.9$ (or old price $-$ old price $\div 10$).\end{recallbox}

\step{1}{Make a table of prices over time}
\begin{center}
\begin{tabular}{|c|c|c|c|}
\hline
\rowcolor{Tint}\textbf{Time} & \textbf{Price before} & \textbf{Reduction ($10\%$)} & \textbf{New price}\\
\hline
$9$:$00$ a.m. & -- & -- & $100\,000$\\
$9$:$15$ a.m. & $100\,000$ & $10\,000$ & $90\,000$\\
$9$:$30$ a.m. & $90\,000$ & $9\,000$ & $81\,000$\\
$9$:$45$ a.m. & $81\,000$ & $8\,100$ & $\mathbf{72\,900}$\\
\hline
\end{tabular}
\end{center}

\step{2}{Find the first time the price is below $80\,000$ dong}
At $9$:$30$ a.m. the price is $81\,000$ dong -- still \emph{more} than $80\,000$ dong, so Mai does not buy yet. At $9$:$45$ a.m. the price is $72\,900$ dong $<80\,000$ dong: this is the first time it falls below $80\,000$ dong. Mai buys it at $9$:$45$ a.m.

\begin{warnbox}(1) The reduction is $10\%$ \emph{of the current price}, so the amount taken off gets smaller each time ($10\,000$, then $9\,000$, then $8\,100$\dots); it is not $10\,000$ dong every time. (2) At $9$:$30$ a.m. the price $81\,000$ dong is ``very close'' to $80\,000$ but not yet \emph{below} $80\,000$, so it is not the moment to buy.\end{warnbox}

\verify{$100\,000\times 0.9\times 0.9\times 0.9=72\,900$ (after $3$ reductions, i.e.\ $3\times 15=45$ minutes after $9$:$00$); after $2$ reductions: $100\,000\times0.81=81\,000>80\,000$.}
Choose option \textbf{(A)}.
\ans{(A) $9$:$45$ a.m.}
}

\begin{probbox}{28}{\relax}
The number $2018$ is used to create six-digit positive integers: each of them must contain the digits $2,0,1,8$ \emph{together and in this order}. For example, $720\,186$ is allowed, but $209\,318$ and $210\,893$ are not. How many of these six-digit integers are divisible by $9$?
\choices{$28$}{$27$}{$31$}{$34$}{$22$}
\end{probbox}
\probsol{
\textbf{Understanding the problem.} The six-digit number contains the block ``$2018$'' plus two more digits (call them $a$ and $b$). We must count those divisible by $9$.

\begin{recallbox}\textbf{Divisibility by $9$:} a number is divisible by $9$ exactly when the sum of its digits is divisible by $9$. The first digit of a six-digit number cannot be $0$.\end{recallbox}

\step{1}{Condition on the digit sum}
The digits of the block add up to $2+0+1+8=11$. The whole number has digit sum $11+a+b$. Since $a,b$ are digits, $0\le a+b\le 18$, so $11\le 11+a+b\le 29$. The multiples of $9$ in this range are $18$ and $27$. Therefore
\[a+b=7\qquad\text{or}\qquad a+b=16.\]
The only pairs with sum $16$ are $(7,9),\ (8,8),\ (9,7)$.

\step{2}{Three possible positions of the block $2018$}
\begin{center}
\begin{tabular}{|c|c|c|c|c|}
\hline
\rowcolor{Tint}\textbf{Form} & \textbf{Extra condition} & \textbf{$a+b=7$} & \textbf{$a+b=16$} & \textbf{Total}\\
\hline
$\nd{2018ab}$ & none ($a$ may be $0$) & $8$ pairs: $a=0,1,\dots,7$ & $3$ & $11$\\
$\nd{a2018b}$ & $a\neq 0$ & $7$ pairs: $a=1,\dots,7$ & $3$ & $10$\\
$\nd{ab2018}$ & $a\neq 0$ & $7$ pairs: $a=1,\dots,7$ & $3$ & $10$\\
\hline
\multicolumn{4}{|r|}{\textbf{Total}} & $\mathbf{31}$\\
\hline
\end{tabular}
\end{center}
About the column ``$a+b=7$'': the pairs are $(0,7),(1,6),(2,5),(3,4),(4,3),(5,2),(6,1),(7,0)$ -- that is $8$ pairs. When $a$ is the first digit of the number, we drop $(0,7)$, leaving $7$ pairs. The pairs with sum $16$ all have $a\ge 7$, so none is dropped.

\step{3}{Is any number counted twice?}
A number would be counted twice only if the block $2018$ appeared in two different places in the same six-digit number. Two $4$-digit blocks inside $6$ digits must overlap, and the block $2018$ cannot overlap itself (for example, shifting it by $1$ place would force the digit $0$ to coincide with the digit $2$\dots). So no number is counted twice.

\begin{warnbox}It is easy to forget one of the two cases: many students only consider $a+b=7$ (getting $8+7+7=22$, which is exactly the distractor (E)) and forget $a+b=16$. It is also easy to forget that the first digit cannot be $0$ (that would give $34$).\end{warnbox}

\verify{A few examples: $201\,870$ has digit sum $18$; $720\,186$ has digit sum $24$ -- not divisible by $9$ (correct, since $7+6=13$ is neither $7$ nor $16$); $920\,187$ has digit sum $27$, divisible by $9$. A computer search of all numbers from $100\,000$ to $999\,999$ also gives exactly $31$.}
Choose option \textbf{(C)}.
\ans{(C) $31$}
}

\begin{probbox}{29}{\relax}
The number $N$ is the product of all positive odd integers from $1$ to $99$ that do \emph{not} end in the digit $5$. That is,
\[N=1\times 3\times 7\times 9\times 11\times 13\times 17\times 19\times\cdots\times 91\times 93\times 97\times 99.\]
The units digit of $N$ is
\choices{$1$}{$3$}{$5$}{$7$}{$9$}
\end{probbox}
\probsol{
\textbf{Understanding the problem.} $N$ is a gigantic number that we cannot compute directly. But we only need its \emph{last digit}, and the last digit of a product depends only on the last digits of the factors.

\begin{recallbox}The units digit of a product equals the units digit of the product of the units digits of the factors. Example: $13\times 17=221$ ends in $1$, just like $3\times 7=21$.\end{recallbox}

\step{1}{Which units digits do the factors have?}
Odd numbers that do not end in $5$ can only end in $1,3,7,9$. Split the factors into $10$ groups, one for each ten:
\[(1\cdot3\cdot7\cdot9),\ (11\cdot13\cdot17\cdot19),\ (21\cdot23\cdot27\cdot29),\ \dots,\ (91\cdot93\cdot97\cdot99).\]
Each group has $4$ factors ending in $1,3,7,9$; there are $10$ groups (the tens $0$ to $9$), $40$ factors in all.

\step{2}{Units digit of each group}
$1\cdot3\cdot7\cdot9=189$, which ends in $9$. So the product of each group ends in $9$.

\step{3}{Multiply the $10$ groups}
The units digit of $N$ equals the units digit of $\underbrace{9\cdot 9\cdots 9}_{10\text{ factors}}=9^{10}$. Now $9\cdot 9=81$ ends in $1$. Pair up the $10$ nines into $5$ pairs, each ending in $1$, so $9^{10}$ ends in $1\cdot1\cdot1\cdot1\cdot1=1$.

\begin{tipbox}Powers of $9$ have alternating last digits: $9^1\to 9$, $9^2\to 1$, $9^3\to 9$, $9^4\to 1$,\dots: odd exponents end in $9$, even exponents end in $1$.\end{tipbox}

\begin{warnbox}Do not think ``a product of odd numbers ends in $9$'', and do not choose $5$ by reflex: the numbers ending in $5$ have been left out. Also count the groups correctly: from $1$ to $99$ there are $10$ groups (the first one is $1,3,7,9$).\end{warnbox}

\verify{Step by step: after group $1$ the last digit is $9$, after group $2$ it is $9\cdot9\to1$, after group $3$ it is $9$,\dots, after group $10$ (an even number) it is $1$. A computer computing $N$ exactly also gives last digit $1$.}
Choose option \textbf{(A)}.
\ans{(A) $1$}
}

\begin{probbox}{30}{\relax}
Five students wrote a quiz with a maximum score of $50$. The scores of four of the students were $42$, $43$, $46$ and $49$. The score of the fifth student was $N$. The mean of the five students' scores was the same as the median of the five students' scores. The number of possible values of $N$ is
\choices{$3$}{$4$}{$1$}{$0$}{$2$}
\end{probbox}
\probsol{
\textbf{Understanding the problem.} $N$ is a score from $0$ to $50$. The mean depends on $N$; the median also depends on where $N$ falls when the five scores are put in order. We consider each case for the position of $N$.

\begin{recallbox}The \textbf{median} of five numbers is the \emph{middle} one (the third) when they are arranged from smallest to largest. The \textbf{mean} $=$ sum $\div$ how many numbers.\end{recallbox}

\step{1}{Write the mean in terms of $N$}
The four known scores add to $42+43+46+49=180$. The mean of the five scores is $\dfrac{180+N}{5}$.

\step{2}{What is the median? -- Three cases}
Arrange $42,43,46,49$ together with $N$ in increasing order:
\begin{itemize}[leftmargin=1.4em,itemsep=1pt]
\item If $N\le 43$: the order is $\dots,43,46,49$ with $N$ in one of the first two places; the third number is $43$. Median $=43$.
\item If $43\le N\le 46$: the order is $42,43,N,46,49$. Median $=N$.
\item If $N\ge 46$: the order is $42,43,46,\dots$ with $N$ in one of the last two places. Median $=46$.
\end{itemize}
\begin{figC}
\draw[->,thick,Navy] (0,0)--(10.6,0);
\foreach \x/\lab in {0/0,6/42,6.6/43,8.4/46,10.2/50}{\draw[Navy] (\x,0.1)--(\x,-0.1); \node[below,font=\footnotesize] at (\x,-0.1) {$\lab$};}
\draw[tmBlue,line width=3pt,opacity=0.5] (0,0.35)--(6.6,0.35); \node[above,font=\scriptsize,text=tmBlue] at (3.3,0.4) {median $=43$};
\draw[tmOrange,line width=3pt,opacity=0.6] (6.6,0.35)--(8.4,0.35); \node[above,font=\scriptsize,text=tmOrange!80!black] at (7.5,0.4) {median $=N$};
\draw[tmGreen,line width=3pt,opacity=0.5] (8.4,0.35)--(10.2,0.35); \node[above,font=\scriptsize,text=tmGreen!70!black] at (9.4,0.85) {median $=46$};
\end{figC}

\step{3}{Solve each case}
\begin{itemize}[leftmargin=1.4em,itemsep=1pt]
\item \textbf{Case 1} (median $43$): $\dfrac{180+N}{5}=43\Rightarrow 180+N=215\Rightarrow N=35$. Check: $35\le 43$ -- it is in the case being considered. \textbf{Accept.}
\item \textbf{Case 2} (median $N$): $\dfrac{180+N}{5}=N\Rightarrow 180+N=5N\Rightarrow 4N=180\Rightarrow N=45$. Check: $43\le 45\le 46$. \textbf{Accept.}
\item \textbf{Case 3} (median $46$): $\dfrac{180+N}{5}=46\Rightarrow 180+N=230\Rightarrow N=50$. Check: $50\ge 46$, and $50$ does not exceed the maximum score. \textbf{Accept.}
\end{itemize}

\begin{warnbox}After solving for $N$ in each case, you \emph{must check} that $N$ really belongs to that case (and lies between $0$ and $50$). In this problem all three values are valid, but in many other problems some values get rejected. Do not stop after finding one value!\end{warnbox}

\verify{
$N=35$: scores $35,42,43,46,49$, sum $215$, mean $43$, median $43$.
$N=45$: scores $42,43,45,46,49$, sum $225$, mean $45$, median $45$.
$N=50$: scores $42,43,46,49,50$, sum $230$, mean $46$, median $46$.}
There are $3$ values of $N$, which is option \textbf{(A)}.
\ans{(A) $3$}
}

% ---------------- part c ----------------
% =====================================================================
% Topic 4 — Number Sense — part c: problems 31–40 (extending problems 1–15)
% =====================================================================
\clearpage
\dividerbar{Extension and challenge}

% ---------------------------------------------------------------- Problem 31
\begin{probbox}{31}{\relax}
Consider the numbers: $-1.5;\quad 0.25;\quad -0.75;\quad 1.2;\quad -2$.
\begin{enumerate}[label=\alph*),leftmargin=1.6em,itemsep=0pt,topsep=2pt]
  \item Arrange the numbers in increasing order.
  \item Which numbers lie between $-1$ and $1$ on the number line?
  \item Find the sum of the largest and the smallest number.
\end{enumerate}
\end{probbox}
\probsol{
\textbf{Understanding the problem.} This problem combines two skills from Problems 1 and 3: comparing decimals and comparing
negative numbers.
\begin{recallbox}Negative numbers are always less than positive numbers. Of two negative numbers, the one farther from $0$ is
smaller. For two positive numbers, compare the whole parts, then the decimal places one by one.\end{recallbox}
\step{1}{Split into two groups: negative and positive}
Negative: $-1.5;\ -0.75;\ -2$. \quad Positive: $0.25;\ 1.2$.
\step{2}{Order the negative numbers}
Distance from $0$: $-2$ is $2$ away; \ $-1.5$ is $1.5$ away; \ $-0.75$ is $0.75$ away. The one farthest from $0$ is the smallest:
\[ -2 < -1.5 < -0.75. \]
\step{3}{Order the positive numbers, then put everything together}
$0.25 < 1.2$ (whole parts $0 < 1$). Negative numbers first, positive numbers after:
\[ -2 < -1.5 < -0.75 < 0.25 < 1.2. \]
\begin{figC}\begin{scope}[x=1.6cm]
  \draw[->,thick,Navy] (-2.5,0)--(1.7,0);
  \foreach \x in {-2,-1,0,1}{\draw[Navy] (\x,0.08)--(\x,-0.08); \node[below,font=\scriptsize] at (\x,-0.1) {$\x$};}
  \foreach \x/\l in {-2/{-2},-1.5/{-1.5},-0.75/{-0.75},0.25/{0.25},1.2/{1.2}}{\fill[tmOrange] (\x,0) circle (2.6pt); \node[above,font=\scriptsize] at (\x,0.1) {$\l$};}
  \draw[tmGreen!70!black,line width=2pt,opacity=0.35] (-1,0)--(1,0);
\end{scope}\end{figC}
\step{4}{The numbers between $-1$ and $1$}
On the number line, the green segment from $-1$ to $1$ contains $-0.75$ and $0.25$. ($-1.5$ is to the left of $-1$; $1.2$ is
to the right of $1$.)
\step{5}{Sum of the largest and the smallest}
The largest is $1.2$, the smallest is $-2$:
\[ 1.2 + (-2) = -(2 - 1.2) = -0.8. \]
\begin{recallbox}Adding two numbers with opposite signs: subtract the smaller size from the larger (ignoring signs), then use
the sign of the number farther from $0$.\end{recallbox}
\begin{warnbox}Many students put $-0.75$ before $-1.5$ because ``$0.75 < 1.5$''. For negative numbers it is reversed: $-1.5 < -0.75$.\end{warnbox}
\verify{$-0.8 - 1.2 = -2$. Correct.}
\ans{a) $-2 < -1.5 < -0.75 < 0.25 < 1.2$;\quad b) $-0.75$ and $0.25$;\quad c) $-0.8$}
}

% ---------------------------------------------------------------- Problem 32
\begin{probbox}{32}{\relax}
Consider the expression $2 + 3 \times 4 - 1$.
\begin{enumerate}[label=\alph*),leftmargin=1.6em,itemsep=0pt,topsep=2pt]
  \item Evaluate the expression.
  \item Insert \emph{one} pair of brackets $(\ )$ into the expression (without changing the order of the numbers or the signs)
  to get the largest possible value. What is that largest value?
  \item Same question as b), but for the smallest value.
\end{enumerate}
\end{probbox}
\probsol{
\textbf{Understanding the problem.} This extends Problems 4 and 7: the position of the brackets changes the order of the
operations, and so changes the result.
\begin{recallbox}Brackets first; then multiplication and division; finally addition and subtraction (from left to right).\end{recallbox}
\step{1}{Value without brackets}
$2 + 3\times 4 - 1 = 2 + 12 - 1 = 13$.
\step{2}{List systematically every way to place one pair of brackets}
A pair of brackets must enclose at least two neighbouring numbers. Call the four numbers, in order, number $1, 2, 3, 4$
($2;\ 3;\ 4;\ 1$). The brackets can start at number $i$ and end at number $j$ with $i < j$: $6$ ways.
\begin{center}
\begin{tabular}{|c|l|c|}
\hline
\rowcolor{tmBlue!10} Brackets around & Expression & Value\\ \hline
numbers 1 -- 2 & $(2+3)\times 4 - 1 = 5\times 4 - 1$ & $\mathbf{19}$\\ \hline
numbers 1 -- 3 & $(2+3\times 4) - 1 = 14 - 1$ & $13$\\ \hline
numbers 1 -- 4 & $(2+3\times 4-1)$ & $13$\\ \hline
numbers 2 -- 3 & $2 + (3\times 4) - 1$ & $13$\\ \hline
numbers 2 -- 4 & $2 + (3\times 4 - 1) = 2 + 11$ & $13$\\ \hline
numbers 3 -- 4 & $2 + 3\times(4 - 1) = 2 + 9$ & $\mathbf{11}$\\ \hline
\end{tabular}
\end{center}
\step{3}{Read off the results}
The largest is $19$, with $(2+3)\times 4 - 1$: the sum $2 + 3$ is ``blown up'' $4$ times.\\
The smallest is $11$, with $2 + 3\times(4-1)$: $3$ is multiplied by a smaller number ($3$ instead of $4$).
\begin{tipbox}For a large result, let the multiplication act on a large sum; for a small result, let it act on a difference.
But to be sure, the best way is still to list every case, as in the table.\end{tipbox}
\verify{$(2+3)\times 4 - 1 = 20 - 1 = 19$ and $2 + 3\times 3 = 11$.}
\ans{a) $13$;\quad b) $(2+3)\times 4 - 1 = 19$;\quad c) $2 + 3\times(4-1) = 11$}
}

% ---------------------------------------------------------------- Problem 33
\begin{probbox}{33}{\relax}
An's storybook has $360$ pages. On the first day An reads $\dfrac{1}{4}$ of the book. On the second day An reads
$\dfrac{1}{3}$ \emph{of the remaining pages}. After two days, how many pages has An not read yet, and what fraction of the
book is that?
\end{probbox}
\probsol{
\textbf{Understanding the problem.} Careful: the fraction $\dfrac13$ on the second day is taken of the \emph{remaining
pages}, not of the whole book. This is a ``two-step'' version of Problem 9.
\begin{recallbox}$\dfrac{m}{n}$ of $a$ equals $a \div n \times m$.\end{recallbox}
\step{1}{First day}
An reads $360 \div 4 = 90$ (pages). Remaining: $360 - 90 = 270$ (pages).
\step{2}{Second day (taken of $270$ pages)}
An reads $270 \div 3 = 90$ (pages). Remaining: $270 - 90 = 180$ (pages).
\step{3}{Fraction of the book left}
$\dfrac{180}{360} = \dfrac12$.
\begin{barfigC}
  \blk{0}{2}{0}{0.6}{tmBlue}{$90$}
  \blk{2}{4}{0}{0.6}{tmGreen}{$90$}
  \blk{4}{6}{0}{0.6}{tmOrange}{$90$}
  \blk{6}{8}{0}{0.6}{tmOrange}{$90$}
  \brcup{0}{2}{0.5}{day 1}
  \brcup{2}{4}{0.5}{day 2}
  \brcup{4}{8}{0.5}{not read}
  \brcdn{2}{8}{-0.5}{left after day 1: $270$ pages, split into $3$ parts}
\end{barfigC}
\begin{warnbox}A common mistake: adding $\dfrac14 + \dfrac13 = \dfrac{7}{12}$ and saying $\dfrac{5}{12}$ of the book
($150$ pages) is left. Wrong, because $\dfrac13$ is of the \emph{remaining part} ($270$ pages), not of the whole book ($360$ pages).\end{warnbox}
\step{4}{Another way: multiply the ``remaining'' fractions}
After day 1, $\dfrac34$ of the book remains; after day 2, $\dfrac23$ of that remains, i.e.\ $\dfrac23\times\dfrac34 = \dfrac{6}{12} = \dfrac12$
of the book, which is $360 \div 2 = 180$ pages.
\verify{$90 + 90 + 180 = 360$. Correct.}
\ans{$180$ pages, which is $\dfrac12$ of the book}
}

% ---------------------------------------------------------------- Problem 34
\begin{probbox}{34}{\relax}
The first test has $40$ questions, and Bình answers $32$ of them correctly.
\begin{enumerate}[label=\alph*),leftmargin=1.6em,itemsep=0pt,topsep=2pt]
  \item What percentage of the questions in the first test did Bình answer correctly?
  \item The second test has $60$ questions. At least how many questions must Bình answer correctly in the second test so
  that, over both tests together, his correct answers make up at least $85\%$ of all the questions?
\end{enumerate}
\end{probbox}
\probsol{
\textbf{Understanding the problem.} Part a) asks for a percentage (the reverse of Problem 5). Part b) asks for a percentage of a total.
\begin{recallbox}$a$ as a percentage of $b$ is $\dfrac{a}{b}\times 100\%$. \ $p\%$ of $a$ is $a\times\dfrac{p}{100}$.\end{recallbox}
\step{1}{Percentage in the first test}
$\dfrac{32}{40}\times 100\% = 0.8\times 100\% = 80\%$.
\step{2}{Total number of questions in both tests}
$40 + 60 = 100$ (questions).
\step{3}{Correct answers needed overall}
$85\%$ of $100$ questions is $100\times\dfrac{85}{100} = 85$ (questions).
\step{4}{Correct answers needed in the second test}
Bình already has $32$ correct answers, so he needs at least $85 - 32 = 53$ (questions) in the second test.
\begin{warnbox}Do not take ``the average of the two percentages'': solving $(80\% + x) \div 2 = 85\%$ gives $x = 90\%$ of $60$
questions, i.e.\ $54$ questions -- wrong! The tests have different numbers of questions, so we must work with the
\emph{total number of questions}.\end{warnbox}
\verify{With $53$ correct: $(32 + 53) \div 100 = 85\%$ (just enough). With $52$: $84 \div 100 = 84\% < 85\%$ (not enough).}
\ans{a) $80\%$;\quad b) at least $53$ questions}
}

% ---------------------------------------------------------------- Problem 35
\begin{probbox}{35}{\relax}
A shop sells cashew nuts in packs: a $3$ kg pack costs $255$ thousand dong, a $5$ kg pack costs $410$ thousand dong (no loose sales).
\begin{enumerate}[label=\alph*),leftmargin=1.6em,itemsep=0pt,topsep=2pt]
  \item Which pack is cheaper per kilogram?
  \item Ms Hằng needs to buy \emph{exactly} $15$ kg. How should she buy them to spend as little as possible? How much does it cost?
  \item How much does it cost to buy \emph{exactly} $13$ kg?
\end{enumerate}
\end{probbox}
\probsol{
\textbf{Understanding the problem.} This extends Problem 6: compare unit prices, then find how to combine packs to get the exact mass.
\begin{recallbox}Unit price $=$ total cost $\div$ quantity.\end{recallbox}
\step{1}{Unit price of each pack}
$3$ kg pack: $255 \div 3 = 85$ (thousand dong/kg). \quad $5$ kg pack: $410 \div 5 = 82$ (thousand dong/kg).\\
So the $5$ kg pack is cheaper ($82 < 85$).
\step{2}{Ways to make exactly $15$ kg}
Let $a$ be the number of $3$ kg packs and $b$ the number of $5$ kg packs: $3a + 5b = 15$. Try $b = 0, 1, 2, 3$:
\begin{center}
\begin{tabular}{|c|c|c|c|}
\hline
\rowcolor{tmBlue!10} $b$ & $15 - 5b$ & $a$ & Cost (thousand dong)\\ \hline
$0$ & $15$ & $5$ & $5\times 255 = 1275$\\ \hline
$1$ & $10$ & not divisible by $3$ & --\\ \hline
$2$ & $5$ & not divisible by $3$ & --\\ \hline
$3$ & $0$ & $0$ & $3\times 410 = \mathbf{1230}$\\ \hline
\end{tabular}
\end{center}
Buying $3$ packs of $5$ kg is cheapest: $1\,230$ thousand dong $= 1\,230\,000$ dong.
\step{3}{Exactly $13$ kg}
$3a + 5b = 13$: \ $b = 0 \Rightarrow 13$ (not divisible by $3$); \ $b = 1 \Rightarrow 8$ (not divisible by $3$);
\ $b = 2 \Rightarrow 3$, $a = 1$. There is only one way: $1$ pack of $3$ kg and $2$ packs of $5$ kg.\\
Cost: $255 + 2\times 410 = 255 + 820 = 1075$ (thousand dong) $= 1\,075\,000$ dong.
\begin{warnbox}Part c) cannot be answered by computing $13\times 82$, because the shop does not sell loose nuts: the packs
must add up to exactly $13$ kg.\end{warnbox}
\verify{$3 + 5 + 5 = 13$ kg; $5 + 5 + 5 = 15$ kg.}
\ans{a) the $5$ kg pack ($82$ thousand dong/kg);\quad b) $3$ packs of $5$ kg, $1\,230\,000$ dong;\quad c) $1\,075\,000$ dong}
}

% ---------------------------------------------------------------- Problem 36
\begin{probbox}{36}{\relax}
The five members of a group weigh (in kg) $38;\ 41;\ 40;\ 37$ and $x$. Their mean weight is $40$ kg.
\begin{enumerate}[label=\alph*),leftmargin=1.6em,itemsep=0pt,topsep=2pt]
  \item Find $x$. Is the fifth member heavier or lighter than the mean, and by how many kilograms?
  \item Find the median weight of the group.
  \item A sixth member joins the group, and the mean weight of all six becomes $41$ kg. How much does the sixth member weigh?
\end{enumerate}
\end{probbox}
\probsol{
\textbf{Understanding the problem.} This extends Problem 14: use the sum and the deviations to find the missing number; add
the median and an ``extra member'' question.
\begin{recallbox}Sum $=$ mean $\times$ how many numbers. The deviations from the mean add up to $0$.\end{recallbox}
\step{1}{Find $x$ using deviations}
Compared with $40$: $38 \to -2$; \ $41 \to +1$; \ $40 \to 0$; \ $37 \to -3$. Sum: $-2 + 1 + 0 - 3 = -4$.\\
For the deviations to add up to $0$, the fifth member must be $+4$: $x = 40 + 4 = 44$ (kg), $4$ kg heavier than the mean.\\
Check with the sum: $5\times 40 = 200$; \ $200 - (38 + 41 + 40 + 37) = 200 - 156 = 44$.
\step{2}{Median}
In increasing order: $37;\ 38;\ \mathbf{40};\ 41;\ 44$. The middle (third) number is $40$. The median is $40$ kg.
\step{3}{The sixth member}
Total weight of six members: $6\times 41 = 246$ (kg). The first five weigh $200$ kg in total. The sixth weighs $246 - 200 = 46$ (kg).
\begin{tipbox}Think in deviations: the mean rises by $1$ kg for all $6$ people, so $6$ kg extra above the level of $40$ are needed;
the new member must weigh $40 + 6 = 46$ kg.\end{tipbox}
\begin{warnbox}It is not ``the sixth member weighs $41$ kg'' -- then the mean would rise only a little: $241 \div 6 \approx 40.17$ kg.\end{warnbox}
\verify{$(37 + 38 + 40 + 41 + 44 + 46) \div 6 = 246 \div 6 = 41$. Correct.}
\ans{a) $x = 44$ kg, $4$ kg heavier than the mean;\quad b) $40$ kg;\quad c) $46$ kg}
}

% ---------------------------------------------------------------- Problem 37
\begin{probbox}{37}{\stars{1}}
A running track around the schoolyard is $200$ m long. Mai runs $6$ laps in $4$ minutes. Tú runs at a constant $18$ km/h.
\begin{enumerate}[label=\alph*),leftmargin=1.6em,itemsep=0pt,topsep=2pt]
  \item Who is faster?
  \item Keeping the same speeds, how many minutes does each of them need to run $3$ km?
  \item They start at the same time from the same point of the track, but run in \emph{opposite directions}. After how many
  seconds do they meet for the first time?
  \item One day Tú increases his speed to $21.6$ km/h while Mai runs as before. They start at the same time from the same
  point, running in the \emph{same direction}. How long is it until Tú is exactly one lap ahead of Mai?
\end{enumerate}
\end{probbox}
\probsol{
\textbf{Understanding the problem.} This extends Problem 12: find a speed from laps run, convert km/h to m/s, then study two
runners on the same circular track.
\begin{recallbox}$v = s \div t$, \ $t = s \div v$. \ $1$ km/h $= \dfrac{5}{18}$ m/s, so to convert km/h to m/s multiply by $5$ and divide by $18$.\end{recallbox}
\step{1}{Mai's speed}
Distance: $6\times 200 = 1200$ (m). Time: $4$ minutes $= 240$ seconds. \ $v_{\text{Mai}} = 1200 \div 240 = 5$ (m/s).
\step{2}{Tú's speed (in m/s)}
$v_{\text{Tú}} = 18\times\dfrac{5}{18} = 5$ (m/s). They run \textbf{equally fast}!
\begin{warnbox}Comparing ``$5$'' with ``$18$'' directly is meaningless, because the units differ (m/s and km/h). Convert to the
same unit first.\end{warnbox}
\step{3}{Time to run $3$ km}
$3$ km $= 3000$ m. \ $t = 3000 \div 5 = 600$ (seconds) $= 600 \div 60 = 10$ (minutes). Both need $10$ minutes.
\step{4}{Opposite directions}
Every second Mai runs $5$ m and Tú runs $5$ m the other way, so the ``combined'' distance they cover grows by
$5 + 5 = 10$ m per second. They first meet when the \emph{total} of their distances is exactly one lap, $200$ m:
\[ t = 200 \div 10 = 20 \text{ (seconds)}. \]
\begin{figC}
  \draw[Slate!60,line width=4pt] (0,0) circle (1.1);
  \fill[Navy] (0,1.1) circle (2.5pt); \node[above=4pt,font=\scriptsize] at (0,1.35) {start};
  \draw[->,tmBlue,thick] (95:1.4) arc (95:265:1.4);
  \draw[->,tmOrange,thick] (85:1.4) arc (85:-85:1.4);
  \node[left,font=\scriptsize,text=tmBlue!80!black] at (-1.5,0) {Mai: $100$ m};
  \node[right,font=\scriptsize,text=tmOrange!80!black] at (1.5,0) {Tú: $100$ m};
  \fill[AnsRed] (0,-1.1) circle (2.5pt); \node[below,font=\scriptsize,text=AnsRed] at (0,-1.5) {they meet after $20$ s};
\end{figC}
After $20$ seconds each has run $5\times 20 = 100$ m, exactly half a lap: they meet at the point opposite the start.
\step{5}{Same direction, Tú faster}
$v_{\text{Tú}} = 21.6\times\dfrac{5}{18} = \dfrac{108}{18} = 6$ (m/s). Every second Tú runs $6 - 5 = 1$ m more than Mai. Tú is
exactly one lap ahead when he has run $200$ m more than Mai:
\[ t = 200 \div 1 = 200 \text{ (seconds)} = 3 \text{ minutes } 20 \text{ seconds}. \]
\verify{After $200$ seconds: Tú has run $6\times 200 = 1200$ m ($6$ laps), Mai $5\times 200 = 1000$ m ($5$ laps). Tú is exactly $1$ lap ahead.}
\begin{parentbox}You can act out parts c) and d) with two coins on a round plate: in opposite directions the ``gap to close''
shrinks by the \emph{sum} of the speeds; in the same direction it shrinks by the \emph{difference} of the speeds.\end{parentbox}
\ans{a) equally fast ($5$ m/s);\quad b) $10$ minutes;\quad c) $20$ seconds;\quad d) $200$ seconds ($3$ minutes $20$ seconds)}
}

% ---------------------------------------------------------------- Problem 38
\begin{probbox}{38}{\stars{1}}
\begin{enumerate}[label=\alph*),leftmargin=1.6em,itemsep=0pt,topsep=0pt]
  \item Find the digits $a$, $b$ such that the number $\nd{4a7b}$ is divisible by $2$, $5$ and $9$.
  \item How many four-digit numbers of the form $\nd{5a3b}$ are divisible by $9$? List them.
\end{enumerate}
\end{probbox}
\probsol{
\textbf{Understanding the problem.} This extends Problem 10: use several divisibility rules at once and count systematically.
\begin{recallbox}Divisible by $2$ and $5$ $\iff$ the last digit is $0$. \ Divisible by $9$ $\iff$ the digit sum is divisible by $9$.\end{recallbox}
\begin{proofbox}Why is ``divisible by both $2$ and $5$'' the same as ``ends in $0$''? Divisibility by $2$ needs an even last digit;
divisibility by $5$ needs a last digit of $0$ or $5$. The only digit that is even and in $\{0;5\}$ is $0$.\end{proofbox}
\step{1}{Part a) -- find $b$}
$\nd{4a7b}$ is divisible by $2$ and $5$, so $b = 0$.
\step{2}{Part a) -- find $a$}
Digit sum: $4 + a + 7 + 0 = 11 + a$. Since $0 \le a \le 9$, we have $11 \le 11 + a \le 20$. The only multiple of $9$ in this
range is $18$. So $11 + a = 18$, $a = 7$. The number is $4770$.
\verify{$4770 = 9\times 530$, and it ends in $0$, so it is divisible by $2$ and $5$.}
\step{3}{Part b) -- the condition}
The digit sum of $\nd{5a3b}$ is $5 + a + 3 + b = 8 + a + b$. Since $0 \le a + b \le 18$, we have $8 \le 8 + a + b \le 26$. The
multiples of $9$ in this range are $9$ and $18$. Hence $a + b = 1$ or $a + b = 10$.
\step{4}{Part b) -- count systematically}
$a + b = 1$: $(a;b) = (0;1)$ or $(1;0)$ \ -- \ $2$ numbers.\\
$a + b = 10$: $a$ runs from $1$ to $9$, $b = 10 - a$ (from $9$ down to $1$) \ -- \ $9$ numbers. ($a = 0$ would give $b = 10$, not a digit.)\\
In total: $2 + 9 = 11$ numbers:
\[ 5031,\ 5130,\ 5139,\ 5238,\ 5337,\ 5436,\ 5535,\ 5634,\ 5733,\ 5832,\ 5931. \]
\begin{warnbox}Do not forget that $a$ may be $0$ (the first digit is $5$), and do not miss the case where the digit sum is $9$
(that is, $a + b = 1$).\end{warnbox}
\verify{Each number above has digit sum $9$ or $18$, e.g.\ $5\,436$: $5+4+3+6 = 18$.}
\ans{a) $a = 7$, $b = 0$ (the number $4770$);\quad b) $11$ numbers}
}

% ---------------------------------------------------------------- Problem 39
\begin{probbox}{39}{\stars{1}}
In the addition below, $M$ and $N$ are two different non-zero digits. Find all pairs $(M;N)$ that work.
\[ \begin{array}{r@{\;}c@{\;}c} & M & N\\ & N & M\\ + & N & N\\ \hline 1 & 8 & 7 \end{array} \]
\end{probbox}
\probsol{
\textbf{Understanding the problem.} This extends Problem 13. This time, instead of trying column by column, we use place value
to get one simple equation, and then find \emph{all} solutions.
\begin{recallbox}$\nd{MN} = 10M + N$, \ $\nd{NM} = 10N + M$, \ $\nd{NN} = 10N + N = 11N$.\end{recallbox}
\step{1}{Write the sum using place value}
\[ \nd{MN} + \nd{NM} + \nd{NN} = (10M + N) + (10N + M) + 11N = 11M + 22N = 11\times(M + 2N). \]
\step{2}{Set up the equation}
$11\times(M + 2N) = 187$, so $M + 2N = 187 \div 11 = 17$.
\step{3}{Find every pair of digits}
$M = 17 - 2N$ must be a digit from $1$ to $9$, so $17 - 2N \le 9$, i.e.\ $N \ge 4$; and $17 - 2N \ge 1$, i.e.\ $N \le 8$.
\begin{center}
\begin{tabular}{|c|c|c|c|c|c|}
\hline
\rowcolor{tmBlue!10} $N$ & $4$ & $5$ & $6$ & $7$ & $8$\\ \hline
$M = 17 - 2N$ & $9$ & $7$ & $5$ & $3$ & $1$\\ \hline
Addition & $94+49+44$ & $75+57+55$ & $56+65+66$ & $37+73+77$ & $18+81+88$\\ \hline
\end{tabular}
\end{center}
All five pairs have $M \ne N$, and every sum equals $187$.
\begin{tipbox}When letters appear ``symmetrically'', as in $\nd{MN}$ and $\nd{NM}$, try place value: there is often a common
factor ($11$ here) that makes the problem much simpler.\end{tipbox}
\begin{warnbox}The problem asks for \emph{all} pairs. Finding one pair (say $M = 9$, $N = 4$) and stopping is not enough.\end{warnbox}
\verify{$94 + 49 + 44 = 187$; \ $75 + 57 + 55 = 187$; \ $56 + 65 + 66 = 187$; \ $37 + 73 + 77 = 187$; \ $18 + 81 + 88 = 187$.}
\ans{$(M;N) \in \{(9;4),\ (7;5),\ (5;6),\ (3;7),\ (1;8)\}$}
}

% ---------------------------------------------------------------- Problem 40
\begin{probbox}{40}{\stars{2}}
A sequence is built by the Fibonacci rule: from the third number on, each number is the sum of the two numbers just before it.
\begin{enumerate}[label=\alph*),leftmargin=1.6em,itemsep=0pt,topsep=2pt]
  \item The first number is $2$ and the fifth number is $19$. Find the second and the eighth numbers of the sequence.
  \item Find the mean of the first five numbers of the sequence in part a).
  \item In the familiar Fibonacci sequence $1, 1, 2, 3, 5, 8, \ldots$, is the $2026$th number even or odd? Explain.
\end{enumerate}
\end{probbox}
\probsol{
\textbf{Understanding the problem.} This extends Problem 15. In part a) the second number is unknown, so we call it $b$ and
express the following numbers in terms of $b$. In part c) we cannot write out $2026$ numbers; we need a pattern of even and odd.
\step{1}{Express the numbers in terms of $b$}
\begin{center}
\begin{tabular}{|c|c|c|c|c|c|}
\hline
\rowcolor{tmBlue!10} Position & 1 & 2 & 3 & 4 & 5\\ \hline
Number & $2$ & $b$ & $2 + b$ & $b + (2 + b) = 2 + 2b$ & $(2 + b) + (2 + 2b) = 4 + 3b$\\ \hline
\end{tabular}
\end{center}
\step{2}{Find $b$}
$4 + 3b = 19 \Rightarrow 3b = 15 \Rightarrow b = 5$. The second number is $5$.
\step{3}{Continue up to the eighth number}
$2,\ 5,\ 7,\ 12,\ 19,\ 31,\ 50,\ 81$ \quad ($12 + 19 = 31$; \ $19 + 31 = 50$; \ $31 + 50 = 81$). The eighth number is $81$.
\step{4}{Mean of the first five numbers}
$(2 + 5 + 7 + 12 + 19) \div 5 = 45 \div 5 = 9$.
\step{5}{The even--odd pattern of the Fibonacci sequence}
\begin{recallbox}odd $+$ odd $=$ even; \ odd $+$ even $=$ odd; \ even $+$ odd $=$ odd.\end{recallbox}
Write O (odd) and E (even) for the numbers $1, 1, 2, 3, 5, 8, 13, 21, 34, \ldots$:
\begin{center}
\begin{tabular}{|c|c|c|c|c|c|c|c|c|c|}
\hline
\rowcolor{tmBlue!10} Position & 1 & 2 & 3 & 4 & 5 & 6 & 7 & 8 & 9\\ \hline
Number & 1 & 1 & 2 & 3 & 5 & 8 & 13 & 21 & 34\\ \hline
Even/odd & O & O & \textbf{E} & O & O & \textbf{E} & O & O & \textbf{E}\\ \hline
\end{tabular}
\end{center}
The pattern ``O, O, E'' repeats forever.
\begin{proofbox}Suppose at some point we have two consecutive numbers O, O. Then the next number is O $+$ O $=$ E; the one after is
O $+$ E $=$ O; the one after that is E $+$ O $=$ O. We again have two consecutive numbers O, O, so the cycle ``O, O, E'' starts
again. The sequence begins with $1, 1$ (O, O), so this pattern repeats all the way: the numbers in positions divisible by $3$
are even, all the others are odd.\end{proofbox}
\step{6}{Look at position $2026$}
$2026 = 3\times 675 + 1$, i.e.\ $2026$ leaves remainder $1$ when divided by $3$. This position corresponds to the first ``O'' in
a cycle. So the $2026$th number is \textbf{odd}.
\begin{warnbox}You cannot conclude ``$2026$ is even, so the $2026$th number is even too''. The parity of the \emph{position} is
different from the parity of the \emph{number} in that position; for example, the $2$nd number is $1$ (odd).\end{warnbox}
\verify{Part a): $2 + 5 = 7$, $5 + 7 = 12$, $7 + 12 = 19$ -- indeed the fifth number is $19$. Part c): test the pattern on small
positions: position $4 = 3\times 1 + 1$ holds $3$ (odd); position $7 = 3\times 2 + 1$ holds $13$ (odd).}
\ans{a) the second number is $5$, the eighth is $81$;\quad b) $9$;\quad c) odd}
}

% ---------------- part d ----------------
% ============================================================
% Topic 4 -- Number Sense -- part d: problems 41--50 (EN)
% Advanced extensions of problems 16--30
% ============================================================

\begin{probbox}{41}{\stars{1}}
\textbf{Sums of consecutive numbers} (extends Problem 16).
\begin{enumerate}[label=\alph*),leftmargin=1.6em,itemsep=1pt]
\item Find \emph{all} the ways of writing $100$ as a sum of at least two consecutive positive integers.
\item Prove that $64$ cannot be written as a sum of at least two consecutive positive integers.
\end{enumerate}
\end{probbox}
\probsol{
\textbf{Understanding the problem.} A sum of consecutive positive integers is determined by two pieces of information: \emph{how many numbers} (call it $k$, with $k\ge 2$) and \emph{the first number} (call it $a$, with $a\ge 1$). We need a quick formula for such a sum, and then we use it to ``narrow down'' the possibilities.

\begin{recallbox}\textbf{Sum of an arithmetic sequence} (Gauss's method): sum $=\dfrac{(\text{first}+\text{last})\times\text{number of terms}}{2}$.\end{recallbox}

\begin{proofbox}Write the sum twice, once forwards and once backwards, and add column by column: for example $S=9+10+\dots+16$ and $S=16+15+\dots+9$. Each column adds up to $9+16=25$ and there are $8$ columns, so $2S=25\cdot 8$, $S=100$. The same argument works for any arithmetic sequence.\end{proofbox}

\step{1}{Set up the key equation}
With $k$ consecutive numbers starting at $a$, the last number is $a+k-1$. Let $m=(\text{first}+\text{last})=2a+k-1$. By Gauss's formula: sum $=\dfrac{m\cdot k}{2}$. If the sum is $100$, then
\[k\cdot m=200.\]

\step{2}{Two important observations about $k$ and $m$}
\begin{itemize}[leftmargin=1.4em,itemsep=1pt]
\item $m-k=2a-1$ is odd, so \textbf{$k$ and $m$ have opposite parity} (one is even, the other odd).
\item $m-k=2a-1\ge 1$, so \textbf{$m>k$}.
\end{itemize}

\step{3}{List the factorisations $200=k\cdot m$ with $2\le k<m$}
Since $k<m$ and $k\cdot m=200$, we get $k\cdot k<200$, so $k\le 14$. The divisors of $200$ from $2$ to $14$ are $2,4,5,8,10$.
\begin{center}
\begin{tabular}{|c|c|c|c|}
\hline
\rowcolor{Tint}$k$ & $m=200\div k$ & Opposite parity? & Result\\
\hline
$2$ & $100$ & no (both even) & rejected\\
$4$ & $50$ & no & rejected\\
$5$ & $40$ & yes & $2a+4=40\Rightarrow a=18$\\
$8$ & $25$ & yes & $2a+7=25\Rightarrow a=9$\\
$10$ & $20$ & no & rejected\\
\hline
\end{tabular}
\end{center}

\step{4}{Write out the sums}
\[100=18+19+20+21+22\qquad\text{and}\qquad 100=9+10+11+12+13+14+15+16.\]
There are exactly \textbf{two} ways.

\begin{tipbox}Remember Problem 16: with $5$ consecutive numbers (an odd number of terms), the middle number equals the mean $=100\div 5=20$, so the five numbers are $18,19,20,21,22$. The ``middle number'' idea gives a very quick check.\end{tipbox}

\step{5}{Part b): why $64$ is impossible}
Suppose $64$ is a sum of $k\ge 2$ consecutive positive integers. Arguing as above: $k\cdot m=2\cdot 64=128$ with $m>k\ge 2$ and $k$, $m$ of opposite parity. So one of $k$, $m$ is odd and divides $128$. But $128=2^7$ has only \emph{one} odd divisor, namely $1$, while both $k$ and $m$ are $\ge 2$. Contradiction! So $64$ cannot be written in this way.

\begin{proofbox}The same argument shows that \emph{every power of $2$} ($2,4,8,16,32,64,\dots$) cannot be written as a sum of at least two consecutive positive integers. Conversely, every number with an odd divisor greater than $1$ can be.\end{proofbox}

\begin{warnbox}Do not forget the word ``positive'': if $0$ and negative numbers were allowed there would be extra ``fake'' ways, e.g.\ $(-17)+(-16)+\dots+21+22=100$ (the numbers from $-17$ to $17$ cancel, leaving $18+19+20+21+22$). Here every number must be $\ge 1$, so $a\ge 1$, which gives $m>k$.\end{warnbox}

\verify{$18+19+20+21+22=100$; $9+10+\dots+16=\dfrac{(9+16)\cdot 8}{2}=100$. Trying $64$ by hand: $2$ consecutive numbers have an odd sum; $3$ consecutive numbers have a sum divisible by $3$; \dots\ no way gives $64$.}
\ans{a) Two ways: $18+19+20+21+22$ and $9+10+\dots+16$. \ b) See the proof.}
}

\begin{probbox}{42}{\stars{1}}
\textbf{Common divisors and counting divisors} (extends Problem 17).
\begin{enumerate}[label=\alph*),leftmargin=1.6em,itemsep=1pt]
\item How many positive integers divide both $720$ and $1080$?
\item Find the smallest positive integer with exactly $12$ divisors.
\end{enumerate}
\end{probbox}
\probsol{
\textbf{Understanding the problem.} Listing the divisors of $720$ and $1080$ would be very long (each has dozens of divisors). We use two powerful tools: (1) the common divisors of two numbers are the divisors of their GCD; (2) the formula counting divisors from the prime factorisation.

\begin{recallbox}\textbf{Counting divisors.} If $n=p^{x}\cdot q^{y}\cdot r^{z}$ ($p,q,r$ different primes), then $n$ has $(x+1)(y+1)(z+1)$ divisors.\end{recallbox}

\begin{proofbox}Each divisor of $n$ has the form $p^{i}\cdot q^{j}\cdot r^{l}$ with $0\le i\le x$, $0\le j\le y$, $0\le l\le z$. The exponent $i$ has $x+1$ choices ($0,1,\dots,x$), $j$ has $y+1$ choices, $l$ has $z+1$ choices. By the multiplication rule there are $(x+1)(y+1)(z+1)$ divisors, and different choices give different divisors (because prime factorisation is unique).\end{proofbox}

\par\smallskip\textbf{Part a).}
\step{1}{Prime factorisation}
$720=2^4\cdot 3^2\cdot 5$ (since $720=16\cdot 45=16\cdot 9\cdot 5$); \quad $1080=2^3\cdot 3^3\cdot 5$ (since $1080=8\cdot 135=8\cdot 27\cdot 5$).

\step{2}{Find the GCD}
Take each common prime with the \emph{smaller} exponent: $\text{GCD}(720,1080)=2^3\cdot 3^2\cdot 5=8\cdot 9\cdot 5=360$.
(Another check: $1080-720=360$ and $720=2\cdot 360$, so the GCD is $360$.)

\step{3}{Count the divisors of $360$}
$360=2^3\cdot 3^2\cdot 5^1$ has $(3+1)(2+1)(1+1)=4\cdot 3\cdot 2=24$ divisors. So $24$ numbers divide both $720$ and $1080$.

\par\smallskip\textbf{Part b).}
\step{4}{The ``shapes'' of a number with exactly $12$ divisors}
The number of divisors is a product of factors (exponent $+1$). We write $12$ as a product of numbers $\ge 2$:
\[12=12,\qquad 12=6\cdot 2,\qquad 12=4\cdot 3,\qquad 12=3\cdot 2\cdot 2.\]
These correspond to four shapes: $p^{11}$, \ $p^{5}q$, \ $p^{3}q^{2}$, \ $p^{2}qr$.

\step{5}{The smallest number of each shape}
To make the number small, use the smallest primes $2,3,5$, and give the \emph{larger} exponent to the \emph{smaller} prime:
\begin{center}
\begin{tabular}{|c|c|c|}
\hline
\rowcolor{Tint}Shape & Smallest number & Value\\
\hline
$p^{11}$ & $2^{11}$ & $2048$\\
$p^{5}q$ & $2^5\cdot 3$ & $96$\\
$p^{3}q^{2}$ & $2^3\cdot 3^2$ & $72$\\
$p^{2}qr$ & $2^2\cdot 3\cdot 5$ & $\mathbf{60}$\\
\hline
\end{tabular}
\end{center}
The smallest is $60$.

\begin{warnbox}Do not give the large exponent to the large prime: for example $2\cdot 3^5=486$ is much bigger than $2^5\cdot 3=96$. And do not forget the shape with \emph{three} primes ($p^2qr$) -- it is the one that gives the answer.\end{warnbox}

\verify{The divisors of $60$: $1,2,3,4,5,6,10,12,15,20,30,60$ -- exactly $12$. A computer check of the numbers from $1$ to $59$: none has $12$ divisors.}
\ans{a) $24$ numbers. \ b) $60$.}
}

\begin{probbox}{43}{\stars{1}}
\textbf{How does a fraction change?} (extends Problems 18 and 19).
\begin{enumerate}[label=\alph*),leftmargin=1.6em,itemsep=1pt]
\item When the same whole number $n$ is added to both the numerator and the denominator of $\dfrac58$, the result is a fraction equal to $\dfrac45$. Find $n$.
\item Let $\dfrac ab$ be a fraction with $0<a<b$. Prove that for every positive integer $n$, $\dfrac{a+n}{b+n}>\dfrac ab$.
\item How many fractions equal to $\dfrac37$ have a two-digit denominator?
\end{enumerate}
\end{probbox}
\probsol{
\textbf{Understanding the problem.} Parts a) and b) are both about ``adding the same number to the numerator and the denominator''. What does \emph{not} change when we do that is the \emph{difference} between the denominator and the numerator -- this is the key.

\par\smallskip\textbf{Part a).}
\step{1}{Find the quantity that stays the same}
For $\dfrac58$, denominator $-$ numerator is $8-5=3$. After adding $n$ to both: $(8+n)-(5+n)=3$, still $3$.

\step{2}{Use a ``parts'' model as in Problem 19}
The new fraction equals $\dfrac45$, so the new numerator is $4$ parts and the new denominator is $5$ parts; the difference is $1$ part. Since the difference is $3$, $1$ part $=3$. New numerator $=4\cdot 3=12$, new denominator $=5\cdot 3=15$.

\step{3}{Find $n$}
$n=12-5=7$ (and also $15-8=7$).

\verify{$\dfrac{5+7}{8+7}=\dfrac{12}{15}=\dfrac{12\div3}{15\div3}=\dfrac45$. Correct.}

\par\smallskip\textbf{Part b).}
\step{4}{Compare ``what is missing to reach $1$''}
Since $a<b$, both fractions are less than $1$. What is missing to make $1$ is
\[1-\frac ab=\frac{b-a}{b},\qquad 1-\frac{a+n}{b+n}=\frac{(b+n)-(a+n)}{b+n}=\frac{b-a}{b+n}.\]
The fractions $\dfrac{b-a}{b}$ and $\dfrac{b-a}{b+n}$ have \emph{the same positive numerator}, and the denominator $b+n>b$. The fraction with the larger denominator is smaller, so $\dfrac{b-a}{b+n}<\dfrac{b-a}{b}$.

\step{5}{Conclusion}
The fraction $\dfrac{a+n}{b+n}$ is ``missing less'' to reach $1$, so it is larger: $\dfrac{a+n}{b+n}>\dfrac ab$.

\begin{recallbox}Two positive fractions with the same numerator: the one with the larger denominator is smaller (the same cake shared among more people gives smaller pieces).\end{recallbox}

\begin{proofbox}Another way (cross-multiplying): with positive denominators, $\dfrac{a+n}{b+n}>\dfrac ab$ exactly when $(a+n)\cdot b>a\cdot(b+n)$, i.e.\ $ab+nb>ab+na$, i.e.\ $nb>na$. This is true because $n>0$ and $b>a$.\end{proofbox}

\begin{warnbox}This is exactly why in Problem 18 we may \emph{not} ``add'' to the numerator and denominator: adding changes the fraction (a fraction less than $1$ gets bigger). Only \emph{multiplying} (or dividing) both by the same non-zero number keeps the fraction the same.\end{warnbox}

\par\smallskip\textbf{Part c).}
\step{6}{General form and condition}
Fractions equal to $\dfrac37$ have the form $\dfrac{3k}{7k}$ ($k$ a positive integer). The denominator $7k$ has two digits: $10\le 7k\le 99$. We have $7\cdot 1=7<10$, $7\cdot 2=14\ge 10$; $7\cdot 14=98\le 99$, $7\cdot 15=105>99$. So $k=2,3,\dots,14$: that is $14-2+1=13$ values.
The fractions are $\dfrac{6}{14},\dfrac{9}{21},\dots,\dfrac{42}{98}$.

\ans{a) $n=7$. \ b) See the proof. \ c) $13$ fractions.}
}

\begin{probbox}{44}{\stars{1}}
\textbf{Poles, a bicycle and a train} (extends Problem 20).
\begin{enumerate}[label=\alph*),leftmargin=1.6em,itemsep=1pt]
\item An cycles steadily at $18$ km/h along a road whose power poles are equally spaced $50$ m apart. How many minutes pass from the moment An passes the first pole to the moment he passes the $25$th pole?
\item A train travelling at constant speed takes $9$ seconds to pass a power pole and $30$ seconds to cross a bridge $350$ m long (from the moment the front of the train reaches the bridge until the end of the train has just left it). Find the length of the train and its speed in km/h.
\end{enumerate}
\end{probbox}
\probsol{
\textbf{Understanding the problem.} Part a) hides the ``fence-post trap'': $25$ poles make how many gaps? Part b) requires understanding how far a train (an object with length) must travel to ``pass'' something.

\begin{recallbox}Distance $=$ speed $\times$ time. Converting $1$ km/h: $1000$ m in $3600$ seconds, so $1$ km/h $=\dfrac{1000}{3600}=\dfrac{5}{18}$ m/s; conversely $1$ m/s $=3.6$ km/h.\end{recallbox}

\par\smallskip\textbf{Part a).}
\step{1}{Count the gaps between poles}
\begin{figC}
\foreach \x in {0,1,2,3}{\draw[Navy,line width=1.5pt] (\x*1.2,0)--(\x*1.2,0.8); \node[below,font=\scriptsize] at (\x*1.2,0) {pole \pgfmathparse{int(\x+1)}\pgfmathresult};}
\foreach \x in {0,1,2}{\draw[<->,tmOrange] (\x*1.2+0.05,0.4)--(\x*1.2+1.15,0.4); \node[above,font=\scriptsize,text=tmOrange!80!black] at (\x*1.2+0.6,0.4) {$50$ m};}
\node at (4.35,0.4) {$\cdots$};
\node[anchor=west,font=\footnotesize,text=Ink] at (5,0.4) {$4$ poles $\to$ $3$ gaps};
\end{figC}
Between the first and the $25$th pole there are $25-1=24$ gaps, with total length $24\cdot 50=1200$ m.

\step{2}{Convert the speed and find the time}
$18$ km/h $=18\cdot\dfrac{5}{18}=5$ m/s. Time: $1200\div 5=240$ seconds $=240\div 60=4$ minutes.

\begin{warnbox}Counting $25$ gaps ($25\cdot 50=1250$ m) gives $250$ seconds -- wrong. The number of gaps is always one less than the number of poles (counting from the first pole to the last).\end{warnbox}

\par\smallskip\textbf{Part b).}
\step{3}{What ``passing a pole'' and ``crossing a bridge'' mean}
\begin{itemize}[leftmargin=1.4em,itemsep=1pt]
\item Passing a pole: from when the front of the train is level with the pole until the end of the train is level with it, the train travels exactly \emph{one train length} $L$.
\item Crossing the bridge: from when the front is at the start of the bridge until the end leaves the far side, the train travels \emph{the bridge length plus the train length}: $350+L$.
\end{itemize}
\begin{figC}
\fill[Slate!25] (2,-0.15) rectangle (7,0); \node[below,font=\scriptsize] at (4.5,-0.15) {bridge $350$ m};
\draw[tmBlue,fill=tmBlue!20,rounded corners=2pt] (0,0.05) rectangle (2,0.45); \node[font=\scriptsize] at (1,0.25) {train (start)};
\draw[tmOrange,fill=tmOrange!20,rounded corners=2pt] (7,0.05) rectangle (9,0.45); \node[font=\scriptsize] at (8,0.25) {train (end)};
\draw[<->,Navy] (2,0.8)--(9,0.8) node[midway,above,font=\scriptsize] {the front travels $350+L$};
\draw[dashed,Slate] (2,0.05)--(2,0.8); \draw[dashed,Slate] (9,0.05)--(9,0.8);
\end{figC}

\step{4}{Compare the two situations}
Crossing the bridge takes $30$ seconds, passing the pole takes $9$ seconds. The difference $30-9=21$ seconds is exactly the time to travel the extra $350$ m (the bridge length). The speed of the train is
\[v=350\div 21=\frac{350}{21}=\frac{50}{3}\ \text{m/s}.\]

\step{5}{Find the train length and convert the speed}
Train length $L=v\cdot 9=\dfrac{50}{3}\cdot 9=150$ m. Speed: $\dfrac{50}{3}\cdot 3.6=\dfrac{50\cdot 3.6}{3}=60$ km/h.

\begin{warnbox}A common mistake: thinking the train travels only $350$ m to cross the bridge (forgetting to add the train length); then the speed would be $350\div 30\approx 11.7$ m/s, which is wrong. Follow \emph{one fixed point} on the train (e.g.\ its front) to know how far the train has travelled.\end{warnbox}

\verify{With $v=\dfrac{50}{3}$ m/s and $L=150$ m: passing the pole takes $150\div\dfrac{50}{3}=9$ seconds; crossing the bridge takes $(350+150)\div\dfrac{50}{3}=500\cdot\dfrac{3}{50}=30$ seconds. Both facts check out.}
\ans{a) $4$ minutes. \ b) The train is $150$ m long; its speed is $60$ km/h.}
}

\begin{probbox}{45}{\stars{2}}
\textbf{Products of digits} (extends Problem 21).
\begin{enumerate}[label=\alph*),leftmargin=1.6em,itemsep=1pt]
\item How many four-digit numbers have digits whose product is $36$?
\item Find the smallest positive integer whose digits have product $360$.
\item Find the largest positive integer with \emph{all different} digits whose product is $360$.
\end{enumerate}
\end{probbox}
\probsol{
\textbf{Understanding the problem.} All three parts start by \emph{factorising into primes}, then ``sharing'' the prime factors among the digits (each digit from $1$ to $9$). Note: no digit can be $0$ (otherwise the product is $0$).

\begin{recallbox}$36=2^2\cdot 3^2$; \ $360=2^3\cdot 3^2\cdot 5$. The digits and their factorisations: $1$; $2$; $3$; $4=2^2$; $5$; $6=2\cdot 3$; $7$; $8=2^3$; $9=3^2$.\end{recallbox}

\par\smallskip\textbf{Part a).}
\step{1}{Find the sets of four digits (ignoring order) with product $36$}
As in Problem 21, list $a\le b\le c\le d$:
\begin{itemize}[leftmargin=1.4em,itemsep=1pt]
\item Two digits equal to $1$: the other two digits have product $36$: $\{1,1,4,9\}$, $\{1,1,6,6\}$.
\item Exactly one digit $1$: the other three digits ($\ge 2$) have product $36$: $\{1,2,2,9\}$, $\{1,2,3,6\}$, $\{1,3,3,4\}$.
\item No digit $1$: four digits $\ge 2$ with product $36$: only $\{2,2,3,3\}$ (since $36=2\cdot2\cdot3\cdot3$ cannot be split further).
\item Three digits equal to $1$: the remaining digit would be $36$ -- impossible.
\end{itemize}

\begin{recallbox}\textbf{Arranging $4$ digits:} all four different: $4\cdot 3\cdot 2\cdot 1=24$ ways. Exactly two equal (the other two different from each other and from them): choose $2$ of the $4$ positions for the equal pair ($6$ ways), then place the other two digits in the $2$ remaining spots ($2$ ways): $12$ ways. Two equal pairs (like $1,1,6,6$): just choose $2$ positions for one pair: $6$ ways.\end{recallbox}

\step{2}{Count}
\begin{center}
\begin{tabular}{|c|c|c|}
\hline
\rowcolor{Tint}Digit set & Type & Arrangements\\
\hline
$\{1,1,4,9\}$ & one equal pair & $12$\\
$\{1,1,6,6\}$ & two equal pairs & $6$\\
$\{1,2,2,9\}$ & one equal pair & $12$\\
$\{1,2,3,6\}$ & all different & $24$\\
$\{1,3,3,4\}$ & one equal pair & $12$\\
$\{2,2,3,3\}$ & two equal pairs & $6$\\
\hline
\multicolumn{2}{|r|}{\textbf{Total}} & $\mathbf{72}$\\
\hline
\end{tabular}
\end{center}

\par\smallskip\textbf{Part b).}
\step{3}{Use as few digits as possible}
A number with fewer digits is smaller. One or two digits give a product of at most $9\cdot 9=81<360$, so we need at least $3$ digits.

\step{4}{Find three digits with product $360$}
$360$ has the prime factor $5$, and the only digit containing the factor $5$ is $5$ itself. So one digit is $5$, and the other two have product $360\div 5=72=8\cdot 9$ (the only way with digits $\le 9$). Digit set: $\{5,8,9\}$.

\step{5}{Arrange to get the smallest number}
Arrange the digits from smallest to largest: $589$.

\par\smallskip\textbf{Part c).}
\step{6}{For the largest number: as many digits as possible}
The digits are different, non-zero, and one of them must be $5$. Six different digits have product at least $1\cdot2\cdot3\cdot4\cdot5\cdot6=720>360$, so there are at most $5$ digits.

\step{7}{Find the sets of $5$ different digits with product $360$}
One digit is $5$; the other four are different, not $5$, with product $72=2^3\cdot 3^2$. Four different digits with a product as small as $72$ must include $1$ (since $2\cdot3\cdot4\cdot6=144>72$); the remaining three different digits have product $72$: $72=2\cdot4\cdot9=3\cdot4\cdot6$ (while $2\cdot 6\cdot 6$ and $3\cdot 3\cdot 8$ repeat a digit, and $1\cdot 8\cdot 9$ repeats the digit $1$). Two sets: $\{1,2,4,5,9\}$ and $\{1,3,4,5,6\}$.

\step{8}{Arrange to get the largest number}
Arrange from largest to smallest: $95421$ and $65431$. The larger is $95421$.

\begin{warnbox}In part c), do not jump to a number with a big first digit such as $985$ (product $360$): a $5$-digit number is always larger than a $3$-digit number. Maximise the \emph{number of digits} first, and only then the leading digit.\end{warnbox}

\verify{$5\cdot 8\cdot 9=360$; $9\cdot5\cdot4\cdot2\cdot1=360$. A computer search over all numbers with up to $6$ digits confirms: there are $72$ four-digit numbers with digit product $36$; the smallest number with digit product $360$ is $589$; the largest with all different digits and digit product $360$ is $95\,421$.}
\ans{a) $72$ numbers. \ b) $589$. \ c) $95\,421$.}
}

\begin{probbox}{46}{\stars{1}}
\textbf{Arithmetic grids} (extends Problem 22). In a square grid, the numbers in every row and in every column form arithmetic sequences.
\begin{enumerate}[label=\alph*),leftmargin=1.6em,itemsep=1pt]
\item A $4\times 4$ grid has these four corner numbers: top left $1$, top right $10$, bottom left $19$, bottom right $37$. Fill in the whole grid.
\item Find the sum of all $16$ numbers in the grid of part a). Compare it with the sum of the four corners: what do you notice? Explain.
\item Prove that in a $3\times 3$ grid of this kind, the number in the centre always equals the mean of the four corner numbers.
\end{enumerate}
\begin{figC}
\foreach \rr in {0,1,2,3}{\foreach \cc in {0,1,2,3}{\draw[Navy,line width=0.9pt,fill=white] (\cc*1.1,-\rr*0.8) rectangle (\cc*1.1+1.1,-\rr*0.8-0.8);}}
\node at (0.55,-0.4) {$1$}; \node at (3.85,-0.4) {$10$};
\node at (0.55,-2.8) {$19$}; \node at (3.85,-2.8) {$37$};
\end{figC}
\end{probbox}
\probsol{
\textbf{Understanding the problem.} Only the four corners are given. The first and last columns (and the first and last rows) each have a known first and last term, so they can be found immediately. Then we use the middle rows or columns.

\begin{recallbox}Arithmetic sequence with $4$ terms: difference $d=(\text{last}-\text{first})\div 3$. Sum of the $4$ terms $=2\times(\text{first}+\text{last})$ (pair term $1$ with term $4$, term $2$ with term $3$: both pairs have the same sum). With $3$ terms: middle term $=$ mean of the two outer terms.\end{recallbox}

\par\smallskip\textbf{Part a).}
\step{1}{First and last rows}
Row $1$: from $1$ to $10$, $d=(10-1)\div 3=3$: \ $1,4,7,10$.\\
Row $4$: from $19$ to $37$, $d=(37-19)\div 3=6$: \ $19,25,31,37$.

\step{2}{The four columns (first and last terms of each are known)}
Column $1$: $1\to 19$, $d=6$: $1,7,13,19$. \quad Column $2$: $4\to 25$, $d=7$: $4,11,18,25$.\\
Column $3$: $7\to 31$, $d=8$: $7,15,23,31$. \quad Column $4$: $10\to 37$, $d=9$: $10,19,28,37$.

\step{3}{Check the middle rows}
Row $2$: $7,11,15,19$ (difference $4$); row $3$: $13,18,23,28$ (difference $5$). The grid satisfies every condition:
\begin{figC}
\foreach \rr in {0,1,2,3}{\foreach \cc in {0,1,2,3}{\draw[Navy,line width=0.9pt,fill=white] (\cc*1.1,-\rr*0.8) rectangle (\cc*1.1+1.1,-\rr*0.8-0.8);}}
\foreach \cc/\val in {0/1,1/4,2/7,3/10}{\node at (\cc*1.1+0.55,-0.4) {$\val$};}
\foreach \cc/\val in {0/7,1/11,2/15,3/19}{\node at (\cc*1.1+0.55,-1.2) {$\val$};}
\foreach \cc/\val in {0/13,1/18,2/23,3/28}{\node at (\cc*1.1+0.55,-2.0) {$\val$};}
\foreach \cc/\val in {0/19,1/25,2/31,3/37}{\node at (\cc*1.1+0.55,-2.8) {$\val$};}
\foreach \rr/\val in {0/22,1/52,2/82,3/112}{\node[anchor=west,font=\footnotesize,text=Slate] at (4.55,-\rr*0.8-0.4) {row sum: $\val$};}
\end{figC}

\par\smallskip\textbf{Part b).}
\step{4}{Compute the sum}
Sum of the rows: $22+52+82+112=268$. Sum of the corners: $1+10+19+37=67$. We see that $268=4\cdot 67$: \textbf{the sum of the whole grid is $4$ times the sum of the corners.}

\step{5}{Explanation (true for every such $4\times4$ grid)}
Each row is an arithmetic sequence of $4$ terms, so each row sum $=2\times$(first term $+$ last term of the row). Adding the $4$ rows:
\[\text{Grid sum}=2\times(\text{sum of column }1+\text{sum of column }4).\]
Columns $1$ and $4$ are also arithmetic sequences of $4$ terms: sum of column $1=2\times(1+19)$, sum of column $4=2\times(10+37)$. So
\[\text{Grid sum}=2\times\big(2\cdot(1+19)+2\cdot(10+37)\big)=4\times(1+10+19+37)=4\times 67=268.\]

\par\smallskip\textbf{Part c).}
\step{6}{Name the cells of the $3\times 3$ grid}
\begin{figC}
\foreach \rr in {0,1,2}{\foreach \cc in {0,1,2}{\draw[Navy,line width=0.9pt,fill=white] (\cc*1.1,-\rr*0.8) rectangle (\cc*1.1+1.1,-\rr*0.8-0.8);}}
\node at (0.55,-0.4) {$a$}; \node at (1.65,-0.4) {$\cdot$}; \node at (2.75,-0.4) {$b$};
\node at (0.55,-1.2) {$u$}; \node[text=AnsRed] at (1.65,-1.2) {$x$}; \node at (2.75,-1.2) {$v$};
\node at (0.55,-2.0) {$c$}; \node at (1.65,-2.0) {$\cdot$}; \node at (2.75,-2.0) {$d$};
\end{figC}
The corners are $a,b,c,d$; the centre is $x$; $u$ and $v$ are the middle cells of the left and right columns.

\step{7}{Use ``middle term $=$ mean of the outer terms''}
The left column $a,u,c$ is arithmetic, so $u=\dfrac{a+c}{2}$. Right column: $v=\dfrac{b+d}{2}$. The middle row $u,x,v$ is arithmetic, so
\[x=\frac{u+v}{2}=\frac{\frac{a+c}{2}+\frac{b+d}{2}}{2}=\frac{a+b+c+d}{4}.\]
This is exactly the mean of the four corner numbers.

\begin{proofbox}Why does an arithmetic sequence $p,q,r$ satisfy $q=\dfrac{p+r}{2}$? Because $q-p=r-q$ (both equal the difference $d$), so $2q=p+r$.\end{proofbox}

\verify{Apply part c) to the top-left $3\times 3$ block of the grid in part a): corners $1,7,13,23$, mean $\dfrac{1+7+13+23}{4}=11$, exactly the centre $11$.}
\ans{a) See the grid. \ b) Sum $268=4\times 67$ ($4$ times the sum of the corners). \ c) See the proof.}
}

\begin{probbox}{47}{\stars{2}}
\textbf{``Rotating'' numbers} (extends Problem 23). In each part, every letter represents a digit, and $A\neq 0$.
\begin{enumerate}[label=\alph*),leftmargin=1.6em,itemsep=1pt]
\item Find the six-digit number $\nd{2ABCDE}$ such that $\nd{2ABCDE}\times 3=\nd{ABCDE2}$.
\item Find all digits $d$ (from $1$ to $9$) for which there is a six-digit number $\nd{dABCDE}$ with $\nd{dABCDE}\times 3=\nd{ABCDEd}$.
\item Prove that there is no six-digit number $\nd{1ABCDE}$ with $\nd{1ABCDE}\times 2=\nd{ABCDE1}$.
\end{enumerate}
\end{probbox}
\probsol{
\textbf{Understanding the problem.} All three parts are about ``moving the first digit to the end''. We reuse Method 1 of Problem 23: treat the block $\nd{ABCDE}$ as one number $N$ and set up an equation. This time we do it once for every leading digit $d$.

\begin{recallbox}With $N=\nd{ABCDE}$: \ $\nd{dABCDE}=d\cdot 100\,000+N$ \ and \ $\nd{ABCDEd}=10\cdot N+d$. Since $A\neq 0$, $N$ has exactly five digits: $10\,000\le N\le 99\,999$.\end{recallbox}

\step{1}{The general equation for parts a) and b)}
\[3\cdot(d\cdot 100\,000+N)=10N+d\ \Longrightarrow\ 300\,000\,d+3N=10N+d\ \Longrightarrow\ 299\,999\,d=7N.\]
Since $299\,999=7\cdot 42\,857$, dividing both sides by $7$ gives
\[N=42\,857\cdot d.\]

\step{2}{Part a): $d=2$}
$N=42\,857\cdot 2=85\,714$. The number is $285\,714$.

\step{3}{Part b): when does $N$ have five digits?}
$N=42\,857\,d\le 99\,999$. For $d=1$: $N=42\,857$ (five digits, fine). For $d=2$: $N=85\,714$ (fine). For $d=3$: $N=128\,571$ has six digits -- too big to be the block $\nd{ABCDE}$. For $d\ge 3$, $N$ is even bigger. So only $d=1$ and $d=2$ work, giving the two numbers $142\,857$ and $285\,714$.

\begin{tipbox}These two numbers use the same digits $1,4,2,8,5,7$ arranged in a circle. That is no accident: $\dfrac17=0.\overline{142857}$ and $\dfrac27=0.\overline{285714}$. Multiplying by $3$ corresponds to ``rotating'' the circle of digits.\end{tipbox}

\step{4}{Part c): set up the equation}
With $N=\nd{ABCDE}$:
\[2\cdot(100\,000+N)=10N+1\ \Longrightarrow\ 200\,000+2N=10N+1\ \Longrightarrow\ 199\,999=8N.\]

\step{5}{Part c): an odd/even contradiction}
The right-hand side $8N$ is even (since $8$ is even). The left-hand side $199\,999$ is odd. An odd number cannot equal an even number. So no such number exists.

\begin{proofbox}You can see it immediately without an equation: $\nd{1ABCDE}\times 2$ is even, so its last digit must be even; but $\nd{ABCDE1}$ ends in $1$ (odd). Contradiction.\end{proofbox}

\begin{warnbox}In part b), many students stop after finding $d=1$ (the answer of Problem 23) and forget $d=2$; or they accept $d=3,4,\dots$ without checking whether $N$ has more than five digits. Always check the condition ``$N$ has exactly five digits''.\end{warnbox}

\verify{$142\,857\times 3=428\,571$; \ $285\,714\times 3=857\,142$. In both cases the first digit has indeed moved to the end.}
\ans{a) $285\,714$. \ b) $d=1$ or $d=2$ (the numbers $142\,857$ and $285\,714$). \ c) See the proof.}
}

\begin{probbox}{48}{\stars{2}}
\textbf{A frog on the number line} (extends Problem 24). A frog sits at the point $0$ on the number line. On its $k$-th jump ($k=1,2,3,\dots$), the frog jumps exactly $k$ units, to the left or to the right as it likes.
\begin{enumerate}[label=\alph*),leftmargin=1.6em,itemsep=1pt]
\item After exactly $10$ jumps, can the frog be back at $0$? Why?
\item After exactly $10$ jumps, which points can the frog be at? How many such points are there?
\item What is the smallest number of jumps the frog needs to reach the point $20$?
\end{enumerate}
\end{probbox}
\probsol{
\textbf{Understanding the problem.} The frog's position after $n$ jumps is $\pm1\pm2\pm3\cdots\pm n$ (a $+$ sign for a jump to the right, a $-$ sign for a jump to the left). The idea is the same as in Problem 24: start from ``all plus'' and see how switching signs changes the result.

\begin{recallbox}$1+2+\dots+n=\dfrac{n(n+1)}{2}$ (Gauss's formula). For example $1+2+\dots+10=\dfrac{10\cdot 11}{2}=55$.\end{recallbox}

\step{1}{The ``switching'' rule}
If the frog always jumps right, after $10$ jumps it is at $55$. Switching one jump $k$ from right to left lowers the position by $2k$ (it loses $k$ and goes back another $k$). So if the jumps to the left add up to $s$, the final position is
\[55-2s.\]

\par\smallskip\textbf{Part a).}
\step{2}{Odd or even}
$55-2s$ is always \emph{odd} (odd minus even). The point $0$ is even, so the frog \textbf{cannot} be back at $0$ after $10$ jumps.

\par\smallskip\textbf{Part b).}
\step{3}{Which values can $s$ take?}
$s$ is the sum of a group (possibly empty) of different numbers from $\{1,2,\dots,10\}$, so $0\le s\le 55$. We show that \emph{every} integer $s$ from $0$ to $55$ can be made.

\begin{proofbox}Take the largest numbers one by one: $10, 9, 8, \dots$ The totals $0$, $10$, $10+9=19$, $19+8=27$, $27+7=34$, $34+6=40$, $40+5=45$, $45+4=49$, $49+3=52$, $52+2=54$, $54+1=55$ cut the interval $[0,55]$ into pieces. A number $s$ between two consecutive marks, e.g.\ $27<s\le 34$, is written as $s=10+9+8+(s-27)$ with $s-27$ from $1$ to $7$ -- a number not yet used. The same works for every other piece (for $0<s\le 10$ just take the number $s$). So every $s$ from $0$ to $55$ can be made.\end{proofbox}

\step{4}{The final positions}
For $s=0,1,2,\dots,55$, the position $55-2s$ takes the values $55,53,51,\dots,-53,-55$: that is, \textbf{every odd number from $-55$ to $55$}. Counting: from $-55$ to $55$ there are $28$ positive odd numbers ($1,3,\dots,55$) and $28$ negative ones, $56$ points altogether (the same as the number of values of $s$: $56$).

\begin{figC}
\draw[->,thick,Navy] (-5.9,0)--(5.9,0);
\foreach \x in {-5.5,-4.5,-3.5,-2.5,-1.5,-0.5,0.5,1.5,2.5,3.5,4.5,5.5}{\fill[KidOrange] (\x,0) circle (2.3pt);}
\draw[Navy] (0,0.1)--(0,-0.1); \node[below,font=\scriptsize] at (0,-0.1) {$0$};
\node[below,font=\scriptsize] at (-5.5,-0.1) {$-55$}; \node[below,font=\scriptsize] at (5.5,-0.1) {$55$};
\node[below,font=\scriptsize] at (0.5,-0.1) {$1$}; \node[below,font=\scriptsize] at (-0.5,-0.1) {$-1$};
\node[above,font=\scriptsize,text=Slate] at (0,0.15) {odd numbers only (sketch, not to scale)};
\end{figC}

\par\smallskip\textbf{Part c).}
\step{5}{Two necessary conditions}
After $n$ jumps, let $T=1+2+\dots+n$. The position is $T-2s$, so: (1) $T\ge 20$ (the frog cannot get farther than $T$); (2) $T-20$ must be even, i.e.\ $T$ must be even like $20$.

\step{6}{Try $n$ one at a time}
\begin{center}
\begin{tabular}{|c|c|c|c|c|c|c|c|}
\hline
\rowcolor{Tint}$n$ & $1$ & $2$ & $3$ & $4$ & $5$ & $6$ & $7$\\
\hline
$T$ & $1$ & $3$ & $6$ & $10$ & $15$ & $21$ & $28$\\
\hline
Reach $20$? & \multicolumn{5}{c|}{no ($T<20$)} & no ($T$ odd) & \textbf{possible}\\
\hline
\end{tabular}
\end{center}

\step{7}{Show a way with $n=7$}
We need $28-2s=20$, i.e.\ $s=4$: just jump left on the $4$th jump:
\[1+2+3-4+5+6+7=20.\]
So the smallest number of jumps is $7$.

\begin{warnbox}Checking only ``$T\ge 20$'' is not enough: for $n=6$, $T=21\ge 20$, but the position is always odd and can never be $20$. This is exactly the odd/even idea of Problem 24.\end{warnbox}

\verify{$1+2+3-4+5+6+7=20$. For $n=6$ the position is always $21-2s$, an odd number, never $20$. A computer search of all jump patterns also confirms: after $10$ jumps there are exactly $56$ positions, not including $0$; the smallest number of jumps to reach $20$ is $7$.}
\ans{a) No (the position is always odd). \ b) Every odd number from $-55$ to $55$: $56$ points. \ c) $7$ jumps.}
}

\begin{probbox}{49}{\stars{1}}
\textbf{Numbers containing the block $2025$} (extends Problem 28). Consider the six-digit numbers that contain the four digits $2,0,2,5$ \emph{together and in this order} (for example: $720\,251$, $202\,513$, $132\,025$). How many of these numbers are divisible by $45$?
\end{probbox}
\probsol{
\textbf{Understanding the problem.} The number consists of the block $2025$ plus two more digits $a,b$ in different positions. ``Divisible by $45$'' should be split into two conditions that are easier to check.

\begin{recallbox}$45=5\cdot 9$, and $5$ and $9$ have no common divisor other than $1$. So a number is divisible by $45$ exactly when it is divisible by both $5$ and $9$.
\textbf{Divisible by $5$:} the last digit is $0$ or $5$.
\textbf{Divisible by $9$:} the digit sum is divisible by $9$.\end{recallbox}

\step{1}{The condition for divisibility by $9$}
The digits of the block add up to $2+0+2+5=9$. The whole number has digit sum $9+a+b$, which is divisible by $9$ when $a+b$ is divisible by $9$, i.e.\ $a+b\in\{0,\ 9,\ 18\}$.

\step{2}{The three positions of the block, case by case}
\begin{itemize}[leftmargin=1.4em,itemsep=2pt]
\item \textbf{Form $\nd{2025ab}$} (the last digit is $b$, so $b\in\{0,5\}$; $a$ is any digit):
  $b=0$: $a\in\{0,9\}$ $\to$ $202\,500$, $202\,590$. \ $b=5$: $a+5\in\{9,18\}$ $\to$ $a=4$ $\to$ $202\,545$. \ \textbf{$3$ numbers.}
\item \textbf{Form $\nd{a2025b}$} ($a\neq 0$, $b\in\{0,5\}$):
  $b=0$: $a=9$ $\to$ $920\,250$. \ $b=5$: $a=4$ $\to$ $420\,255$. \ \textbf{$2$ numbers.}
\item \textbf{Form $\nd{ab2025}$} (ends in $5$ -- always divisible by $5$; $a\neq 0$):
  $a+b=9$: $a=1,2,\dots,9$ (with $b=9-a$) $\to$ $9$ numbers: $182\,025$, $272\,025$, \dots, $902\,025$. \ $a+b=18$: $a=b=9$ $\to$ $992\,025$. \ ($a+b=0$ is impossible because $a\neq0$.) \ \textbf{$10$ numbers.}
\end{itemize}

\step{3}{Check for double counting}
A number would be counted twice if the block $2025$ appeared twice among its six digits, i.e.\ two overlapping blocks. Shifting the block by $1$ place: ``$025$'' would have to match ``$202$'' -- false; shifting by $2$ places: ``$25$'' would have to match ``$20$'' -- false. So no number is counted twice.

\step{4}{Add up}
$3+2+10=15$ numbers.

\begin{warnbox}(1) Do not forget the case $a+b=0$ in the form $\nd{2025ab}$ (the number $202\,500$). (2) In the forms $\nd{a2025b}$ and $\nd{ab2025}$, the first digit $a$ cannot be $0$. (3) Check \emph{both} conditions, divisibility by $5$ and by $9$, not just one.\end{warnbox}

\verify{Example $992\,025$: it ends in $5$, and its digit sum $9+9+2+0+2+5=27$ is divisible by $9$, so it is divisible by $45$ ($992\,025=45\cdot 22\,045$). A computer search over all six-digit numbers also gives exactly $15$.}
\ans{$15$ numbers.}
}

\begin{probbox}{50}{\stars{2}}
\textbf{Mean and median} (extends Problem 30). A list of five positive integers (repeats allowed) has mean $12$ and median $10$.
\begin{enumerate}[label=\alph*),leftmargin=1.6em,itemsep=1pt]
\item What is the largest possible value of the largest number in the list?
\item What is the smallest possible value of the largest number in the list?
\item Answer both questions again if the five numbers must all be different.
\end{enumerate}
\end{probbox}
\probsol{
\textbf{Understanding the problem.} Arrange the five numbers in increasing order: $a\le b\le c\le d\le e$. The median is the middle number, so $c=10$. The mean is $12$, so the sum is $5\cdot 12=60$. We want to make $e$ as large or as small as possible. This is an ``optimisation'' problem: we must both \emph{prove} that a limit cannot be passed and \emph{give an example} that reaches exactly that limit.

\begin{recallbox}Mean $=$ sum $\div$ how many numbers, so sum $=12\cdot 5=60$. The median of five numbers (in increasing order) is the third one.\end{recallbox}

\step{1}{The key equation}
\[a+b+10+d+e=60\quad\Longrightarrow\quad a+b+d+e=50,\qquad a\le b\le 10\le d\le e.\]
\begin{barfig}
\foreach \x/\lab/\col in {0/a/tmBlue,1.3/b/tmBlue,2.6/10/tmGreen,3.9/d/tmOrange,5.2/e/tmOrange}{\blk{\x}{\x+1.2}{0}{0.6}{\col}{$\lab$}}
\node[anchor=west,font=\footnotesize] at (6.6,0) {sum $=60$, \ median $=10$};
\end{barfig}

\par\smallskip\textbf{Part a): make $e$ as large as possible.}
\step{2}{Upper limit}
For $e$ to be large, the other numbers must be as small as possible: $a\ge 1$, $b\ge 1$, $d\ge 10$. Hence $e=50-a-b-d\le 50-1-1-10=38$.
\step{3}{An example that reaches it}
The list $1,1,10,10,38$: sum $60$, median $10$. So the largest possible $e$ is $38$.

\par\smallskip\textbf{Part b): make $e$ as small as possible.}
\step{4}{Lower limit}
For $e$ to be small, $a$ and $b$ must be as large as possible, but $a\le b\le 10$, so $a+b\le 20$. Hence $d+e=50-(a+b)\ge 30$. Since $d\le e$, $2e\ge d+e\ge 30$, so $e\ge 15$.
\step{5}{An example that reaches it}
The list $10,10,10,15,15$: sum $60$, median $10$. So the smallest possible $e$ is $15$.

\par\smallskip\textbf{Part c): five different numbers} ($a<b<10<d<e$).
\step{6}{Largest $e$}
As small as possible: $a\ge 1$, $b\ge 2$, $d\ge 11$. So $e\le 50-1-2-11=36$. Example: $1,2,10,11,36$ (sum $60$). So the largest $e$ is $36$.
\step{7}{Smallest $e$}
As large as possible: $b\le 9$, $a\le 8$, so $a+b\le 17$ and $d+e\ge 50-17=33$. Since $d<e$, $d\le e-1$, so $2e-1\ge d+e\ge 33$, i.e.\ $e\ge 17$. Example: $8,9,10,16,17$ (sum $60$). So the smallest $e$ is $17$.

\begin{warnbox}A ``largest/smallest'' problem needs \emph{two} parts: (1) a proof that the limit cannot be passed, and (2) a concrete example reaching exactly that limit. An example alone may not be the best; a limit alone may not be reachable.\end{warnbox}

\begin{parentbox}Let your child share $60$ sweets among $5$ ``boxes'' kept in order, with the middle box always holding $10$. Ask: to make the last box as full as possible, should the other boxes hold few or many? And to make the last box as small as possible? This game helps children ``feel'' optimisation reasoning.\end{parentbox}

\verify{$1+1+10+10+38=60$; $10+10+10+15+15=60$; $1+2+10+11+36=60$; $8+9+10+16+17=60$; all have median $10$. A computer search over all lists also confirms the values $38$, $15$, $36$, $17$.}
\ans{a) $38$. \ b) $15$. \ c) Largest $36$, smallest $17$.}
}

\end{document}
